% !TeX root = ../Book5.tex
% 纲要标题：# 第四编　高级表示论与代数对称
% 纲要位置：工作记录/计划纲要.md，第 4358 行。

\BookPart{高级表示论与代数对称}{b5:part:04}
\BookPartGuide
前面建立了有限群表示的分解方法、ADE 箭图的不可分解分类与复半单 Lie 代数的有限维最高权分类。本编进一步讨论无限维最高权模组成的范畴、群的坐标环上的张量结构，以及不变多项式与商的构造。范畴 \(\mathcal O\) 中的中心特征标与 BGG 消解用来研究最高权模；坐标环的余乘法记录群的乘法，表示对应的余模则说明群作用怎样在张量积上组合；不变环的生成元与关系具体描述多项式函数怎样形成商。

\ProseParagraphBreak
\chapref{b5:ch:20}有两组问题。研究无限维最高权模时，可读第20.1、20.2及20.4节：范畴 \(\mathcal O\) 的有限性条件、中心作用和 BGG 方法依次服务于这一问题。为完成\secref{b5:sec:1204}留下的 Levi 分解与共轭问题，可以直接读第20.3节所证明的 Whitehead 引理\footnote{亨利·怀特海德（John Henry Constantine Whitehead，1904-1960），英国数学家，是同伦论的奠基者之一；此处的怀特海德不是其叔父、数学家兼哲学家阿尔弗雷德·诺思·怀特海德。}及可解根归纳，不必先学范畴 \(\mathcal O\) 和 Harish--Chandra 同构\footnote{哈里什-钱德拉（Harish-Chandra，1923-1983），印度裔美国数学家，在半单 Lie 群的表示论与调和分析方面作出重要贡献。}。第20.4节的一般 BGG 证明概要使用相对 Chevalley--Eilenberg 复形\footnote{舍瓦莱（Claude Chevalley，1909-1984）是法国数学家，研究类域论和代数群；艾伦伯格（Samuel Eilenberg，1913-1998）是波兰裔美国数学家，创立范畴论并研究代数拓扑。二人共同建立了 Lie 代数上同调的复形。}，追踪该概要时仍需补读第20.3节。第三编的根系及特征标公式只用此前已证的结构与表示论结果，可以独立阅读。

\ProseParagraphBreak
\chapref{b5:ch:21}从余乘法、对极和张量表示开始，再把代数群写成坐标环。前半只需张量积与代数同态，后半使用仿射几何的基本语言。群代数与函数代数、模与余模、群点与群概形分别比较；古典群的型与对合例子最后说明，基域和中心商仍会影响表示及点群。

\ProseParagraphBreak
\chapref{b5:ch:22}的前四节可以在学完第一编并具备多项式基础后阅读。先证明不变环有限生成，再用次数界和 Molien 公式\footnote{莫利恩（Theodor Georg Andreas Molien，1861-1941），波罗的海德裔俄国数学家，研究结合代数及有限群多项式不变量。}寻找生成元，最后求出生成元之间的全部关系。经典群不变量与仿射商随后解释这些计算的几何意义。几何不变量理论部分证明线性约化群的有限生成、幂提升与射影商的仿射图构造，并给出 Haboush 定理\footnote{哈布什（William Joseph Haboush，1942-），美国数学家，证明约化代数群的几何约化性。}与 Nagata 定理\footnote{永田雅宜（假名：ながた まさよし，罗马音：Masayoshi Nagata，1927-2008），日本数学家，研究可换代数和代数几何，给出 Hilbert 第十四问题的反例。}的具体概要及低维算例。

\ProseParagraphBreak
附录 A 的组合与反射群工具承接第6、7章；附录 B 的量子群、簇变异和范畴化主要通过小模型说明方法。附录 C 的 Lie 群、算术与几何表示论另需相应的拓扑、分析或几何知识。附录 D.2.1 节使用附录 E 的 Picard--Vessiot 术语\footnote{皮卡（Charles Émile Picard，1856-1941）与维西奥（Ernest Vessiot，1865-1952）都是法国数学家，发展了线性微分方程的 Galois 理论；皮卡还研究复分析，维西奥还研究微分方程对称性。}与例子；附录 D.1 节的微分算子模计算则可独立于附录 E 进行。这些专题仍须逐项核对其底域和适用条件。

% !TeX root = ../../Book5.tex
% 纲要标题：## 第20章*　范畴 O 与 Lie 代数上同调
% 纲要位置：工作记录/计划纲要.md，第 4360 行。

\OptionalChapter{范畴 O 与 Lie 代数上同调}{b5:ch:20}
\BookChapterNumber{b5:ch:20}

\begin{chapterabstract}
范畴 \(\mathcal O\) 把无限维最高权模及其扩张纳入一个具有有限性条件的交换范畴。中心作用给出无穷小特征标，秩一计算已经显示它按移位 Weyl 轨道分组。Lie 代数上同调则记录截面与同态的误差：本章证明 Whitehead 两个引理，并完成一般 Levi 分解的存在性与共轭性。最后介绍 BGG 消解，明确区分 Verma 消解和投射消解，以秩一正合列、对偶扩张及微分算子实例说明这些区别。
\end{chapterabstract}

\begin{learninggoals}
\begin{itemize}
\item 检查范畴 \(\mathcal O\) 的三个定义条件，证明权空间有限性及核、余核的封闭性。
\item 从 PBW 投影计算中心作用，完整求出 \(\mathfrak{sl}_2\) 包络代数的中心。
\item 使用 Chevalley--Eilenberg 微分，解释一、二阶上同调并证明 Whitehead 引理与一般 Levi--Malcev 定理。
\item 核验秩一 BGG 正合列，区分中心特征标、简单模参数以及不同范畴中的扩张。
\end{itemize}
\end{learninggoals}

第20.1、20.2及20.4节研究无限维最高权模：先确定范畴 \(\mathcal O\) 中的对象，再研究中心特征标和 BGG 消解。这部分需要 PBW、Verma 模与最高权分类；一般 Harish--Chandra 同构和一般 BGG 消解只给证明概要，秩一中心与秩一消解在本卷证明。本章并未建立一般范畴 \(\mathcal O\) 的有限长度、投射覆盖及标准滤过理论，因此这里的讨论不构成其完整的最高权范畴理论。

\ProseParagraphBreak
第20.3节以 Lie 代数上同调处理\secref{b5:sec:1204}的 Levi 分解问题，所需的是有限维 Lie 代数、模、短正合列以及第三编已经证明的完全可约性与 Casimir 结果；不以范畴 \(\mathcal O\) 或一般中心同构为前提。交换范畴与上同调长正合列的基础见第四卷，本节也写出所用复形及低阶公式。Whitehead 引理与 Levi 定理均在本卷证明。第20.4节的一般 BGG 证明概要还要在第20.3节的 Chevalley--Eilenberg 复形基础上使用相对版本，故两组问题在这里有一处实际联系。

\ProseParagraphBreak
参考\cite[Sections 8.4 and 8.8]{Kirillov2008LieGroups}可对照 BGG 消解与中心的两种描述。Lie 代数积分和旗簇微分算子的进一步例子放在\appref{b5:app:C}，本章的代数证明不依赖这些分析或几何补充。

% !TeX root = ../../Book5.tex
% 纲要标题：### 20.1 范畴 O
% 纲要位置：工作记录/计划纲要.md，第 4362 行。

\section{范畴 O}
\label{b5:sec:2001}

% 纲要范围：
%
% * 有限生成性、权分解及局部有限性
% * Verma 模和简单模
% * 本节需要第四卷的交换范畴语言
%

复半单 Lie 代数的有限维表示完全可约，研究 Verma 模中的非分裂扩张却必须允许无限维对象。全部包络代数模过于宽泛：Cartan 作用可能没有权分解，权空间也可能无限维。范畴 O 选取有限生成的权模，并限制正根方向的作用，使最高权方法仍能控制权集和各权空间。

\subsection{有限生成性、权分解与局部有限性}
\label{b5:subsec:2001:01}

\begin{definition}\label{b5:def:category-o}
设 \(\mathfrak g\) 为有限维复半单 Lie 代数，固定三角分解
\(\mathfrak g=\mathfrak n^-\oplus\mathfrak h\oplus\mathfrak n^+\)，其中 \(\mathfrak h\) 为 Cartan 子代数，\(\mathfrak n^\pm\) 分别为所选正、负根空间之和。在幺性左 \(U(\mathfrak g)\) 模范畴中，选取满足下列条件的模 \(M\)：
\begin{enumerate}
\item 模 \(M\) 作为 \(U(\mathfrak g)\) 模有限生成；
\item \(M=\bigoplus_{\mu\in\mathfrak h^*}M_\mu\)，其中 \(M_\mu\) 是共同特征值为 \(\mu\) 的 Cartan 权空间；
\item 对每个 \(v\in M\)，向量空间 \(U(\mathfrak n^+)v\) 有限维。
\end{enumerate}
这些模及它们之间的全部模同态组成的全子范畴，称为\term[fan-chou-O]{范畴 \(\mathcal O\)}[category O]。第三条称 \(\mathfrak n^+\) 的作用\term[jubu-youxian-zuoyong]{局部有限}[locally finite]。
\end{definition}
\symidx{\(\mathcal O\)：半单 Lie 代数的范畴 O}

正根方向的局部有限性有独立的作用。例如取 \(\mathfrak g=\mathfrak{sl}_2(\mathbb C)\)，标准基为 \(e,f,h\)，并固定 \(z\in\mathbb C\setminus\mathbb Z\)。在
\(W=\bigoplus_{j\in\mathbb Z}\mathbb Cv_j\) 上规定
\[
hv_j=(z-2j)v_j,\qquad fv_j=v_{j+1},\qquad
ev_j=j(z-j+1)v_{j-1}.
\]
前两条括号关系由权的升降直接得到，而
\[
(ef-fe)v_j=\bigl((j+1)(z-j)-j(z-j+1)\bigr)v_j
=(z-2j)v_j=hv_j,
\]
所以这确为一个模。它由 \(v_{-1}\) 生成，每个权空间一维，却有
\[
e^rv_{-1}=(-1)^r r!\prod_{s=1}^r(z+s+1)v_{-r-1}\ne0
\qquad(r\ge1).
\]
这些向量的权互不相同，故 \(U(\mathfrak n^+)v_{-1}\) 无限维。这个模满足前两条定义条件，但不属于 \(\mathcal O\)；其权向上没有界。

\begin{proposition}\label{b5:prop:o-weight-finiteness}
设 \(\mathfrak g\) 为有限维复半单 Lie 代数，固定 Cartan 子代数 \(\mathfrak h\) 及简单根 \(\alpha_1,\ldots,\alpha_r\)，令 \(\mathcal O\) 为相应的范畴 O，\(Q_+=\sum_i\mathbb Z_{\ge0}\alpha_i\)。设 \(M\in\mathcal O\)，记 \(M_\mu\) 为权 \(\mu\in\mathfrak h^*\) 的空间，\(\operatorname{wt}M\) 为全部非零权空间的权集合。则每个 \(M_\mu\) 有限维，且存在有限集 \(F\subset\mathfrak h^*\)，使
\[
\operatorname{wt}M\subseteq\bigcup_{\nu\in F}(\nu-Q_+).
\]
若 \(M\ne0\)，它含有一个非零最高权向量。
\end{proposition}
\begin{proof}
把一组有限生成元分别拆成有限个权分量，得到有限个权生成元 \(v_j\)。每个 \(U(\mathfrak n^+)v_j\) 有限维，而且由权向量生成，故被 \(\mathfrak h\) 保持。因此
\(E=\sum_jU(\mathfrak b)v_j=\sum_jU(\mathfrak n^+)v_j\)
是有限维 \(\mathfrak b\) 子模。PBW 给出 \(M=U(\mathfrak n^-)E\)。取 \(E\) 的权集为 \(F\)，即得权集包含式。

对固定 \(\mu\)，每个 \(\nu\in F\) 所贡献的负根单项式满足
\(\sum_{\alpha>0}a_\alpha\alpha=\nu-\mu\)。右端若不在 \(Q_+\)，没有贡献；若在 \(Q_+\)，所有 \(a_\alpha\) 之和至多为其高度，故只有有限种可能。\(E_\nu\) 也有限维，所以 \(M_\mu\) 有限维。

任取实际出现的权 \(\mu\)。处在 \(\mu+Q_+\) 中的权只有有限多个：对每个 \(\nu\in F\)，若 \(\mu+\beta=\nu-\gamma\)，则 \(\beta,\gamma\in Q_+\) 且 \(\beta+\gamma=\nu-\mu\)，高度受到固定上界约束。取这些权中的极大者 \(\eta\)，则全部正根向量都零化 \(M_\eta\)。任意非零 \(v\in M_\eta\) 就是所需向量。
\end{proof}

局部有限不表示整个模有限维。对 \(\mathfrak{sl}_2\)，Verma 模的 \(f\) 方向无限延伸，而 \(e\) 反复作用于任一向量最终为零。另一方面，去掉有限生成条件后，\(\bigoplus_{j\ge0}\mathbb C\) 虽有权分解且正根作用为零，零权空间却无限维。

\subsection{Verma 模与简单模}
\label{b5:subsec:2001:02}

\begin{proposition}\label{b5:prop:o-simple-objects}
设 \(\mathfrak g\) 为有限维复半单 Lie 代数，固定 Cartan 子代数 \(\mathfrak h\) 和正根集合，令 \(\mathcal O\) 为相应的范畴 O。对 \(\lambda\in\mathfrak h^*\)，记 \(M(\lambda)\) 为最高权 \(\lambda\) 的 Verma 模，\(L(\lambda)\) 为其唯一简单商。则这两个模都属于 \(\mathcal O\)，而 \(\mathcal O\) 的所有简单对象恰为这些 \(L(\lambda)\)，参数 \(\lambda\) 不重复。
\end{proposition}
\begin{proof}
由\cref{b5:prop:verma-pbw-universal}，\(M(\lambda)\) 是循环权模，权在 \(\lambda-Q_+\) 中且权空间有限维。对权为 \(\lambda-\beta\) 的向量，正根方向只能升到 \(\lambda-\beta+\gamma\le\lambda\)，其中 \(\gamma\in Q_+\) 的高度至多为 \(\beta\) 的高度；只有有限个这样的权。因此其 \(U(\mathfrak n^+)\) 轨道有限维。一般向量只有有限个权分量，结论仍成立。商模 \(L(\lambda)\) 继承三个条件。

若 \(S\) 是 \(\mathcal O\) 中的简单对象，\cref{b5:prop:o-weight-finiteness}给出最高权向量 \(v\)。子模 \(U(\mathfrak g)v\) 本身是最高权模，满足 \(\mathcal O\) 的条件，所以简单性使它等于 \(S\)。于是 Verma 泛性质和\cref{b5:thm:verma-simple-quotient}给出 \(S\cong L(\lambda)\)，并保证参数唯一。
\end{proof}

这里的分类允许任意复最高权。有限维简单模只占其中的支配整权部分。对于 \(\mathfrak{sl}_2\)，\eqref{b5:eq:sl2-verma-actions}说明：当 \(z\notin\mathbb Z_{\ge0}\) 时，\(M(z)\) 已经简单；当 \(z=m\ge0\) 为整数时，\(f^{m+1}v\) 生成它的唯一非零真子模。这一变化不能由有限维完全可约性解释，因为 Verma 模从来不是有限维模。

\subsection{交换范畴基本概念回顾}
\label{b5:subsec:2001:03}

一个范畴中的核与余核必须仍处在该范畴内，才可以在其中讨论正合列。包络代数的 Noether 性\footnote{诺特（Amalie Emmy Noether，1882-1935），德国数学家，发展抽象代数中的理想理论，并证明对称性与守恒律的联系。}保证有限生成模的子模仍有限生成；下面先证明这一点，再与 \(\mathcal O\) 的情形比较。

\begin{lemma}\label{b5:lem:enveloping-noetherian}
若 \(\mathfrak g\) 是有限维复 Lie 代数，则 \(U(\mathfrak g)\) 是左、右 Noether 环。因此有限生成 \(U(\mathfrak g)\) 模的每个子模仍有限生成。
\end{lemma}
\begin{proof}
PBW 给出 \(\operatorname{gr}U(\mathfrak g)\cong\operatorname{Sym}(\mathfrak g)\)，后者为有限个变量的多项式环，由 Hilbert 基定理 Noether。对左理想 \(I\) 取诱导滤过，其主符号构成分次理想 \(\operatorname{gr}I\)。选择有限个齐次生成元，并在 \(I\) 中提升为 \(u_1,\ldots,u_s\)。任取 \(u\in I\)，先用 \(u_i\) 的左倍数消去最高次主符号，所得差次数严格下降；按次数归纳可把 \(u\) 写成这些 \(u_i\) 的左线性组合。因此左理想有限生成。右侧同理。最后对有限秩自由模的秩归纳，再把一般有限生成模表示为自由模的商，即得子模结论。
\end{proof}

\begin{proposition}\label{b5:prop:o-abelian}
设 \(\mathfrak g\) 为有限维复半单 Lie 代数，固定 Cartan 子代数及正根集合，令 \(\mathcal O\) 为相应的范畴 O。则 \(\mathcal O\) 是交换范畴，核、余核、有限直和以及像均按通常的 \(U(\mathfrak g)\) 模构造。
\end{proposition}
\begin{proof}
若 \(f:M\to N\) 是模同态，核与商继承权分解：对一个向量的有限个权分量，选 \(h\in\mathfrak h\) 区分这些权，再用 \(h\) 的插值多项式提取各分量，故子模也按权分解。局部 \(\mathfrak n^+\) 有限性对子模和商显然保持；有限生成性对商保持，对核由\cref{b5:lem:enveloping-noetherian}保持。有限直和也满足三个条件。于是模范畴中的核、余核都在 \(\mathcal O\) 内，且余像到像的自然映射仍为同构，这正是交换范畴的判据。
\end{proof}

不能据此断言它在全部 \(U(\mathfrak g)\) 模中对任意扩张封闭。例如令 \(\mathfrak g=\mathfrak{sl}_2\)，在二维 \(\mathfrak b\) 模 \(E\) 上令 \(e=0\)、\(h=\lambda I+N\)，其中 \(N\ne0\)、\(N^2=0\)。诱导模 \(U(\mathfrak g)\otimes_{U(\mathfrak b)}E\) 是两个 \(M(\lambda)\) 的扩张；诱导正合性来自 PBW 的右自由性。但其中 \(h\) 不是半单算子，所以中间项不在 \(\mathcal O\)。讨论 \(\operatorname{Ext}_{\mathcal O}\) 时，允许的中间对象也必须属于 \(\mathcal O\)。

% !TeX root = ../../Book5.tex
% 纲要标题：### 20.2 中心作用
% 纲要位置：工作记录/计划纲要.md，第 4368 行。

\section{中心作用}
\label{b5:sec:2002}

% 纲要范围：
%
% * 无穷小特征标
% * 包络代数的中心
% * Harish–Chandra 同构的标准形式
%

Casimir 元在最高权模上给出一个标量，但不同最高权可能给出同一标量。中心的全部元素可以同时测量一个模，不过仍不能区分所有简单最高权模。本节将这种限制具体算出。

\subsection{无穷小特征标}
\label{b5:subsec:2002:01}

\begin{definition}\label{b5:def:infinitesimal-character}
设 \(\mathfrak g\) 是复 Lie 代数，\(Z=Z\bigl(U(\mathfrak g)\bigr)\)，\(M\ne0\) 是左 \(U(\mathfrak g)\) 模。若存在含幺代数同态 \(\chi:Z\to\mathbb C\)，使 \(zm=\chi(z)m\) 对所有 \(z\in Z,m\in M\) 成立，则称 \(M\) 具有\term[wuqiongxiao-tezhengbiao]{无穷小特征标}[infinitesimal character] \(\chi\)。这里所说的是包络代数中心的特征标，与群上取矩阵迹的特征标不同。
\end{definition}

\begin{proposition}\label{b5:prop:highest-central-character}
设 \(\mathfrak g\) 是有限维复半单 Lie 代数，固定 Cartan 子代数 \(\mathfrak h\) 和正根集合，\(\lambda\in\mathfrak h^*\)。设非零左 \(\mathfrak g\) 模 \(M=U(\mathfrak g)v\) 由权 \(\lambda\) 的最高权向量 \(v\) 生成。则 \(M\) 有无穷小特征标，即包络代数中心 \(Z(U(\mathfrak g))\) 在 \(M\) 上按一个含幺代数同态 \(\chi_\lambda:Z(U(\mathfrak g))\to\mathbb C\) 作用。这个同态只依赖 \(\lambda\)，所有同一最高权的非零商模具有同一个 \(\chi_\lambda\)。
\end{proposition}
\begin{proof}
设 \(v\) 为最高权生成元。中心元 \(z\) 与 Cartan 子代数交换，所以保持一维空间 \(M_\lambda=\mathbb Cv\)，从而 \(zv=\chi_\lambda(z)v\)。对 \(u\in U(\mathfrak g)\)，有 \(zuv=uzv=\chi_\lambda(z)uv\)，故这个标量作用于全模。和、积及单位元在 \(v\) 上的作用说明 \(\chi_\lambda\) 是含幺代数同态。最后每个最高权模都是 \(M(\lambda)\) 的商，故标量与所选商无关。
\end{proof}

这个论证利用最高权空间的一维性及其生成全模的性质，无须将有限维 Schur 引理推广到无限维模。

\subsection{包络代数的中心}
\label{b5:subsec:2002:02}

固定三角分解。PBW 将包络代数写成有序负根、Cartan、正根单项式的线性组合。保留没有正负根因子的单项式，得到线性投影
\[
\operatorname{pr}:U(\mathfrak g)\longrightarrow U(\mathfrak h)
=\operatorname{Sym}(\mathfrak h).
\]
对中心元 \(z\)，它的伴随 Cartan 权为零。含正根因子的项零化最高权向量；只含非零负根因子的项不可能有权零。因此
\begin{equation}\label{b5:eq:central-character-projection}
\chi_\lambda(z)=\operatorname{pr}(z)(\lambda).
\end{equation}
投影 \(\operatorname{pr}\) 在整个 \(U(\mathfrak g)\) 上不是代数同态，但在中心上的乘法性由这个标量公式推出：两个多项式在全部 \(\lambda\in\mathfrak h^*\) 上相同，就确实相等。

\begin{theorem}\label{b5:thm:sl2-enveloping-center}
设 \(\mathfrak g=\mathfrak{sl}_2(\mathbb C)\)，取基 \(e,f,h\)，满足 \([h,e]=2e\)、\([h,f]=-2f\)、\([e,f]=h\)，并记 \(C\) 为 Killing 型对应的 Casimir 元。令
\[
\Omega=h^2+2h+4fe=8C.
\]
则 \(Z\bigl(U(\mathfrak g)\bigr)=\mathbb C[\Omega]\)。在最高权为 \(z\in\mathbb C\) 的模上，\(\Omega\) 作用为 \(z(z+2)\)。
\end{theorem}
\begin{proof}
\(\Omega=8C\) 的中心性来自\cref{b5:prop:casimir-central}和\eqref{b5:eq:sl2-killing-casimir}，作用标量由最高向量直接算得。现在证明没有遗漏其他中心元。

若 \(u\) 是非零中心元，取 PBW 滤过中的最高次主符号 \(p\in\mathbb C[e,f,h]\)。对每个 \(x\in\mathfrak g\)，\([x,u]=0\) 使 \(\operatorname{ad}x\) 零化 \(p\)。由 \(h\) 的权，\(p\) 中 \(e,f\) 次数相等，所以 \(p=F(t,h)\)，\(t=ef\)。再用 \(\operatorname{ad}e\)，得到
\[
0=e\left(h\frac{\partial F}{\partial t}-2\frac{\partial F}{\partial h}\right).
\]
多项式环无零因子，故括号为零。换元 \(q=h^2+4t\)，则
\(\mathbb C[t,h]=\mathbb C[q,h]\)，上述导子在新坐标中就是 \(-2\partial/\partial h\)。因此 \(p\in\mathbb C[q]\)。由于 \(p\) 齐次且 \(q\) 次数为二，必有 \(p=a q^r\)。\(\Omega\) 的主符号正是 \(q\)，于是 \(u-a\Omega^r\) 仍中心，且滤过次数严格下降。归纳至常数，得到 \(u\in\mathbb C[\Omega]\)。不同幂的主符号次数不同，也说明 \(\Omega\) 没有多项式关系。
\end{proof}

因此 \(\mathfrak{sl}_2\) 的两个最高权给出相同无穷小特征标，当且仅当
\[
z(z+2)=w(w+2),\qquad
w=z\ \text{或}\ w=-z-2.
\]
例如 \(L(0)\) 与无限维简单模 \(L(-2)=M(-2)\) 有同一中心特征标，却不同构。中心参数识别的是一个有限反射群轨道，而不是单个最高权。

\subsection{Harish–Chandra 同构的标准形式}
\label{b5:subsec:2002:03}

令 \(\rho=\dfrac12\sum_{\alpha>0}\alpha\)，在 \(\mathfrak h^*\) 上定义移位作用
\[
w\mathbin{\cdot}\lambda=w(\lambda+\rho)-\rho.
\]
这确为群作用，因为先平移 \(\rho\)、作通常 Weyl 作用、再平移回来。引入 \(\rho\) 是为了把刚才的 \(z\mapsto-z-2\) 化成关于原点的反射：\(z+1\mapsto-(z+1)\)。

\begin{theorem}[Harish--Chandra 同构]\label{b5:thm:harish-chandra}
设 \(\mathfrak g\) 是有限维复半单 Lie 代数，固定 Cartan 子代数 \(\mathfrak h\) 和正根集合 \(\Phi^+\)，\(W\) 为相应 Weyl 群，\(\rho=\dfrac12\sum_{\alpha\in\Phi^+}\alpha\)。记 \(\operatorname{pr}:U(\mathfrak g)\to U(\mathfrak h)=\operatorname{Sym}(\mathfrak h)\) 为 PBW 三角分解中保留纯 Cartan 项的线性投影，把 \(\operatorname{Sym}(\mathfrak h)\) 识别为 \(\mathfrak h^*\) 上的多项式函数，\(W\) 按其自然作用作用于这些函数。则映射
\[
\gamma:Z\bigl(U(\mathfrak g)\bigr)\longrightarrow
\operatorname{Sym}(\mathfrak h)^W,
\qquad
\gamma(z)(\mu)=\operatorname{pr}(z)(\mu-\rho)
\]
是代数同构。对 \(\lambda\in\mathfrak h^*\)，记 \(\chi_\lambda\) 为最高权 \(\lambda\) 的 Verma 模的无穷小特征标。则对任意中心元 \(z\)，有 \(\chi_\lambda(z)=\gamma(z)(\lambda+\rho)\)。对 \(\lambda,\nu\in\mathfrak h^*\)，\(\chi_\lambda=\chi_\nu\) 当且仅当 \(\nu=w\mathbin{\cdot}\lambda=w(\lambda+\rho)-\rho\) 对某个 \(w\in W\) 成立。
\end{theorem}

\begin{proofsketch}
前面已经由最高权向量的作用证明 \(\operatorname{pr}\) 在中心上是代数同态，且 \(\chi_\lambda(z)=\operatorname{pr}(z)(\lambda)\)。先核对平移后的 \(W\) 不变性。取支配整权 \(\lambda\)，令 \(m_i=\lambda(h_i)\)。\cref{b5:lem:integrable-verma-quotient}的证明已核验 \(f_i^{m_i+1}v_\lambda\) 被全部正简单根向量零化；PBW 又保证它非零。因 \(\rho(h_i)=1\)，这个最高权向量的权恰为
\[
\lambda-(m_i+1)\alpha_i
=s_i(\lambda+\rho)-\rho=s_i\mathbin{\cdot}\lambda.
\]
因此 \(\rho\) 平移正好记录了已有奇异向量的最高权。它生成一个非零最高权子模，同一个中心元在母模和子模上作用的标量相同，故
\[
 \operatorname{pr}(z)(\lambda)
 =\operatorname{pr}(z)(s_i\mathbin{\cdot}\lambda).
\]
支配整权在 \(\mathfrak h^*\) 中 Zariski 稠密，而两端都是关于 \(\lambda\) 的多项式，所以等式对全部权成立。每个简单反射都满足这一式，故 \(\gamma(z)\) 属于 \(\operatorname{Sym}(\mathfrak h)^W\)。

再给包络代数取 PBW 滤过 \(F_dU(\mathfrak g)\)，给中心取诱导滤过 \(F_dZ=Z\cap F_dU(\mathfrak g)\)，给对称代数取次数滤过。若 \(z\in Z\) 的最高次主符号为 \(p\)，则 \([x,z]=0\) 对每个 \(x\in\mathfrak g\) 成立，所以 \(p\in\operatorname{Sym}(\mathfrak g)^{\mathfrak g}\)。主符号映射因而给出单射的分次代数同态
\[
 \operatorname{gr}Z\bigl(U(\mathfrak g)\bigr)
 \longrightarrow\operatorname{Sym}(\mathfrak g)^{\mathfrak g}.
\]
反过来，PBW 对称化
\[
 \mathrm{sym}:\operatorname{Sym}(\mathfrak g)\longrightarrow U(\mathfrak g)
\]
是滤过 \(\mathfrak g\)-模同构：伴随作用按 Leibniz 规则作用于各因子，因此与对排列取平均相容；在伴随分次上，对称化诱导恒等映射。若 \(p\) 是齐次不变元，则 \(\mathrm{sym}(p)\) 与全部 \(x\in\mathfrak g\) 交换，故为中心元，且其主符号为 \(p\)。这证明上述主符号映射满射，于是
\[
 \operatorname{gr}Z\bigl(U(\mathfrak g)\bigr)
 \cong\operatorname{Sym}(\mathfrak g)^{\mathfrak g}
\]
为分次代数同构；这里乘法的相容性由主符号映射给出。若 \(z\) 的 PBW 次数为 \(d\)、主符号为 \(p\)，则 \(\gamma(z)\) 的 \(d\) 次齐次部分由 \(p\) 沿三角分解的投影 \(\mathfrak g\to\mathfrak h\) 得到，因为 PBW 重排产生的交换子以及平移 \(\rho\) 都只改变较低次数项。将 \(\mathfrak h^*\) 作为零化 \(\mathfrak n^+\oplus\mathfrak n^-\) 的线性泛函嵌入 \(\mathfrak g^*\)，这个投影就是多项式函数的限制。Chevalley 限制定理给出
\[
 \operatorname{Sym}(\mathfrak g)^{\mathfrak g}
 \xrightarrow{\ \sim\ }
\operatorname{Sym}(\mathfrak h)^W.
\]
这里可以用已经建立的有限维最高权表示说明限制同构。记 Killing 型为 \(B\)。由\cref{b5:prop:root-killing-pairing}，\(B\) 及 \(B|_{\mathfrak h}\) 均非退化，给出同构
\[
 \begin{aligned}
 b_{\mathfrak g}:\mathfrak g&\xrightarrow{\sim}\mathfrak g^*,
 &x&\longmapsto B(x,-),\\
 b_{\mathfrak h}:\mathfrak h&\xrightarrow{\sim}\mathfrak h^*,
 &h&\longmapsto B(h,-)|_{\mathfrak h}.
 \end{aligned}
\]
前者与伴随及余伴随作用相容，后者与 \(W\) 作用相容。它们分别诱导 \(\operatorname{Sym}(\mathfrak g)\cong\operatorname{Sym}(\mathfrak g^*)\) 和 \(\operatorname{Sym}(\mathfrak h)\cong\operatorname{Sym}(\mathfrak h^*)\)。由于 \(B(\mathfrak h,\mathfrak n^+\oplus\mathfrak n^-)=0\)，这两份识别把上面的限制映射化为 \(\mathfrak g\) 上的多项式函数限制到 \(\mathfrak h\) 的映射。因此下面在 \(\operatorname{Sym}(\mathfrak g^*)\) 与 \(\operatorname{Sym}(\mathfrak h^*)\) 中证明限制同构，再通过 \(b_{\mathfrak g}\) 与 \(b_{\mathfrak h}\) 转回定理中的对称代数。若一个伴随不变多项式在 \(\mathfrak h\) 上为零，它在 \(\mathfrak h\) 的全部幂零伴随指数像上也为零；\cref{b5:thm:cartan-conjugacy}证明中的开像计算说明这些像含有非空 Zariski 开集，因此原多项式为零，得到单射。

为证满射，记 \(P\subseteq\mathfrak h^*\) 为整权格。对每个次数 \(m\ge0\)，\(\operatorname{Sym}^m(\mathfrak h^*)\) 由 \(\nu^m\)、\(\nu\in P\)，线性生成。事实上，设
\[
 \ell\in\bigl(\operatorname{Sym}^m(\mathfrak h^*)\bigr)^*,
 \qquad \ell(\nu^m)=0\quad(\nu\in P).
\]
函数 \(\nu\mapsto\ell(\nu^m)\) 是 \(\mathfrak h^*\) 上的 \(m\) 次齐次多项式。以基本权为坐标基，\(P\) 就是整数格；逐变量使用无限多个整数根，便知该多项式恒为零。展开 \(\ell((\sum_i t_i\omega_i)^m)\) 并比较系数，得到 \(\ell\) 在全部次数 \(m\) 的单项式上为零，故 \(\ell=0\)。有限维线性代数于是给出所述生成性。再对 \(W\) 平均，便知不变部分由轨道和 \(\sum_{\nu\in W\lambda}\nu^m\) 生成。

支配整权 \(\lambda\) 的简单模特征标具有三角展开
\[
 \operatorname{ch}L(\lambda)
 =\sum_{\nu\in W\lambda}e^\nu
 +\sum_{\mu<\lambda}a_{\lambda\mu}
   \sum_{\nu\in W\mu}e^\nu,
\]
其中 \(\mu\) 支配且 \(\lambda-\mu\) 是简单根的非负整数组合。低于固定 \(\lambda\) 的支配权有限：正定根内积给出 \(\|\mu\|\le\|\lambda\|\)，而权格离散。按这个有限偏序归纳，每个形式轨道和都是有限维模特征标的线性组合。记 \(L(\lambda)\) 的表示为 \(\pi_\lambda:\mathfrak g\to\operatorname{End}_{\mathbb C}L(\lambda)\)。在线性映射 \(e^\nu\mapsto\nu^m\) 下，其特征标的像恰是
\(h\mapsto\operatorname{tr}(\pi_\lambda(h)^m)\)。它延拓为 \(x\mapsto\operatorname{tr}(\pi_\lambda(x)^m)\)，这是 \(\mathfrak g\) 上的伴随不变多项式，因为其沿 \([y,x]\) 的导数是交换子的迹，等于零。于是所有所需轨道和均能延拓，限制映射满射。

故 \(\gamma\) 的伴随分次映射 \(\operatorname{gr}\gamma\) 是同构。对次数归纳，依次提升最高次不变多项式并减去已提升项，得到 \(\gamma\) 满射；若 \(\gamma(z)=0\)，其最高次符号为零，按次数下降得到 \(z=0\)，故它也单射。

最后，有限群的两个不同轨道可由不变多项式区分：先在两条有限轨道上作多项式插值，使取值分别为零和一，再对 \(W\) 平均。于是全部中心元在两个最高权模上取相同标量，当且仅当 \(\lambda+\rho\) 与 \(\nu+\rho\) 属于同一 \(W\) 轨道。移位不变性调用前章已验证的奇异向量；限制同构则由上面的稠密性、特征标三角性和迹多项式论证得到。限制论证可对照\cite[Section 23.1]{Humphreys1972LieAlgebras}，滤过比较和移位不变性见\cite[Section 8.8, Theorems 8.50, 8.52 and 8.54]{Kirillov2008LieGroups}。
\end{proofsketch}


\begin{example}\label{b5:exm:hc-sl3-orbit}
对 \(\mathfrak{sl}_3\)，把 \(\mathfrak h^*\) 写成 \(a+b+c=0\) 的三元组，\(W=S_3\) 排列三个坐标。其不变多项式由
\(a^2+b^2+c^2\) 与 \(a^3+b^3+c^3\) 生成：对称多项式基本定理把任一不变式写为初等对称式的多项式，而第一初等对称式为零，另外两个分别为
\[
-\frac12(a^2+b^2+c^2),\qquad \frac13(a^3+b^3+c^3).
\]
二次 Casimir 只看到第一个参数，完整中心还需要三次参数。这里的取值点是 \(\lambda+\rho\)，不能漏掉平移。
\end{example}

无穷小特征标也提示了如何分组研究 \(\mathcal O\) 中的扩张。若中心并非按纯标量作用，可以要求每个 \(z-\chi(z)\) 在每个向量上局部幂零；这称为广义中心特征标条件。扩张可能保留广义特征标而失去纯标量作用，因此两个条件不能混写。

% !TeX root = ../../Book5.tex
% 纲要标题：### 20.3 Lie 代数上同调
% 纲要位置：工作记录/计划纲要.md，第 4374 行。

\section{Lie 代数上同调}
\label{b5:sec:2003}

% 纲要范围：
%
% * Chevalley–Eilenberg 复形
% * 低阶上同调、导子及扩张
% * Whitehead 引理与 Levi 分解的证明联系
% * 在本节证明 Whitehead 第一、第二引理，并沿可解根的交换理想归纳，完成第12.4节预告的 Levi 一般存在性与共轭性；不能只以交换核的单次修正代替一般证明
%

本节回到\secref{b5:sec:1204}留下的两个问题：线性截面怎样修正成 Lie 代数同态，不同的半单补怎样互相共轭。误差分别满足二阶和一阶余圈恒等式。先建立这些恒等式所属的复形，再证明消失结果，最后处理一般可解根中的逐层修正。

\subsection{Chevalley–Eilenberg 复形}
\label{b5:subsec:2003:01}

\begin{definition}\label{b5:def:chevalley-eilenberg}
设 \(\mathfrak a\) 为复 Lie 代数，\(M\) 为左 \(\mathfrak a\) 模，\(q\ge0\) 为整数。令
\(C^q(\mathfrak a,M)=\operatorname{Hom}_{\mathbb C}(\Lambda^q\mathfrak a,M)\)，并置 \(C^0=M\)。定义升次映射
\begin{align*}
(\mathrm{d}\varphi)(x_0,\ldots,x_q)
={}&\sum_{i=0}^q(-1)^i x_i\cdot
\varphi(x_0,\ldots,\widehat{x_i},\ldots,x_q)\\
&+\sum_{0\le i<j\le q}(-1)^{i+j}
\varphi\bigl([x_i,x_j],x_0,\ldots,\widehat{x_i},\ldots,
\widehat{x_j},\ldots,x_q\bigr).
\end{align*}
帽号表示删去对应变量。此上链复形称为\term[Chevalley-Eilenberg-fuxing]{Chevalley--Eilenberg 复形}[Chevalley--Eilenberg complex]；其上同调
\(H^q(\mathfrak a,M)=\ker \mathrm{d}^q/\operatorname{im}\mathrm{d}^{q-1}\)
称为\term[lie-daishu-shangtongdiao]{Lie 代数上同调}[Lie algebra cohomology]，其中约定 \(\mathrm{d}^{-1}=0\)。
\end{definition}
\symidx{\(C^q(\mathfrak a,M),H^q(\mathfrak a,M)\)：Chevalley--Eilenberg 余链空间与 Lie 代数上同调}

\begin{proposition}\label{b5:prop:ce-differential}
设 \(\mathfrak a\) 为复 Lie 代数，\(M\) 为左 \(\mathfrak a\) 模，\(C^q(\mathfrak a,M)=\operatorname{Hom}_{\mathbb C}(\Lambda^q\mathfrak a,M)\)，\(\mathrm{d}\) 为 Chevalley--Eilenberg 微分。则 \(\mathrm{d}^2=0\)。对每个整数 \(q\ge0\)，\(C^q(\mathfrak a,-)\) 是正合函子；短正合系数列由此诱导上同调长正合列。
\end{proposition}
\begin{proof}
展开 \(\mathrm{d}(\mathrm{d}\varphi)\)。两次各取一个作用项时，每对 \(i<j\) 合并为带符号的
\(x_i\cdot(x_j\cdot v)-x_j\cdot(x_i\cdot v)\)；先取括号再让该括号作用的项给出相反符号的 \([x_i,x_j]\cdot v\)。它们由模公理抵消。一次作用与一个不含该变量的括号对应两种先后次序；删除变量的次序相差一次换位，符号相反，故两两抵消。两对互不相交的括号也按交换先后次序两两抵消。剩余项只涉及三个变量，其内部组合为
\[
\bigl[[x_i,x_j],x_k\bigr]+\bigl[[x_j,x_k],x_i\bigr]+\bigl[[x_k,x_i],x_j\bigr],
\]
在 \(\varphi\) 的同一位置出现，外部排列符号相同；Jacobi 恒等式使之为零。这覆盖了展开中的全部项，故 \(\mathrm{d}^2=0\)。

忘掉 Lie 代数作用后，系数短正合列是分裂的向量空间短正合列。因此余链函子 \(C^q(\mathfrak a,-)\) 正合，这里不要求 \(\mathfrak a\) 有限维。逐次提升一个余圈、对提升取微分，再在核中读取结果，就得到连接同态；不同提升在核中相差一个余链，其微分使结果相差余边，故连接同态良定义。核与像的追逐与一般复形的长正合列构造相同。
\end{proof}

低次数公式值得单独写出：
\begin{align*}
(\mathrm{d}m)(x)&=x\cdot m,\\
(\mathrm{d}b)(x,y)&=x\cdot b(y)-y\cdot b(x)-b\bigl([x,y]\bigr),\\
(\mathrm{d}c)(x,y,z)&=x\cdot c(y,z)-y\cdot c(x,z)+z\cdot c(x,y)\\
&\quad-c\bigl([x,y],z\bigr)+c\bigl([x,z],y\bigr)
-c\bigl([y,z],x\bigr).
\end{align*}
它们与\cref{b5:prop:levi-abelian-obstruction}中的截面误差完全一致。

\subsection{低阶上同调、导子及扩张}
\label{b5:subsec:2003:02}

由定义，\(H^0(\mathfrak a,M)=M^{\mathfrak a}\)，而 \(H^1\) 是满足
\(b\bigl([x,y]\bigr)=x\cdot b(y)-y\cdot b(x)\)
的线性映射模掉 \(b(x)=x\cdot m\) 的映射。系数为伴随模 \(\mathfrak a\) 时，前者正是导子；后者等于 \(x\mapsto[x,m]=-\operatorname{ad}(m)x\)，所以
\[
H^1(\mathfrak a,\mathfrak a)
\cong\operatorname{Der}(\mathfrak a)/\operatorname{ad}(\mathfrak a).
\]

\begin{proposition}\label{b5:prop:lie-h2-extensions}
设 \(\mathfrak a\) 为复 Lie 代数，\(M\) 为其左模。交换核扩张
\[
0\longrightarrow M\longrightarrow\mathfrak e\longrightarrow\mathfrak a
\longrightarrow0
\]
的核 \(M\) 是满足 \([M,M]=0\) 的 Lie 理想，诱导的 \(\mathfrak a\) 作用固定为所给作用，并且扩张同构要求在 \(M\) 和 \(\mathfrak a\) 上均为恒等。则其同构类与 \(H^2(\mathfrak a,M)\) 一一对应，零类对应半直积扩张。
\end{proposition}
\begin{proof}
选商映射 \(\pi:\mathfrak e\to\mathfrak a\) 的线性截面 \(s\)，令
\[
x\cdot m=\bigl[s(x),m\bigr],\qquad
c(x,y)=\bigl[s(x),s(y)\bigr]-s\bigl([x,y]\bigr).
\]
因为核 \(M\) 交换，前式不依赖提升的选择，且 Jacobi 恒等式使之满足模公理；这里正取命题要求固定的作用。第二式的像落在 \(\ker\pi=M\)。将它代入 \(\mathfrak e\) 的 Jacobi 恒等式，得到
\[
\begin{aligned}
0={}&x\cdot c(y,z)-y\cdot c(x,z)+z\cdot c(x,y)\\
&-c\bigl([x,y],z\bigr)+c\bigl([x,z],y\bigr)
-c\bigl([y,z],x\bigr)=(\mathrm{d}c)(x,y,z).
\end{aligned}
\]
若 \(s'=s+b\)、\(b:\mathfrak a\to M\) 线性，则展开括号并用 \([M,M]=0\)，有
\[
c'(x,y)-c(x,y)
=x\cdot b(y)-y\cdot b(x)-b\bigl([x,y]\bigr)
=(\mathrm{d}b)(x,y).
\]
这些正是\cref{b5:prop:levi-abelian-obstruction}中的截面障碍计算；该处虽在有限维条件下陈述，逐项计算并未使用维数限制，因此同样适用于当前的一般扩张。反过来，对余圈 \(c\) 在 \(M\oplus\mathfrak a\) 上定义
\[
\bigl[(m,x),(n,y)\bigr]
=\bigl(x\cdot n-y\cdot m+c(x,y),[x,y]\bigr).
\]
反对称性成立。涉及两个 \(M\) 变量的 Jacobi 式为零；一个 \(M\) 变量的情形正是模公理；三个 \(\mathfrak a\) 变量的 \(M\) 分量为 \(\mathrm{d}c\)，另一个分量为 \(\mathfrak a\) 的 Jacobi 式。因此这是 Lie 代数。

改变截面等价于改变坐标 \((m,x)\mapsto\bigl(m+b(x),x\bigr)\)，相应余圈相差 \(\mathrm{d}b\)。任何在两端为恒等的扩张同构都具有这种坐标形式，故余圈类相等恰好等价于扩张同构。最后 \(c=0\) 时上述括号是半直积；一个类为零的扩张可通过改换截面使 \(c=0\)。
\end{proof}

例如 \(\mathfrak a=\mathbb Cx\oplus\mathbb Cy\) 交换、\(M=\mathbb Cz\) 平凡时，所有微分为零，故 \(H^2(\mathfrak a,M)\) 一维。取 \(c(x,y)=z\)，得到 \([x,y]=z\)、\(z\) 中心的三维 Heisenberg Lie 代数。这一扩张不分裂，因为若有交换补，提升 \(x,y\) 后无论加上多少中心分量，括号仍是 \(z\)。

\subsection{Whitehead 引理与 Levi 分解}
\label{b5:subsec:2003:03}

\begin{theorem}[Whitehead 第一引理]\label{b5:thm:whitehead-first}
设 \(\mathfrak s\) 为有限维复半单 Lie 代数，\(M\) 为有限维复 \(\mathfrak s\) 模。则 \(H^1(\mathfrak s,M)=0\)。
\end{theorem}
\begin{proof}
一阶余圈条件正是\cref{b5:lem:semisimple-crossed-map}中的交叉映射条件。该引理已经用 Killing Casimir 证明存在 \(m\in M\) 使 \(b(x)=x\cdot m\)，而这里的右端就是 \((\mathrm{d}m)(x)\)。故所有一阶余圈都是余边。引理要求的半单性与有限维性恰为当前假设。
\end{proof}

第二引理还需要把 Casimir 修正推广到更高阶。下面给出可直接检验的链同伦公式；它同时说明非平凡简单系数与平凡系数为什么应分别处理。

\begin{lemma}\label{b5:lem:casimir-ce-homotopy}
设 \(\mathfrak s\) 为有限维复半单 Lie 代数，\(M\) 为左 \(\mathfrak s\) 模，\((x_i)\) 和 \((x^i)\) 为 \(\mathfrak s\) 的两组 Killing 对偶基。记 \(C^q(\mathfrak s,M)=\operatorname{Hom}_{\mathbb C}(\Lambda^q\mathfrak s,M)\)，\(\mathrm{d}\) 为 Chevalley--Eilenberg 微分，\(C_M:m\mapsto\sum_i x_i\cdot(x^i\cdot m)\) 为 Killing Casimir 在 \(M\) 上的作用。对整数 \(q\ge1\)，在 \(C^q(\mathfrak s,M)\) 上定义
\[
(H\varphi)(y_1,\ldots,y_{q-1})
=\sum_i x_i\cdot\varphi(x^i,y_1,\ldots,y_{q-1}),
\]
且 \(H=0\) 于次数零。则 \(\mathrm{d}H+H\mathrm{d}=C_M\)，其中右边只让 Killing Casimir 作用于 \(\varphi\) 的取值。若 \(C_M\) 可逆，则所有 \(H^q(\mathfrak s,M)\) 均为零。
\end{lemma}
\begin{proof}
先在一阶余链 \(b\) 上看抵消的内容。此时 \(Hb=\sum_i x_i\cdot b(x^i)\)，直接使用低次数微分得到
\[
\begin{aligned}
(\mathrm{d}Hb+H\mathrm{d}b)(y)
={}&\sum_i x_i\cdot\bigl(x^i\cdot b(y)\bigr)\\
&+\sum_i\left([y,x_i]\cdot b(x^i)
+x_i\cdot b([y,x^i])\right).
\end{aligned}
\]
第二行正是把不变张量恒等式
\(\sum_i([y,x_i]\otimes x^i+x_i\otimes[y,x^i])=0\)
送入映射 \(a\otimes z\mapsto a\cdot b(z)\)，故为零。留下的第一行是 Casimir 对取值 \(b(y)\in M\) 的作用；它没有同时对变量 \(y\) 作伴随作用。

用 \(R_x\) 表示只对取值施加 \(x\) 作用，用 \(i_x\) 表示在交错形式第一个变量插入 \(x\)。若 \(a_j\) 是任意基，\(\epsilon^j\) 表示与其对偶线性形式作外乘，则直接把微分的作用项和括号项分开，得到
\begin{align*}
\mathrm{d}R_x-R_x\mathrm{d}&=\sum_j\epsilon^j R_{[a_j,x]},\\
\mathrm{d}i_x+i_x\mathrm{d}&=L_x
=R_x-\sum_j\epsilon^j i_{[x,a_j]}.
\end{align*}
第二式的右边在 \(q\) 个变量上就是
\(x\cdot\varphi(y_1,\ldots,y_q)
-\sum_r\varphi\bigl(y_1,\ldots,[x,y_r],\ldots,y_q\bigr)\)：微分中未涉及插入变量的项因插入位置移动而抵消，留下这一式。

因为 \(H=\sum_i R_{x_i}i_{x^i}\)，代入上面两式得到
\begin{align*}
\mathrm{d}H+H\mathrm{d}
={}&\sum_i R_{x_i}R_{x^i}\\
&+\sum_{i,j}\epsilon^j
\bigl(R_{[a_j,x_i]}i_{x^i}+R_{x_i}i_{[a_j,x^i]}\bigr).
\end{align*}
第二行由\eqref{b5:eq:casimir-tensor-invariance}为零，第一行就是 \(C_M\)。\(C_M\) 为模自同态，因而与 \(\mathrm{d}\) 交换。若 \(C_M\) 可逆，且 \(\varphi\in C^q(\mathfrak s,M)\) 满足 \(\mathrm{d}\varphi=0\)，则当 \(q\ge1\) 时，
\(\varphi=\mathrm{d}(C_M^{-1}H\varphi)\)，所以 \(\varphi\) 是余边。当 \(q=0\) 时，\(H\varphi=0\)，故
\[
C_M\varphi=H(\mathrm{d}\varphi)=0;
\]
由 \(C_M\) 可逆得 \(\varphi=0\)，所以零次上同调也为零。
\end{proof}

\begin{theorem}[Whitehead 第二引理]\label{b5:thm:whitehead-second}
设 \(\mathfrak s\) 为有限维复半单 Lie 代数，\(M\) 为有限维复 \(\mathfrak s\) 模。则 \(H^2(\mathfrak s,M)=0\)。
\end{theorem}
\begin{proof}
由\cref{b5:thm:weyl-complete-reducibility}，\(M\) 是有限个简单模的直和；上同调与这一有限直和相容。非平凡简单分量上的 Killing Casimir 是非零标量，由\cref{b5:lem:casimir-positive-simple}和\cref{b5:lem:casimir-ce-homotopy}，这些分量的二阶上同调为零。因此只需处理平凡系数 \(\mathbb C\)。\(\mathfrak s=0\) 时结论直接成立，以下允许用非退化 Killing 型 \(B\)。

取交错\(2\)-余圈 \(\omega:\Lambda^2\mathfrak s\to\mathbb C\)。由 \(B\) 非退化，存在唯一线性算子 \(D\) 满足
\(B(Dx,y)=\omega(x,y)\)。交错性说明 \(D\) 关于 \(B\) 斜对称。平凡系数的余圈式给出
\[
\omega\bigl([x,y],z\bigr)=\omega\bigl(x,[y,z]\bigr)+\omega\bigl(y,[z,x]\bigr).
\]
结合 Killing 型的不变性，右边等于
\(B\bigl([Dx,y],z\bigr)+B\bigl([x,Dy],z\bigr)\)。非退化性因此给出
\(D[x,y]=[Dx,y]+[x,Dy]\)，即 \(D\) 是导子。

将\cref{b5:thm:whitehead-first}用于有限维伴随模 \(\mathfrak s\)，得所有导子都是内导子。于是 \(D=\operatorname{ad}a\) 对某个 \(a\in\mathfrak s\) 成立。令线性形式 \(b(x)=-B(a,x)\)，则平凡系数时
\[
(\mathrm{d}b)(x,y)=-b\bigl([x,y]\bigr)=B\bigl(a,[x,y]\bigr)
=B\bigl([a,x],y\bigr)=\omega(x,y).
\]
故平凡系数的每个\(2\)-余圈也是余边，完成证明。
\end{proof}

这两个结果合称\term[Whitehead-yinli]{Whitehead 引理}[Whitehead lemmas]。它们不声称半单 Lie 代数在任意次数、任意系数下的上同调都消失；例如平凡系数的三阶上同调可以非零。有限维假设也不能从上述 Casimir 分解中删除。

\begin{theorem}[Levi--Malcev 定理]\label{b5:thm:levi-malcev}
设 \(\mathfrak g\) 为有限维复 Lie 代数，\(\mathfrak r=\operatorname{rad}\mathfrak g\)。则存在半单子代数 \(\mathfrak s\)，使
\(\mathfrak g=\mathfrak r\rtimes\mathfrak s\)。任意两个这样的子代数可由若干
\(\exp(\operatorname{ad}z)\)、\(z\in[\mathfrak g,\mathfrak r]\)
的有限乘积互相变换；这些指数均为有限多项式。
\end{theorem}
\begin{proof}
先说明指数为何只有有限项。记 \(\mathfrak n=\operatorname{nilrad}\mathfrak g\)。由\cref{b5:prop:radical-commutator-nilpotent}，\(z\in[\mathfrak g,\mathfrak r]\) 属于 \(\mathfrak n\)。由于 \(\mathfrak n\) 是幂零理想，
\[
(\operatorname{ad}z)^j(\mathfrak g)\subseteq\gamma_j(\mathfrak n)
\qquad(j\ge1).
\]
确实，第一次括号落在 \(\mathfrak n\)，以后每次与 \(z\in\mathfrak n\) 取括号，都进入下降中心列的下一项。这些项最终为零，故 \(\operatorname{ad}z\) 在整个 \(\mathfrak g\) 上幂零，\(\exp(\operatorname{ad}z)\) 是有限多项式。

先证存在性，对 \(\dim\mathfrak r\) 归纳。若 \(\mathfrak r=0\)，取 \(\mathfrak s=\mathfrak g\)。否则取 \(\mathfrak r\) 导出列的最后一个非零项 \(A\)。它是非零交换理想；导出列各项被保持 \(\mathfrak r\) 的导子保持，故 \(A\) 也是 \(\mathfrak g\) 的理想。

由\cref{b5:prop:radical-quotient-semisimple}，\(\mathfrak g/A\) 的可解根为 \(\mathfrak r/A\)。归纳假设给出该商中的 Levi 子代数 \(\overline{\mathfrak s}\)。设 \(\mathfrak e\) 是其在 \(\mathfrak g\) 中的逆像，则
\[
0\longrightarrow A\longrightarrow\mathfrak e
\longrightarrow\overline{\mathfrak s}\longrightarrow0
\]
有交换核。商通过括号在 \(A\) 上作用，与提升选择无关。由\cref{b5:thm:whitehead-second}，这个扩张的二阶类为零；\cref{b5:prop:lie-h2-extensions}便给出 Lie 代数截面。其像 \(\mathfrak s\) 同构于 \(\overline{\mathfrak s}\)，所以半单。商中 \(\overline{\mathfrak s}\) 与 \(\mathfrak r/A\) 相交为零，且截面像与 \(A\) 相交为零，故 \(\mathfrak s\cap\mathfrak r=0\)。维数或满射到 \(\mathfrak g/\mathfrak r\) 又给出 \(\mathfrak g=\mathfrak r+\mathfrak s\)。

再证共轭性，同样对 \(\dim\mathfrak r\) 归纳。取两个 Levi 子代数 \(\mathfrak s_1,\mathfrak s_2\)；在上述商 \(\mathfrak g/A\) 中，它们的像都是 Levi 子代数。归纳假设把两个像用若干指数共轭起来。商中的每个共轭元属于
\([\mathfrak g/A,\mathfrak r/A]\)，可提升为 \([\mathfrak g,\mathfrak r]\) 中的元素。由刚才说明的幂零性，提升后的指数是良定义的自同构，并诱导原商中的指数。对 \(\mathfrak s_1\) 先施加这些自同构后，可以假定二者在商中有相同的像。

于是它们都位于同一个子代数 \(\mathfrak e=\mathfrak s_1+A\) 中，并且都与 \(A\) 互补。以 \(\mathfrak s_1\) 为基准，\(\mathfrak s_2\) 写成图
\[
\bigl\{\,x+b(x)\mid x\in\mathfrak s_1\,\bigr\},
\qquad b:\mathfrak s_1\longrightarrow A.
\]
\(A\) 交换使子代数条件化为 \(b\bigl([x,y]\bigr)=x\cdot b(y)-y\cdot b(x)\)。由第一引理，\(b(x)=x\cdot m\) 对某个 \(m\in A\) 成立。

还须保证共轭元处在定理要求的 \([\mathfrak g,\mathfrak r]\) 中。由 Weyl 完全可约性，有限维 \(\mathfrak s_1\) 模 \(A\) 分解为平凡部分与非平凡简单部分之和，后者恰为 \([\mathfrak s_1,A]\)。把 \(m\) 的平凡分量删去不改变 \(b\)，故可取
\(m\in[\mathfrak s_1,A]\subseteq[\mathfrak g,\mathfrak r]\)。由于 \([m,\mathfrak g]\subseteq A\) 且 \([m,A]=0\)，有 \((\operatorname{ad}m)^2=0\)。因而
\[
\exp(-\operatorname{ad}m)x=x-[m,x]=x+x\cdot m=x+b(x),
\]
这个最后的指数把 \(\mathfrak s_1\) 送到 \(\mathfrak s_2\)。连同已经提升的有限个指数，得到一般共轭性。存在性和共轭性的每次归纳都严格降低可解根维数，故过程终止。
\end{proof}

这就证明了\cref{b5:thm:levi-malcev-preview}。二阶上同调消失使每一层交换核的扩张分裂，一阶上同调消失给出同一层中两个补的共轭；沿可解根归纳，便得到任意有限维复 Lie 代数的 Levi 分解与共轭性。

% !TeX root = ../../Book5.tex
% 纲要标题：### 20.4 BGG 方法
% 纲要位置：工作记录/计划纲要.md，第 4381 行。

\section{BGG 方法}
\label{b5:sec:2004}

% 纲要范围：
%
% * BGG 消解的典型版本
% * Ext 与最高权范畴
% * 几何表示论所需的进一步工具
%
% ---
%

取最高权模的简单商时，要商去较低最高权生成的子模。若出现多个这样的子模，还须考虑它们的交及相互关系。BGG 消解通过一个正合复形逐次描述这些关系。

\subsection{BGG 消解的典型版本}
\label{b5:subsec:2004:01}

\begin{theorem}[BGG 消解]\label{b5:thm:bgg-resolution}
设 \(\mathfrak g\) 为有限维复半单 Lie 代数，固定 Cartan 子代数 \(\mathfrak h\) 及正根集合 \(\Phi^+\)，\(\lambda\in\mathfrak h^*\) 为相应的支配整权，\(W\) 为 Weyl 群，\(\rho=\dfrac12\sum_{\alpha\in\Phi^+}\alpha\)。记 \(M(\mu)\) 为最高权 \(\mu\in\mathfrak h^*\) 的 Verma 模，\(L(\mu)\) 为其唯一简单商，\(\ell(w)\) 为 \(w\in W\) 的简单反射长度，\(w_0\) 为最长元，\(N=\ell(w_0)\)。对 \(0\le j\le N\)，记
\[
D_j=\bigoplus_{\ell(w)=j}M(w\mathbin{\cdot}\lambda),
\qquad w\mathbin{\cdot}\lambda=w(\lambda+\rho)-\rho.
\]
可以选择模同态 \(D_j\to D_{j-1}\)，使
\[
0\longrightarrow D_N\longrightarrow\cdots\longrightarrow D_1
\longrightarrow D_0\longrightarrow L(\lambda)\longrightarrow0
\]
正合。它称为\term[BGG-xiaojie]{BGG 消解}[Bernstein--Gelfand--Gelfand resolution]。
\end{theorem}
\symidx{\(w\mathbin{\cdot}\lambda\)：Weyl 群在权上的移位作用}

\begin{proofsketch}
先取 \(\lambda=0\)。令 \(\mathfrak b=\mathfrak h\oplus\mathfrak n_+\)，以 \(\mathfrak b\) 的伴随作用把 \(\mathfrak g/\mathfrak b\) 看成 \(\mathfrak b\) 模。考虑相对 Chevalley--Eilenberg 复形的各项
\[
 E_j=U(\mathfrak g)\otimes_{U(\mathfrak b)}
 \bigwedge^j(\mathfrak g/\mathfrak b),
 \qquad 0\le j\le N.
\]
以 \(\bar x=x+\mathfrak b\) 记商空间中的类。前两次微分具体为
\[
\begin{aligned}
d_1(u\otimes\bar x)&=ux\otimes1,\\
d_2(u\otimes\bar x\wedge\bar y)
&=ux\otimes\bar y-uy\otimes\bar x-u\otimes\overline{[x,y]}.
\end{aligned}
\]
括号取在 \(\mathfrak g\) 中，再取商类；\(\mathfrak g/\mathfrak b\) 本身并非商 Lie 代数。若把 \(x\) 换成 \(x+b\)、\(b\in\mathfrak b\)，第一式的差为 \(ub\otimes1=u\otimes b\cdot1=0\)，第二式的差为
\[
ub\otimes\bar y-u\otimes\overline{[b,y]}=0,
\]
因为张量积是在 \(U(\mathfrak b)\) 上取的，且 \(b\cdot\bar y=\overline{[b,y]}\)。因此提升改变造成的作用项与括号项恰好抵消。一般微分同样由删去一个外幂因子的作用项与两两括号项组成，Jacobi 恒等式使相邻微分复合为零。PBW 分解 \(\mathfrak g=\mathfrak n_-\oplus\mathfrak b\) 将此复形作为向量空间识别为 \(\mathfrak n_-\) 的标准消解；给 \(U(\mathfrak n_-)\) 取 PBW 滤过后，其分次复形是对称代数的 Koszul 消解，所以除了末端的平凡模外没有同调。

这里还要使用广义中心特征标的分解。对 \(M\in\mathcal O\)，以 \(M_\chi\) 记具有广义中心特征标 \(\chi\) 的向量组成的子模，则
\[
 M=\bigoplus_\chi M_\chi,
 \qquad \operatorname{pr}_\chi:M\longmapsto M_\chi
\]
只有有限多个非零项，且 \(\operatorname{pr}_\chi\) 为正合函子。其依据是中心保持每个有限维权空间，交换算子的广义共同特征空间分解与模同态相容；有限生成性又将出现的特征标限制为有限多个。对短正合列逐权作这一分解，便得到投影的正合性。这一分解见\cite[Section 1.12]{Humphreys2008BGGCategoryO}。\footnote{本节关于中心分解、张量滤过与 BGG 消解证明的定位按 Humphreys 的 2008 年作者稿核对，所引节号与 AMS 版目录中的同名各节一致。}

这些 \(E_j\) 具有 Verma 滤过，投影 \(\operatorname{pr}_{\chi_0}\) 因而仍给出正合复形。\(\bigwedge^j(\mathfrak g/\mathfrak b)\) 的权是 \(j\) 个不同负根的和；外幂权引理说明，其中与零权处于同一移位 \(W\) 轨道的，恰为 \(w\mathbin{\cdot}0=w\rho-\rho\)、\(\ell(w)=j\)，各出现一次。于是投影后的第 \(j\) 项具有每个 \(M(w\mathbin{\cdot}0)\) 恰一次的 Verma 滤过。为使它成为直和，需要 Verma 扩张判据的下述推论：若 \(\mu\) 为支配整权，\(x,y\in W\) 满足 \(\ell(x)=\ell(y)\)，则
\[
 \operatorname{Ext}^1_{\mathcal O}
 \bigl(M(x\mathbin{\cdot}\mu),M(y\mathbin{\cdot}\mu)\bigr)=0.
\]
事实上，非零扩张要求 \(y<x\) 成立于 Bruhat 次序，因而 \(\ell(y)<\ell(x)\)；该判据见\cite[Section 6.5]{Humphreys2008BGGCategoryO}。逐次分裂滤过中的短正合列，就得到 \(\lambda=0\) 时的 \(D_j\)。

例如取 \(\mathfrak{sl}_3\)，简单根记为 \(\alpha,\beta\)。商空间 \(\mathfrak g/\mathfrak b\) 的三个权为 \(-\alpha,-\beta,-\alpha-\beta\)。外幂次数零至三的权依次是
\[
\begin{aligned}
j=0:&\quad 0,\\
j=1:&\quad-\alpha,\ -\beta,\ -\alpha-\beta,\\
j=2:&\quad-\alpha-\beta,\ -2\alpha-\beta,\ -\alpha-2\beta,\\
j=3:&\quad-2\alpha-2\beta.
\end{aligned}
\]
其中零权的移位轨道为
\(0,-\alpha,-\beta,-2\alpha-\beta,-\alpha-2\beta,-2\alpha-2\beta\)。所以中心投影在次数一删去 \(-\alpha-\beta\) 对应的 Verma 因子，在次数二也删去同权的因子，保留的因子数依次为 \(1,2,2,1\)。这首先是投影后的 Verma 滤过；还要调用刚才的同长度 \(\operatorname{Ext}^1\) 消失，才可把每一项写成这些 Verma 模的直和。中心投影本身不提供这一步分裂。

一般支配整权 \(\lambda\) 可从这一复形出发，先逐项张量有限维模 \(L(\lambda)\)，再施行 \(\operatorname{pr}_{\chi_\lambda}\)。这两步均正合，末端变为 \(L(\lambda)\)。张量与诱导的标准恒等式给出 \(M(w\mathbin{\cdot}0)\otimes L(\lambda)\) 的 Verma 滤过：因子为
\[
 M(w\rho-\rho+\nu),
 \qquad \text{重数为}\;\dim L(\lambda)_\nu,
\]
其中 \(\nu\) 遍历 \(L(\lambda)\) 的权；见\cite[Section 3.6]{Humphreys2008BGGCategoryO}。现在核对投影保留哪些因子。由\cref{b5:thm:harish-chandra}，若这个因子的中心特征标是 \(\chi_\lambda\)，就存在 \(t\in W\) 使
\[
 t(w\rho+\nu)=\rho+\lambda.
\]
记 \(Q_+\) 为简单根的非负整数组合。因 \(\rho\) 支配，\(\rho-tw\rho\in Q_+\)；又因有限维模的权在 \(W\) 下不变且不高于最高权，\(\lambda-t\nu\in Q_+\)。上式使这两项之和为零，所以两项均为零。\(\rho\) 的稳定子为平凡群，故 \(t=w^{-1}\)，从而 \(\nu=w\lambda\)。这个极端权的重数为一。因此投影只留下 \(M(w\mathbin{\cdot}\lambda)\) 一个 Verma 因子。投影与有限直和相容，所得第 \(j\) 项正是 \(D_j\)，微分由原复形诱导并保持正合。

相对微分的良定义及滤过消解核验、外幂权引理和 Verma 扩张判据的证明见\cite[Sections 6.2--6.5]{Humphreys2008BGGCategoryO}。\cite[Section 8.4, Theorem 8.30]{Kirillov2008LieGroups}给出了同一定理的陈述，但未展开证明。
\end{proofsketch}

下面在秩一情形给出具体映射，并直接证明正合性。

\begin{proposition}\label{b5:prop:bgg-sl2-explicit}
设 \(\mathfrak g=\mathfrak{sl}_2(\mathbb C)\)，\(m\in\mathbb Z_{\ge0}\)，记 \(M(\mu)\) 为 \(\mathfrak g\) 的最高权 \(\mu\) 的 Verma 模，\(L_m\) 为最高权 \(m\)、维数 \(m+1\) 的不可约 \(\mathfrak g\) 模。令 \(v,u\) 分别为 \(M(m),M(-m-2)\) 的非零最高权生成元，\(f=E_{21}\)，其中 \(E_{21}\) 为矩阵单位。则下式给出一个不分裂的短正合列，\(\iota\) 按所有整数 \(j\ge0\) 时的公式定义：
\[
0\longrightarrow M(-m-2)\overset{\iota}{\longrightarrow}
M(m)\longrightarrow L_m\longrightarrow0,
\qquad \iota(f^ju)=f^{j+m+1}v,
\]
\end{proposition}
\begin{proof}
向量 \(f^{m+1}v\) 的权为 \(-m-2\)，且
\(ef^{m+1}v=(m+1)(m-m)f^mv=0\)。Verma 泛性质给出 \(\iota\)；PBW 基表明所有 \(f^{j+m+1}v\) 线性无关，所以 \(\iota\) 单射。商模以 \(v,fv,\ldots,f^mv\) 为基，\(e,f,h\) 的作用正是有限维简单模 \(L_m\) 的标准公式，故商为 \(L_m\)。

若正合列分裂，商中最高向量的像必须是 \(M(m)\) 中权 \(m\) 的向量，因此是 \(v\) 的非零倍数。这个向量生成整个 \(M(m)\)，却在 \(L_m\) 中应被 \(f^{m+1}\) 零化，矛盾。因此它不分裂。
\end{proof}

在一般消解中，逐权取维数仍合法，因为每个权空间有限维。欧拉交错和给出
\[
\operatorname{ch}L(\lambda)
=\sum_{w\in W}(-1)^{\ell(w)}\operatorname{ch}M(w\mathbin{\cdot}\lambda).
\]
再用 PBW 的
\(\operatorname{ch}M(\mu)=e^\mu\prod_{\alpha>0}(1-e^{-\alpha})^{-1}\)，即可恢复 Weyl 公式。这个推导从 BGG 消解解释前章的公式。

\subsection{Ext 与最高权范畴}
\label{b5:subsec:2004:02}

\begin{definition}\label{b5:def:yoneda-ext-o}
设 \(\mathfrak g\) 为有限维复半单 Lie 代数，固定 Cartan 子代数及正根集合，令 \(\mathcal O\) 为相应的范畴 O，\(M,N\in\mathcal O\)。中间项 \(E\) 也在 \(\mathcal O\) 内的短正合列
\(0\to N\to E\to M\to0\)，按在两端为恒等的交换图取等价类，并用拉回、推出定义 Baer 加法，所得群记作 \(\operatorname{Ext}^1_{\mathcal O}(M,N)\)。零类恰为分裂列。
\end{definition}
\symidx{\(\operatorname{Ext}^1_{\mathcal O}(M,N)\)：范畴 O 中短正合扩张的等价类群}

\cref{b5:prop:bgg-sl2-explicit}因此给出
\(\operatorname{Ext}^1_{\mathcal O}\bigl(L_m,M(-m-2)\bigr)\ne0\)。这里不能把 BGG 消解误称为投射消解：它的各项是 Verma 模，而 Verma 模并非都投射。

\begin{example}\label{b5:exm:o-costandard-sl2}
在 \(M(m)\) 的每个一维权空间上取对偶，令 \(\varphi_j\) 与 \(f^jv\) 对偶，取其代数直和 \(D M(m)=\bigoplus_{j\ge0}\mathbb C\varphi_j\)。具体地，令 \(\tau(e)=f\)、\(\tau(f)=e\)、\(\tau(h)=h\)，并按 \(\tau(uv)=\tau(v)\tau(u)\) 延拓为包络代数的反自同构，定义
\[
(x\varphi)(v)=\varphi(\tau(x)v).
\]
这个作用没有普通对偶中的负号：它是借 \(\tau\) 反转乘法，故 \(h\) 权保持为原来的权。直和中的泛函只在有限多个权空间上非零；整个代数对偶 \(M(m)^*\cong\prod_{j\ge0}\mathbb C\varphi_j\) 则允许在无限多个权空间上取非零值。因此这里取的是受限对偶，而非整个代数对偶。按定义逐基计算得到
\[
h\varphi_j=(m-2j)\varphi_j,\qquad
e\varphi_j=\varphi_{j-1},\qquad
f\varphi_j=(j+1)(m-j)\varphi_{j+1},
\]
其中 \(\varphi_{-1}=0\)。这些公式也可以直接验证三个 \(\mathfrak{sl}_2\) 括号。
若 \(m\ge0\)，前 \(m+1\) 个向量生成子模 \(L_m\)；商的最高权为 \(-m-2\)，后续 \(f\) 作用均非零，故经逐项缩放基后商同构于 \(M(-m-2)\)。中间模由 \(\varphi_{m+1}\) 生成，具有有限维权空间且 \(e\) 局部幂零，所以属于 \(\mathcal O\)。得到
\[
0\longrightarrow L_m\longrightarrow D M(m)
\longrightarrow M(-m-2)\longrightarrow0.
\]
该列不分裂：商的最高向量只有 \(\varphi_{m+1}\) 的标量倍提升能保持权，而
\(e\varphi_{m+1}=\varphi_m\ne0\)。于是 \(M(-m-2)\) 不是 \(\mathcal O\) 中的投射对象。
\end{example}

最高权范畴把简单模、标准模和投射模放在同一个偏序下。一个常用的有限版本如下：设有限长度的 \(\mathbb C\) 线性交换范畴具有有限维同态空间，有有限多个简单对象 \(L(\lambda)\)，由有限偏序集标记，并有投射覆盖 \(P(\lambda)\)。若可选择标准对象 \(\Delta(\lambda)\)，使其顶为 \(L(\lambda)\)、其余合成因子参数严格小于 \(\lambda\)，且
\(\operatorname{End}\Delta(\lambda)=\mathbb C\)，同时
\(P(\lambda)\twoheadrightarrow\Delta(\lambda)\)
的核有参数严格大于 \(\lambda\) 的标准对象滤过，就得到一种\term[zuigaoquan-fanchou]{最高权范畴}[highest weight category]结构。这个定义还要求投射覆盖与上述滤过确实存在；简单模有最高权本身并不足以推出这些条件。

\ProseParagraphBreak
可以直接看见投射性如何依赖范畴范围。对 \(\mathfrak{sl}_2\)，在权只允许为
\(m,m-2,m-4,\ldots\) 的 \(\mathcal O\) 全子范畴中，
\[
\operatorname{Hom}\bigl(M(m),N\bigr)\cong N_m.
\]
因为 \(N\) 中没有权 \(m+2\)，任何权 \(m\) 的向量都被 \(e\) 零化，Verma 泛性质便给出这个同构。取一个固定权空间是正合函子，故 \(M(m)\) 在这个受限范畴中投射。这个计算与上面的非投射例并不矛盾：投射性的判断始终包含允许哪些对象这一条件。

\ProseParagraphBreak
一般 \(\mathcal O\) 的分块与最高权结构还需要有限长度、投射覆盖和标准滤过理论，本节暂不展开。计算高阶 \(\operatorname{Ext}\) 时，应使用相应范畴中的投射或内射消解，或先证明替代消解的各项相对于目标函子为零调对象。BGG 消解的正合性本身尚不足以说明对它施加 \(\operatorname{Hom}\) 就能计算 \(\operatorname{Ext}\)。

\subsection{几何表示论所需的进一步工具}
\label{b5:subsec:2004:03}

几何方法从旗簇上的局部对象构造表示。在 \(\mathfrak{sl}_2\) 情形，旗簇就是 \(\mathbb P^1\)。次数为 \(m\ge0\) 的齐次多项式是线丛 \(\mathcal O(m)\) 的整体截面，而算子
\[
e=X\frac{\partial}{\partial Y},\qquad
f=Y\frac{\partial}{\partial X},\qquad
h=X\frac{\partial}{\partial X}-Y\frac{\partial}{\partial Y}
\]
使它成为 \(L_m\)。在坐标 \(z=Y/X\) 下，这些算子变成
\(e=\partial_z\)、\(h=m-2z\partial_z\)、\(f=mz-z^2\partial_z\)。
\appref{b5:app:C}将从坐标变换核验这些局部微分算子以及相应中心关系。

\ProseParagraphBreak
高维旗簇的 Schubert 分解\footnote{舒伯特（Hermann Cäsar Hannibal Schubert，1848-1911），德国数学家，开创了代数几何中的枚举几何研究。}按 Weyl 群编号，但胞腔闭包可能奇异。一般简单模的重数问题因而需要层的导出函子、交叉上同调和微分算子模等工具；光滑胞腔与其可能奇异的闭包也须区分。关于 BGG 与旗簇、几何表示的关系，可参见\cite[Section 8.4 and Further reading, Geometric representation theory]{Kirillov2008LieGroups}。这些几何方法进一步研究表示的构造与重数。


\begin{chaptersummary}
范畴 \(\mathcal O\) 的有限生成性和正根局部有限性，结合 PBW，控制了权集的上界及各权空间维数。最高权简单模仍由所有复最高权参数化，但中心特征标只能记录移位 Weyl 轨道；\(\mathfrak{sl}_2\) 中这一差别由 \(z(z+2)\) 的二次对称性直接表现。

\ProseParagraphBreak
一阶上同调衡量交叉映射是否可消去，二阶上同调衡量交换核扩张是否分裂。Casimir 链同伦与 Killing 型给出 Whitehead 消失；再沿可解根的交换理想归纳，得到一般 Levi 分解及规定范围内的共轭性。交换核的单步计算是归纳的一层，并不能单独代替一般定理。

\ProseParagraphBreak
BGG 消解以 Verma 模记录简单模的关系，它给出形式特征标的交错和，但不自动计算任意 \(\operatorname{Ext}\)。秩一的正向、对偶两个不分裂扩张说明，无限维最高权模未必完全可约，有限维半单理论的结论不能直接推广到这里。
\end{chaptersummary}

\BookExercises

\begin{optionalexercise}\label{b5:ex:20-o-finite-generation}
在复半单 Lie 代数上令 \(M=\bigoplus_{n\ge0}\mathbb C\) 为平凡模。判断范畴 \(\mathcal O\) 的各项条件，并说明哪个有限性结论失败。
\end{optionalexercise}
\begin{solution}
它只有权零，正根在每个向量上作用为零，故满足权分解和局部有限性。但包络代数对任何向量只产生其一维张成空间，有限多个生成元只能生成有限维子空间，所以 \(M\) 不有限生成。其零权空间就是整个无限维 \(M\)，说明权空间有限性不能只从后两个条件推出。
\end{solution}

\begin{optionalexercise}\label{b5:ex:20-o-tensor}
设 \(M\in\mathcal O\)，\(V\) 是有限维复半单 Lie 代数模。证明带对角作用的 \(M\otimes V\) 仍属于 \(\mathcal O\)。
\end{optionalexercise}
\begin{solution}
由有限维模的权分解，张量积按权和分解。对 \(m\otimes v\)，其正根轨道包含在 \(U(\mathfrak n^+)m\otimes V\) 中，后者有限维；有限和仍如此。若 \(m_1,\ldots,m_s\) 生成 \(M\)，取 \(V\) 的基 \(v_j\)。这些 \(m_i\otimes v_j\) 生成张量积：按 \(u\in U(\mathfrak g)\) 的次数归纳，恒等式
\(xm\otimes v=x(m\otimes v)-m\otimes xv\)
把第一因子上较长的单项式写成对角作用及较短单项式。因此三个条件均成立。
\end{solution}

\begin{optionalexercise}\label{b5:ex:20-same-character}
证明 \(\mathfrak{sl}_2\) 的 \(M(1/2)\) 与 \(M(-5/2)\) 是不同构的简单模，但具有相同无穷小特征标。
\end{optionalexercise}
\begin{solution}
两个参数都不是非负整数，因而由 \(ef^jv=j(z-j+1)f^{j-1}v\) 可见没有正次数的最高向量；任一非零子模取极大权就产生最高向量，所以两模简单。它们最高权不同，故不同构。两参数的 \(z(z+2)\) 同为 \(5/4\)，而中心由 \(\Omega\) 生成，所以完整中心特征标相同。
\end{solution}

\begin{optionalexercise}\label{b5:ex:20-casimir-projectors}
在 \(L_1\otimes L_1\) 上利用 \(\Omega\) 写出投到对称平方与外平方的投影。
\end{optionalexercise}
\begin{solution}
对称平方是 \(L_2\)，外平方是一维平凡模 \(L_0\)。\(\Omega\) 在两者上分别作用为 \(2(2+2)=8\) 和零。因此投影为 \(P_{\mathrm{sym}}=\Omega/8\)、\(P_{\mathrm{alt}}=I-\Omega/8\)。这里 \(\Omega\) 取张量积对角 Lie 代数作用，不能把它误写为两个因子上 \(\Omega\) 的简单和；其交叉项正是区分两种分量所必需的。
\end{solution}

\begin{optionalexercise}\label{b5:ex:20-sl3-cubic-center}
在 \(a+b+c=0\) 中找两点，其平方和相同，但立方和不同，并解释为什么只用二次 Casimir 不足以区分一般中心特征标。
\end{optionalexercise}
\begin{solution}
取 \((1,-1,0)\) 与 \((t,t,-2t)\)，\(t=1/\sqrt3\)。平方和都是二，立方和分别为零与 \(-6t^3=-2/\sqrt3\)。令两点分别为 \(\lambda+\rho\)、\(\mu+\rho\)。Harish--Chandra 同构给出与三次不变多项式对应的中心元，能够区分它们；二次 Casimir 对应平方和减去同一常数，不能区分。
\end{solution}

\begin{optionalexercise}\label{b5:ex:20-dot-orbit}
对 \(\mathfrak{sl}_3\) 取 \(\rho=(1,0,-1)\)。列出零权的移位 Weyl 轨道，并指出其中哪个权支配。
\end{optionalexercise}
\begin{solution}
排列 \(\rho\) 再减去 \(\rho\)，得到
\[
(0,0,0),\ (0,-1,1),\ (-1,1,0),\ (-1,-1,2),\ (-2,1,1),\ (-2,0,2).
\]
支配条件是第一坐标不小于第二、第二不小于第三，故只有零权支配。六个最高权具有同一中心特征标，但只有零权对应有限维简单模。
\end{solution}

\begin{optionalexercise}\label{b5:ex:20-one-dimensional-cohomology}
设 \(\mathfrak a=\mathbb Cx\)，\(x\) 在模 \(M\) 上按算子 \(T\) 作用。计算所有 \(H^q(\mathfrak a,M)\)。
\end{optionalexercise}
\begin{solution}
外幂在次数二起为零。复形是 \(0\to M\overset{T}{\to}M\to0\)，所以 \(H^0=\ker T\)、\(H^1=M/TM\)、\(H^q=0\) 对 \(q\ge2\)。即使 \(M\) 有限维，非半单的一维交换 Lie 代数也可能有非零一阶上同调，例如 \(T=0\)。
\end{solution}

\begin{optionalexercise}\label{b5:ex:20-affine-cohomology}
设二维 Lie 代数有基 \(x,y\)，\([x,y]=y\)。计算平凡系数上同调。
\end{optionalexercise}
\begin{solution}
令 \(\alpha,\beta\) 是对偶基。\(\mathrm{d}\alpha=0\)、\(\mathrm{d}\beta=-\alpha\wedge\beta\)，且 \(\mathrm{d}:C^0\to C^1\) 为零。因此 \(H^0=\mathbb C\)、\(H^1=\mathbb C[\alpha]\)、\(H^2=0\)，更高次数为零。一阶非零检测到交换化 \(\mathfrak a/[\mathfrak a,\mathfrak a]\) 的一维性。
\end{solution}

\begin{optionalexercise}\label{b5:ex:20-heisenberg-cohomology}
对 \([x,y]=z\)、其余基括号为零的三维 Heisenberg Lie 代数，求平凡系数上同调的维数与代表元。
\end{optionalexercise}
\begin{solution}
对偶基记作 \(\alpha,\beta,\gamma\)。有 \(\mathrm{d}\alpha=\mathrm{d}\beta=0\)、\(\mathrm{d}\gamma=-\alpha\wedge\beta\)。微分为外代数导子，所以全部二形式的微分为零。于是四个次数的维数依次为 \(1,2,2,1\)，可分别取代表
\(1\)；\(\alpha,\beta\)；\(\alpha\wedge\gamma,\beta\wedge\gamma\)；\(\alpha\wedge\beta\wedge\gamma\)。
\end{solution}

\begin{optionalexercise}\label{b5:ex:20-abelian-extension-classes}
设 \(A=\mathbb C^n\) 是交换 Lie 代数，\(M=\mathbb C\) 为平凡模。说明中心扩张在固定两端的意义下由交错型参数化，并解释为何允许任意基变换时参数空间会改变。
\end{optionalexercise}
\begin{solution}
全部微分为零，所以 \(H^2(A,\mathbb C)=\Lambda^2 A^*\)。扩张括号由 \(\bigl[(a,s),(b,t)\bigr]=\bigl(0,\omega(a,b)\bigr)\) 给出。固定两端的等价不能改变 \(\omega\)。若允许在 \(A\) 上作 \(g\in\operatorname{GL}(A)\) 并在中心上缩放 \(c\ne0\)，则型可变为 \(c\,\omega(g^{-1}a,g^{-1}b)\)。因此那时分类的是这种作用的轨道，而不是 \(H^2\) 中逐个向量。
\end{solution}

\begin{optionalexercise}\label{b5:ex:20-sl2-central-cocycle}
设 \(c\) 是 \(\mathfrak{sl}_2\) 上的平凡系数\(2\)-余圈。直接以 \(c(e,f),c(h,e),c(h,f)\) 写出 \(b\)，使 \(\mathrm{d}b=c\)。
\end{optionalexercise}
\begin{solution}
令
\(b(h)=-c(e,f)\)、\(b(e)=-c(h,e)/2\)、\(b(f)=c(h,f)/2\)。平凡系数满足 \(\mathrm{d}b(x,y)=-b\bigl([x,y]\bigr)\)，在三对基向量上分别给出指定的三个值，所以 \(\mathrm{d}b=c\)。实际上任意交错二形式都满足余圈条件，因为在 \((h,e,f)\) 上有 \(-c(2e,f)+c(-2f,e)-c(h,h)=0\)。
\end{solution}

\begin{optionalexercise}\label{b5:ex:20-sl2-third-cohomology}
计算 \(H^3(\mathfrak{sl}_2,\mathbb C)\)，并说明它与 Whitehead 引理不矛盾。
\end{optionalexercise}
\begin{solution}
\(C^3\) 一维且 \(C^4=0\)。上一题的直接计算表明 \(\mathrm{d}:C^2\to C^3\) 为零，故 \(H^3\cong\mathbb C\)。可取在 \((h,e,f)\) 上为一的体积形式。Whitehead 两引理只断言一、二阶消失；Casimir 链同伦也不能在平凡系数上求逆，因为其 Casimir 为零。
\end{solution}

\begin{optionalexercise}\label{b5:ex:20-infinite-whitehead-counterexample}
在 \(\mathfrak{sl}_2\) 模 \(M(-2)\) 中令 \(u\) 为最高向量。证明 \(b(e)=b(h)=0\)、\(b(f)=u\) 给出非零的一阶上同调类。
\end{optionalexercise}
\begin{solution}
对 \([h,e]=2e\) 余圈式两边均为零；对 \([h,f]=-2f\)，两边均为 \(-2u\)；对 \([e,f]=h\)，右边为 \(eu=0\)。故 \(b\) 是余圈。若 \(b(x)=xv\)，则 \(hv=0\)。但 \(M(-2)\) 的权为 \(-2,-4,\ldots\)，所以 \(v=0\)，与 \(b(f)=u\ne0\) 矛盾。这说明有限维系数假设不能删去。
\end{solution}

\begin{optionalexercise}\label{b5:ex:20-abelian-radical-complements}
在 \(\mathfrak g=\mathbb C^2\rtimes\mathfrak{sl}_2\) 中，\(\mathbb C^2\) 是自然模且交换。描述所有 Levi 子代数，并给出将标准补送到每一个补的指数。
\end{optionalexercise}
\begin{solution}
商为 \(\mathfrak{sl}_2\)，每个补是 \(x\mapsto x+b(x)\) 的图。图是子代数恰当 \(b\) 为一余圈，所以 \(b(x)=xv\)。自然模没有不变向量，故 \(v\) 唯一。所求补为 \(\{x+xv\}\)，且 \(\exp(-\operatorname{ad}v)x=x+xv\)，因为交换根使 \((\operatorname{ad}v)^2=0\)。
\end{solution}

\begin{optionalexercise}\label{b5:ex:20-heisenberg-levi-conjugacy}
令 \(\mathfrak r=\langle x,y,z\rangle\) 满足 \([x,y]=z\)，\(z\) 中心。令 \(\mathfrak{sl}_2\) 在 \(x,y\) 上按自然模作用，在 \(z\) 上平凡作用。对 \(a,b\in\mathbb C\)，求 \(\exp\bigl(\operatorname{ad}(ax+by)\bigr)\) 作用于标准 Levi 基 \(e,f,h\) 的结果。
\end{optionalexercise}
\begin{solution}
先核验该作用保持 \([x,y]=z\)：自然模矩阵迹为零，所以在 \(\Lambda^2\langle x,y\rangle\) 上作用为零。令 \(v=ax+by\)，则
\[
e\longmapsto e-bx+\tfrac12b^2z,\qquad
f\longmapsto f-ay-\tfrac12a^2z,\qquad
h\longmapsto h-ax+by+abz.
\]
例如 \([v,e]=-bx\)、\(\bigl[v,[v,e]\bigr]=b^2z\)，更高括号为零；其他两式同样由前两次括号得到。二次中心修正说明非交换可解根中的共轭通常不能只写成一次线性修正。
\end{solution}

\begin{optionalexercise}\label{b5:ex:20-bgg-character}
从秩一 BGG 列计算 \(L_m\) 的形式特征标，不直接列举其基。
\end{optionalexercise}
\begin{solution}
\(M(z)\) 的特征标为 \(e^z/(1-e^{-2})\)。正合性逐权给出
\[
\operatorname{ch}L_m
=\frac{e^m-e^{-m-2}}{1-e^{-2}}
=e^m+e^{m-2}+\cdots+e^{-m}.
\]
最后一步是有限几何级数；这同时给出维数 \(m+1\) 及所有权重数为一。
\end{solution}

\begin{optionalexercise}\label{b5:ex:20-restricted-dual}
对 \(M(0)\) 的受限对偶写出前四个基向量上的 \(e,f,h\) 作用，并指出它与 \(M(0)\) 不同构的一个结构理由。
\end{optionalexercise}
\begin{solution}
有 \(h\varphi_j=-2j\varphi_j\)、\(e\varphi_j=\varphi_{j-1}\)、\(f\varphi_j=-j(j+1)\varphi_{j+1}\)。因此 \(\varphi_0\) 张成平凡子模，\(\varphi_1\) 生成全模，且 \(e\varphi_1=\varphi_0\)。\(M(0)\) 没有平凡子模：其权零空间由最高生成元 \(v\) 张成，但 \(fv\ne0\)。故两模不同构。
\end{solution}

\begin{optionalexercise}\label{b5:ex:20-verma-projective-truncation}
设 \(z\in\mathbb C\)。在权集包含于 \(z-2\mathbb Z_{\ge0}\) 的 \(\mathfrak{sl}_2\) 范畴 \(\mathcal O\) 子范畴中，证明任一满射 \(E\to N\) 都使 \(\operatorname{Hom}\bigl(M(z),E\bigr)\to\operatorname{Hom}\bigl(M(z),N\bigr)\) 满射。
\end{optionalexercise}
\begin{solution}
同态由权 \(z\) 的最高向量决定。在 \(N_z\) 中给定这样的向量，可以沿满射在 \(E_z\) 中提升，因为模同态保持权且每个向量只有有限个权分量。\(E\) 没有权 \(z+2\)，所以提升自动被 \(e\) 零化。由 Verma 泛性质得到所需提升同态。这直接验证了相对于该子范畴的投射性。
\end{solution}

\begin{optionalexercise}\label{b5:ex:20-differential-module-condition}
在二维交换 Lie 代数 \(\mathbb Cx\oplus\mathbb Cy\) 上，若企图用不交换的算子 \(A,B\) 分别表示 \(x,y\)，计算形式上的 \(\mathrm{d}^2m\)，说明复形为何失败。
\end{optionalexercise}
\begin{solution}
形式上 \((\mathrm{d}m)(x)=Am\)、\((\mathrm{d}m)(y)=Bm\)。因括号为零，\((\mathrm{d}^2m)(x,y)=ABm-BAm\)。若 \([A,B]\ne0\)，就存在 \(m\) 使此式非零。所以 \(\mathrm{d}^2=0\) 的验证实质用到了表示保持 Lie 括号，不能对任意一组线性算子写下公式就称上同调。
\end{solution}

\begin{optionalexercise}\label{b5:ex:20-bgg-not-projective}
解释为何秩一 BGG 消解不能仅凭形式直接用于计算任意 \(\operatorname{Ext}^1_{\mathcal O}(L_m,N)\)。给出其中一项不投射的明确证据。
\end{optionalexercise}
\begin{solution}
通常的 \(\operatorname{Hom}\) 复形计算导出函子，需要消解项投射，或先证明消解的各项相对于所用 Hom 函子为零调对象。BGG 的左项是 \(M(-m-2)\)，而正文构造的不分裂列
\(0\to L_m\to D M(m)\to M(-m-2)\to0\)
说明它不投射：否则恒等映射能沿右端满射提升，短正合列就会分裂。因此缺少必要条件，不能作所述直接推断。
\end{solution}

\begin{optionalexercise}\label{b5:ex:20-nontrivial-coefficient-all-degrees}
设 \(\mathfrak s\) 半单，\(M\) 为有限维复模且 \(M^{\mathfrak s}=0\)。证明 \(H^q(\mathfrak s,M)=0\) 对所有 \(q\ge0\) 成立。
\end{optionalexercise}
\begin{solution}
完全可约性把 \(M\) 分解为非平凡简单模之和；无不变向量排除了平凡分量。每个分量上的 Killing Casimir 标量非零，所以它在整个 \(M\) 上可逆。对任意余圈 \(\varphi\)，链同伦公式给出 \(\varphi=\mathrm{d}(C_M^{-1}H\varphi)\)。次数零的同一公式给出不变量为零。注意本题加强了系数假设，因此不与平凡系数三阶非零的例子冲突。
\end{solution}

\begin{optionalexercise}\label{b5:ex:20-differential-operator-casimir}
在 \(\mathbb C[z]\) 上令 \(e=\partial_z\)、\(h=m-2z\partial_z\)、\(f=mz-z^2\partial_z\)，其中 \(m\in\mathbb C\)。计算 \(\Omega=h^2+2h+4fe\) 的作用。
\end{optionalexercise}
\begin{solution}
令 \(D=z\partial_z\)。则 \(D^2=z^2\partial_z^2+D\)，所以
\(h^2+2h=m(m+2)-4mD+4z^2\partial_z^2\)，而
\(4fe=4mD-4z^2\partial_z^2\)。相加得到 \(m(m+2)I\)。对于整数 \(m\ge0\)，次数不超过 \(m\) 的子空间被 \(f\) 保持，正是有限维最高权模；但算子恒等式本身对任意复 \(m\) 都成立。
\end{solution}

% !TeX root = ../../Book5.tex
% 纲要标题：## 第21章*　Hopf 代数与仿射代数群
% 纲要位置：工作记录/计划纲要.md，第 4389 行。

\OptionalChapter{Hopf 代数与仿射代数群}{b5:ch:21}
\BookChapterNumber{b5:ch:21}

\begin{chapterabstract}
本章建立 Hopf 代数的乘法、余乘法、余单位和对极，比较群代数、包络代数与坐标环，并用余模和矩阵系数研究有理表示与群概形例子。
\end{chapterabstract}

\begin{learninggoals}
\begin{itemize}
\item 在具体生成元上核验余乘法、余单位和对极，并构造张量与对偶作用。
\item 区分群代数和坐标代数、模与余模、基域点与所有交换代数上的点。
\item 从余模证明有理表示的局部有限性、环面权分解及仿射群的线性实现。
\item 说明仅由微分表示无法恢复哪些整数与正特征信息，比较等距、相似和对合自同构群。
\end{itemize}
\end{learninggoals}

乘法群的元素可以相乘，也可以取逆。若把它看成由坐标函数 \(t\) 与 \(t^{-1}\) 描述的对象，两个群元素相乘时，\(t\) 的取值随之相乘；在坐标环中，这表现为 \(\Delta(t)=t\otimes t\)。单位元给出 \(\varepsilon(t)=1\)，取逆则给出 \(S(t)=t^{-1}\)。原来作用在点上的三种运算，如今变成作用在函数上的代数映射，并满足相应的恒等式。沿着这个具体例子，Hopf 代数的公理和表示张量积所需的余乘法就有了可核对的来源。

\ProseParagraphBreak
本章使用张量积、代数同态和第三卷的仿射坐标环语言；群概形从可表示函子及坐标方程出发介绍。中心商的局部计算还使用第三卷第24.3节的射影仿射图和第12.5节的忠实平坦性。包络代数例子调用第15章，四元数例子则把所需乘法与范数公式重新列出。

\ProseParagraphBreak
Hopf 代数的结构及对偶可参阅\cite[Sections 5.2--5.3]{EtingofEtAl2015TensorCategories}。该书的 Hopf 定义要求对极可逆，本章采用更一般的定义并在使用对偶时说明实际所需条件；所讨论的群代数、包络代数和交换坐标代数不受这一区别影响。

% !TeX root = ../../Book5.tex
% 纲要标题：### 21.1 余代数与 Hopf 代数
% 纲要位置：工作记录/计划纲要.md，第 4391 行。

\section{余代数与 Hopf 代数}
\label{b5:sec:2101}


% 纲要范围：
%
% * 余乘法、余单位与对极
% * 群代数及坐标环的对偶方向
% * 模与余模的区别
%

群表示与 Lie 代数表示都能作张量积，但作用公式不同：群元素同时作用于两个因子，Lie 代数元素则分别作用后相加。余乘法把这两种规则统一写成代数中的映射；再把单位表示和对偶表示所需的规则一并加入，就得到 Hopf 代数。

\subsection{余乘法、余单位与对极}
\label{b5:subsec:2101:01}

\begin{definition}\label{b5:def:coalgebra-hopf}
设 \(k\) 为域，\(H\) 为 \(k\) 向量空间。线性映射
\(\Delta:H\to H\otimes_kH\)、\(\varepsilon:H\to k\) 满足
\[
(\Delta\otimes1)\Delta=(1\otimes\Delta)\Delta,\qquad
(\varepsilon\otimes1)\Delta=1_H=(1\otimes\varepsilon)\Delta
\]
时，称 \((H,\Delta,\varepsilon)\) 为\term[yudaishu]{余代数}[coalgebra]，称两映射为余乘法与余单位；后两式用 \(k\otimes H=H=H\otimes k\) 识别。
若 \(H\) 还为含幺结合 \(k\) 代数，且 \(\Delta,\varepsilon\) 是含幺代数同态，就称为\term[shuangdaishu]{双代数}[bialgebra]。若另有线性映射 \(S:H\to H\) 满足
\[
m(S\otimes1)\Delta=\eta\varepsilon=m(1\otimes S)\Delta,
\]
其中 \(m\) 是乘法、\(\eta:k\to H\) 是单位映射，就称为\term[Hopf-daishu]{Hopf 代数}[Hopf algebra]，称 \(S\) 为\term[duiji]{对极}[antipode]。
\end{definition}
\symidx{\(\Delta,\varepsilon,S\)：Hopf 代数的余乘法、余单位与对极}
本书不把对极的可逆性额外放进定义；它在本章的群代数、包络代数及交换坐标代数例子中成立。Etingof 等的《Tensor Categories》在定义 Hopf 代数时要求对极可逆，因此引用其张量范畴结论时须保留这一差别。
\ProseParagraphBreak
常把 \(\Delta(h)\) 写为 \(\sum h_{(1)}\otimes h_{(2)}\)，称为 Sweedler 记号。它表示一个有限和，不表示每个 \(h\) 都只有一个纯张量分解。余结合律允许把三重余乘法统一写为 \(\sum h_{(1)}\otimes h_{(2)}\otimes h_{(3)}\)。
\symidx{\(\sum h_{(1)}\otimes h_{(2)}\)：余乘法的 Sweedler 记号}
\begin{proposition}\label{b5:prop:antipode-unique}
设 \(k\) 为域，\(H\) 为 \(k\) 双代数。其对极 \(S:H\to H\) 若存在，则唯一，并满足 \(S(1)=1\)、\(S(ab)=S(b)S(a)\) 对全部 \(a,b\in H\) 成立。
\end{proposition}
\begin{proof}
在 \(\operatorname{Hom}_k(H,H)\) 上定义卷积
\((f*g)(h)=\sum f\bigl(h_{(1)}\bigr)g\bigl(h_{(2)}\bigr)\)。
结合律来自 \(m\) 的结合与 \(\Delta\) 的余结合，单位为 \(\eta\varepsilon\)。对极就是恒等映射的双边卷积逆，故唯一；代入单位即得 \(S(1)=1\)。
对反乘法性质，改在 \(\operatorname{Hom}_k(H\otimes H,H)\) 中卷积，张量积余乘法把 \(a\otimes b\) 送到
\(\sum\bigl(a_{(1)}\otimes b_{(1)}\bigr)\otimes\bigl(a_{(2)}\otimes b_{(2)}\bigr)\)。
映射 \(S\circ m\) 是乘法映射 \(m\) 的卷积逆，因 \(\Delta\) 保持乘法。另一个映射 \(F(a,b)=S(b)S(a)\) 也给出逆：两侧复合分别为
\[
\sum S\bigl(b_{(1)}\bigr)S\bigl(a_{(1)}\bigr)a_{(2)}b_{(2)}
=\varepsilon(a)\varepsilon(b)1
\]
和
\(\sum a_{(1)}b_{(1)}S\bigl(b_{(2)}\bigr)S\bigl(a_{(2)}\bigr)=\varepsilon(a)\varepsilon(b)1\)。
逆的唯一性使两映射相同。
\end{proof}
若再有交换 \(\Delta\) 两因子不改变它，则称余交换。乘法交换与余乘法余交换是两项不同条件，以下两种来自同一个群的代数将说明这种区别。

\subsection{群代数及坐标环的对偶方向}
\label{b5:subsec:2101:02}

\begin{proposition}\label{b5:prop:group-enveloping-hopf}
设 \(k\) 为域，\(G\) 为群，\(\mathfrak g\) 为 \(k\) 上 Lie 代数。下列公式分别给出 \(kG\) 与 \(U(\mathfrak g)\) 上的 Hopf 代数结构：
\[
\begin{array}{c|ccc}
&H:\Delta&\varepsilon&S\\ \hline
kG&g\mapsto g\otimes g&g\mapsto1&g\mapsto g^{-1}\\
U(\mathfrak g)&x\mapsto x\otimes1+1\otimes x&x\mapsto0&x\mapsto-x
\end{array}
\]
第二行的公式在 \(x\in\mathfrak g\) 上规定，再按相应代数或反代数映射延拓。两者都余交换。
\end{proposition}
\begin{proof}
群代数的各式在群基元上核验公理即可，例如
\(m(S\otimes1)\Delta(g)=g^{-1}g=1=\varepsilon(g)1\)。
包络代数中，映射 \(x\mapsto x\otimes1+1\otimes x\) 保持 Lie 括号，因为交叉的两个因子交换，故由泛性质延拓为代数同态。增广零化全部 Lie 括号；映射 \(x\mapsto-x\) 到反代数保持 Lie 括号，因而给出反同态 \(S\)。余结合及余单位等式在生成元上成立，所以在全部代数上成立。对极等式可对字长归纳：若两较短字 \(u,v\) 满足左式，则
\[
\sum S\bigl((uv)_{(1)}\bigr)(uv)_{(2)}
=\sum S\bigl(v_{(1)}\bigr)S\bigl(u_{(1)}\bigr)u_{(2)}v_{(2)}
=\varepsilon(u)\varepsilon(v)1.
\]
右式同理；线性性完成证明。
\end{proof}
若 \(G\) 有限，函数空间 \(k^G\) 则按逐点乘法成为交换 Hopf 代数：
\[
\Delta f(g,h)=f(gh),\qquad \varepsilon(f)=f(1),\qquad
Sf(g)=f(g^{-1}).
\]
在点函数基 \(\delta_x\) 中，
\(\Delta\delta_x=\sum_{ab=x}\delta_a\otimes\delta_b\)。
这是 \(kG\) 的线性对偶 Hopf 代数；函数代数通常不余交换，恰在 \(G\) 交换时余交换。群无限时，任意函数的两变量拉回未必在代数张量积 \(k^G\otimes k^G\) 中，因此不能不加限制照搬有限对偶。

\ProseParagraphBreak
例如在加法群 \(\mathbb Z\) 上取点函数 \(\delta_0\)。乘法拉回在 \((m,n)\) 处为 \(\delta_{m+n,0}\)。若它来自代数张量积，就可写成某个有限和
\[
\delta_{m+n,0}=\sum_{i=1}^r a_i(m)b_i(n).
\]
固定 \(m\) 后，这个关于 \(n\) 的函数必落在 \(b_1,\ldots,b_r\) 的线性生成空间中。但左边随 \(m\) 变化给出所有点函数 \(\delta_{-m}\)，它们线性无关：在各自支集点求值即可逐个读出系数。这与有限维生成空间矛盾，故该两变量函数确实不在 \(k^{\mathbb Z}\otimes k^{\mathbb Z}\) 中。

\subsection{模与余模的区别}
\label{b5:subsec:2101:03}

有限群表示在基下写成 \(gv_j=\sum_i a_{ij}(g)v_i\)。与其对每个 \(g\) 分别保存这个矩阵，也可以把系数看成函数 \(a_{ij}\in k^G\)，一次写成 \(v_j\mapsto\sum_i v_i\otimes a_{ij}\)。群乘法在系数函数上的拉回，正好给出这种展开应满足的相容条件。这便是余模的来源。

\begin{definition}\label{b5:def:right-comodule}
设 \(k\) 为域，\((C,\Delta,\varepsilon)\) 为 \(k\) 余代数。\(k\) 向量空间 \(V\) 配 \(k\) 线性映射 \(\delta:V\to V\otimes_k C\)，使
\[
(\delta\otimes1)\delta=(1\otimes\Delta)\delta,\qquad
(1\otimes\varepsilon)\delta=1_V,
\]
时，称为右 \(C\)\term[yumo]{余模}[comodule]。第一式将两种继续展开余作用的次序联系起来，即下图交换：
\[
\begin{tikzcd}[column sep=large,row sep=large]
V\arrow[r,"\delta"]\arrow[d,"\delta"']
&V\otimes_kC\arrow[d,"1_V\otimes\Delta"]\\
V\otimes_kC\arrow[r,"\delta\otimes1_C"']
&V\otimes_kC\otimes_kC
\end{tikzcd}
\]
给定两个右 \(C\) 余模 \((V,\delta_V)\)、\((W,\delta_W)\)，若 \(k\) 线性映射 \(f:V\to W\) 使下图交换，就称它为余模同态：
\[
\begin{tikzcd}[column sep=large,row sep=large]
V\arrow[r,"\delta_V"]\arrow[d,"f"']
&V\otimes_kC\arrow[d,"f\otimes1_C"]\\
W\arrow[r,"\delta_W"']&W\otimes_kC
\end{tikzcd}
\]
\end{definition}
\symidx{\(\delta:V\to V\otimes_kC\)：右余模的余作用}
模是代数元素作用于向量的规则；余模是把一个向量展开成向量与系数函数的规则。只有适当的对偶条件成立时，两种语言才可直接转换。例如有限维代数 \(A\) 的左模等同于右 \(A^*\) 余模；取对偶基 \(a_i,a^i\)，余作用为
\(v\mapsto\sum_i a_iv\otimes a^i\)。
代数乘法公理与单位公理分别对应余结合与余单位。

\ProseParagraphBreak
一个二维例子已能看清余作用的含义。设 \(k\) 为域，\(C=k[t]\) 的余乘法为 \(\Delta(t)=t\otimes1+1\otimes t\)，余单位为 \(\varepsilon(t)=0\)。在 \(V=kv_1\oplus kv_2\) 上令
\[
\delta(v_1)=v_1\otimes1,\qquad
\delta(v_2)=v_2\otimes1+v_1\otimes t.
\]
在 \(v_2\) 上，余结合律两边都等于
\(v_2\otimes1\otimes1+v_1\otimes t\otimes1+v_1\otimes1\otimes t\)，余单位律也直接成立。把 \(t\) 代成 \(c\in k\)，便得到矩阵
\(\bigl(\begin{smallmatrix}1&c\\0&1\end{smallmatrix}\bigr)\)；余乘法中的两个加项对应参数相加。直线 \(kv_1\) 是子余模，商余模平凡，但将 \(v_2\) 换成 \(v_2+av_1\) 仍留下 \(v_1\otimes t\) 项，故这条子余模没有补。

\begin{proposition}\label{b5:prop:hopf-tensor-dual}
设 \(k\) 为域，\((H,\Delta,\varepsilon,S)\) 为 \(k\) 上的 Hopf 代数，\(V,W\) 为幺性左 \(H\) 模。对 \(h\in H\)，用 Sweedler 记号写 \(\Delta(h)=\sum h_{(1)}\otimes h_{(2)}\)。则张量积 \(V\otimes_kW\) 通过
\[
h(v\otimes w)=\sum h_{(1)}v\otimes h_{(2)}w
\]
成为左模，\(k\) 通过 \(\varepsilon\) 成为单位模。若 \(V\) 有限维，则
\[
(h\varphi)(v)=\varphi\bigl(S(h)v\bigr)
\]
给出 \(V^*\) 的左模结构。对 \(V\) 的基 \(e_1,\ldots,e_r\) 及对偶基 \(e_1^*,\ldots,e_r^*\)，求值与余求值映射
\[
\operatorname{ev}:V^*\otimes_kV\to k,\qquad
\varphi\otimes v\mapsto\varphi(v),
\]
\[
\operatorname{coev}:k\to V\otimes_kV^*,\qquad
1\mapsto\sum_i e_i\otimes e_i^*
\]
均为 \(H\) 模同态，并使 \(V^*\) 成为 \(V\) 的左对偶。
\end{proposition}
\begin{proof}
\(\Delta\) 保持乘法使两次作用等于乘积作用，余结合使三重张量的两种括号相容，余单位使 \(k\otimes V\to V\) 相容。对偶上依次作用 \(a,b\) 得到
\(\varphi\bigl(S(b)S(a)v\bigr)=\varphi\bigl(S(ab)v\bigr)\)，所以是左作用。求值映射的等变性为
\[
\sum\bigl(h_{(1)}\varphi\bigr)\bigl(h_{(2)}v\bigr)
=\varphi\left(\sum S\bigl(h_{(1)}\bigr)h_{(2)}v\right)
=\varepsilon(h)\varphi(v).
\]
为核验余求值，使用典范同构
\[
\Theta:V\otimes_kV^*\xrightarrow{\sim}\operatorname{End}_k(V),\qquad
\Theta(v\otimes\varphi)(x)=\varphi(x)v.
\]
张量 \(\operatorname{coev}(1)\) 对应 \(1_V\)，因而不依赖基的选择。对任意 \(x\in V\)，有
\[
\begin{aligned}
\Theta\bigl(h\operatorname{coev}(1)\bigr)(x)
&=\sum_{i,(h)} e_i^*\bigl(S(h_{(2)})x\bigr)h_{(1)}e_i\\
&=\sum_{(h)}h_{(1)}S(h_{(2)})x
=\varepsilon(h)x.
\end{aligned}
\]
由 \(\Theta\) 单射可知
\(h\operatorname{coev}(1)=\varepsilon(h)\operatorname{coev}(1)\)，故余求值也是 \(H\) 模同态。略去通常的结合与单位同构，两条蛇形复合分别将
\[
v\longmapsto\sum_i e_i\,e_i^*(v)=v,\qquad
\varphi\longmapsto\sum_i\varphi(e_i)e_i^*=\varphi.
\]
这就给出左对偶的两条恒等式；以上计算只用了对极公理。
\end{proof}
在群代数上，张量公式就是对角群作用；在包络代数上，它是 Lie 代数的 Leibniz 型作用。普通交换 \(v\otimes w\mapsto w\otimes v\) 是否等变，要看余乘法是否余交换；\subsecref{b5:subsec:B01:03}给出一个非余交换的例子。上述基础可对照\cite[Definitions 5.2.2, 5.3.2; Propositions 5.3.5--5.3.6; Corollary 5.3.7]{EtingofEtAl2015TensorCategories}。

% !TeX root = ../../Book5.tex
% 纲要标题：### 21.2 代数群的坐标环
% 纲要位置：工作记录/计划纲要.md，第 4397 行。

\section{代数群的坐标环}
\label{b5:sec:2102}


% 纲要范围：
%
% * 仿射代数群
% * 群乘法与 Hopf 结构
% * 一般线性群及可对角化群
% * 特殊线性群到射影一般线性群的中心商与表示下降
%

代数群的乘法也可由函数上的拉回来描述。采用坐标环而非只看一个基域上的点，有两个好处：公式可以同时在不同交换代数上取值，正特征中的幂零信息也不会因为点集太小而消失。

\subsection{仿射代数群}
\label{b5:subsec:2102:01}

\begin{definition}\label{b5:def:affine-group}
设 \(k\) 为域，\(A\) 为有限生成的交换 \(k\) Hopf 代数，结构映射记为 \(\Delta,\varepsilon,S\)。对每个交换含幺 \(k\) 代数 \(R\)，记 \(\eta_R:k\to R\) 为单位映射，并令
\[
G(R)=\operatorname{Hom}_{k\text{-alg}}(A,R),
\]
对 \(a\in A\)，用 Sweedler 记号写 \(\Delta(a)=\sum a_{(1)}\otimes a_{(2)}\)。对 \(g,h\in G(R)\)，用余乘法、余单位及对极分别定义
\[
(gh)(a)=\sum g\bigl(a_{(1)}\bigr)h\bigl(a_{(2)}\bigr),\qquad
1_G=\eta_R\varepsilon,\qquad g^{-1}=g\circ S.
\]
这个群值函子称为一个有限型仿射群概形，记为 \(G=\operatorname{Spec}A\)，其坐标环记为 \(k[G]=A\)。在代数闭域上，若坐标环约化，也把相应的仿射簇及其群运算称为\term[fangshe-daishuqun]{仿射代数群}[affine algebraic group]。
\end{definition}
因为 \(A,R\) 都交换，两个代数同态的卷积仍为代数同态；Hopf 公理正好使以上运算结合、有单位与逆。一个 \(k\) 代数同态 \(R\to R'\) 与这些公式交换，所以这是对 \(R\) 自然的群，而非彼此不相干的一组点集。
\ProseParagraphBreak
两个这类群之间的态射，定义为反向的 Hopf 代数同态。特别地，点上的映射必须在所有 \(R\) 上由同一坐标公式给出；只给 \(G(k)\to H(k)\) 的一个抽象群同态，并不足以定义代数群态射。

\subsection{群乘法与 Hopf 结构}
\label{b5:subsec:2102:02}

坐标环方法与通常的正则映射方法相符。设 \(k\) 代数闭，\(G\) 为仿射簇，乘法、单位及求逆都是正则映射。则乘法拉回给出
\[
\Delta:k[G]\to k[G\times G]=k[G]\otimes k[G],
\]
单位与求逆分别给出 \(\varepsilon,S\)。正则函数在三元组 \((g,h,\ell)\) 上比较时，群乘法结合律就是余结合律，其他公理同理。反向从约化的有限型交换 Hopf 代数出发，正则函数关系恢复同样的运算。
\begin{proposition}\label{b5:prop:hopf-ideals-subgroups}
设 \(k\) 为域，\(A=k[G]\) 为有限生成的交换 \(k\) Hopf 代数，其结构映射为 \(\Delta,\varepsilon,S\)。若理想 \(I\subseteq A\) 满足
\[
\Delta(I)\subseteq I\otimes A+A\otimes I,\qquad
\varepsilon(I)=0,\qquad S(I)\subseteq I,
\]
则 \(A/I\) 自然成为交换 Hopf 代数，且
\(\operatorname{Spec}(A/I)\) 是 \(G\) 的闭子群概形。反过来，使商上的三种映射良定义所需的条件正是以上三式。
\end{proposition}
\begin{proof}
张量商映射 \(A\otimes A\to(A/I)\otimes(A/I)\) 的核为
\(I\otimes A+A\otimes I\)，故第一条件使余乘法下降；其余两条件分别使余单位与对极下降。所有公理在商中继承。对每个 \(R\)，商到 \(R\) 的代数同态就是那些零化 \(I\) 的 \(G(R)\) 元素，以上三式保证它们对乘法、单位和逆封闭。
\end{proof}
例如 \(k[t]\) 上的 \(\Delta(t)=t\otimes1+1\otimes t\) 描述加法群 \(\mathbb G_a\)。特征 \(p\) 中 \(I=(t^p)\) 满足上述条件，因为二项式中间系数为零；所得闭子群概形有非零幂零坐标，不能仅由其 \(k\) 点辨认。

\subsection{一般线性群及可对角化群}
\label{b5:subsec:2102:03}

一般线性群的坐标代数为
\[
k[\operatorname{GL}_n]
=k[x_{ij},d]/\bigl(d\det(x_{ij})-1\bigr).
\]
对任意 \(R\)，它的 \(R\) 点就是可逆矩阵。矩阵乘法给出
\[
\Delta(x_{ij})=\sum_\ell x_{i\ell}\otimes x_{\ell j},
\qquad \varepsilon(x_{ij})=\delta_{ij},
\]
对极则是逆矩阵的条目，由伴随矩阵与行列式逆表示。公式 \(\det(AB)=\det A\det B\) 保证 \(d\) 的余乘法为 \(d\otimes d\)。关系 \(\det X=1\) 又给闭子群 \(\operatorname{SL}_n\)。
\begin{definition}\label{b5:def:diagonalizable-group}
设 \(k\) 为域，\(M\) 为有限生成交换群。群代数 \(k[M]\) 以符号 \(e^m\) 为基，乘法为 \(e^me^n=e^{m+n}\)，余乘法为 \(\Delta(e^m)=e^m\otimes e^m\)。相应的交换仿射群概形记为 \(D(M)\)，称为分裂\term[keduijiaohua-qun]{可对角化群}[diagonalizable group]。当 \(M=\mathbb Z^r\) 时，它是分裂\term[daishu-huanmian]{代数环面}[algebraic torus] \(\mathbb G_m^r\)。
\end{definition}
事实上 \(D(M)(R)=\operatorname{Hom}_{\mathrm{grp}}(M,R^\times)\)，因为代数同态必须把 \(e^m\) 送到相容的单位。特别地，
\[
k[\mathbb G_m]=k[t,t^{-1}],\qquad
D(\mathbb Z/n)=\mu_n,\quad k[\mu_n]=k[t]/(t^n-1).
\]
若 \(k\) 特征为 \(p\)，则 \(t^p-1=(t-1)^p\)，\(\mu_p\) 的坐标环不约化；即使 \(k\) 代数闭，它的 \(k\) 点也只有单位。与此相对，常数群 \(C_p\) 的坐标环为 \(k^{C_p}\)，约化且有 \(p\) 个 \(k\) 点。这是两个不同群概形，不能只因为都带下标 \(p\) 就认同。

\ProseParagraphBreak
可逆矩阵的射影类组成\term[sheying-yiban-xianxingqun]{射影一般线性群}[projective general linear group]
\(\operatorname{PGL}_n\)。作为概形，它是射影矩阵空间
\(\mathbb P_k^{n^2-1}\) 中行列式不为零的仿射开子概形
\(D_+(\det)\)，乘法由矩阵乘法给出，逆由伴随矩阵给出。
在任意交换 \(k\) 代数 \(R\) 上，射影类可在开覆盖上用可逆矩阵表示，重叠处的代表相差一个单位标量；当 \(R\) 为域时，每个类都有全局矩阵代表。
下面的局部计算说明，特殊线性群到这个群的商映射也有坐标环层面的含义。
\begin{proposition}\label{b5:prop:special-projective-quotient}
设 \(k\) 为域，\(n\ge2\)。取矩阵的射影类给出群概形态射
\[
q:\operatorname{SL}_n\longrightarrow\operatorname{PGL}_n.
\]
该态射有限且忠实平坦，核为中心子群概形
\(\mu_n\)，其元素以标量矩阵嵌入。对任意有限维 \(k\) 空间 \(V\)，一个群概形态射
\(r:\operatorname{SL}_n\to\operatorname{GL}(V)\)
唯一地因子化为某个群概形态射
\(\bar r:\operatorname{PGL}_n\to\operatorname{GL}(V)\) 与 \(q\) 的复合，当且仅当 \(r\) 在 \(\mu_n\) 上平凡。
因此 \(\operatorname{PGL}_n\) 是此中心商，记为
\(\operatorname{SL}_n/\mu_n\)。
\end{proposition}
\begin{proof}
在 \(\operatorname{PGL}_n\) 的一个标准仿射开集
\(U=\operatorname{Spec}B\) 上，将某个非零射影坐标归一化为一，得到通用矩阵 \(M\)，其行列式 \(D\) 在 \(B\) 中可逆。
\(q^{-1}(U)\) 上的矩阵恰为 \(tM\)，条件是 \(t^nD=1\)，故其坐标环为
\[
C=B[t]/(t^n-D^{-1}).
\]
它是以 \(1,t,\ldots,t^{n-1}\) 为基的自由 \(B\) 模，而且 \(t\) 可逆。
这些开集覆盖目标，所以 \(q\) 有限且忠实平坦；特别地，它在概形上满射。
单位射影类的原像是 \(tI\)、\(t^n=1\)，即 \(\mu_n\)。
在上述开集上，\(C\otimes_BC\) 中的两个标量参数 \(t,u\) 满足同一个 \(n\) 次幂方程，故 \(z=u/t\) 满足 \(z^n=1\)。这给出
\[
C\otimes_BC\cong C[z]/(z^n-1),\qquad u=tz.
\]
这些同构在重叠处相容，因而
\(\operatorname{SL}_n\times\mu_n\cong
\operatorname{SL}_n\times_{\operatorname{PGL}_n}\operatorname{SL}_n\)，
其中两投影分别为 \((g,z)\mapsto g\) 和 \((g,z)\mapsto gz\)。

还须验证因子化保持正则性，而不只是在点群上取商。
在上述坐标环中，\(\mu_n\) 的作用对应
\[
C\longrightarrow C\otimes_k k[z]/(z^n-1),\qquad
t\longmapsto t\otimes z,\quad b\longmapsto b\otimes1.
\]
若 \(c=\sum_{i=0}^{n-1}b_it^i\) 在此作用下不变，比较
\(1,z,\ldots,z^{n-1}\) 的系数便知 \(b_i=0\) 对所有 \(i>0\) 成立。
因此不变子环恰为 \(B\)，这一计算不要求 \(n\) 在 \(k\) 中可逆。
当 \(r\) 在 \(\mu_n\) 上平凡时，\(r\) 的矩阵条目及行列式逆在该作用下不变，因而都是 \(B\) 中的函数。
它们在 \(U\) 上定义到 \(\operatorname{GL}(V)\) 的态射。
重叠处的两组函数拉回后相同，而这些自由扩张是忠实平坦的，故原函数也相同；于是局部态射拼成唯一的 \(\bar r\)。
在 \(q\times q\) 下拉回后，\(\bar r\) 的乘法相容性就是 \(r\) 的乘法相容性，忠实平坦性再使该等式在目标上成立。
故 \(\bar r\) 是群概形态射。反向若已有因子化，则 \(\ker q=\mu_n\) 显然作用平凡。
\end{proof}

% !TeX root = ../../Book5.tex
% 纲要标题：### 21.3 有理表示
% 纲要位置：工作记录/计划纲要.md，第 4403 行。

\section{有理表示}
\label{b5:sec:2103}


% 纲要范围：
%
% * 余模与有理表示
% * 权和代数环面
% * 与 Lie 代数表示的关系及积分问题
%

代数群表示不能只要求每个群元素给出一个矩阵，还要让矩阵条目随群元素正则地变化。余模的系数恰好记录这些函数，因此能把“有理表示”的条件写成有限个代数等式。这里的有理表示是正则群态射，不是允许任意带极点的有理映射。

\subsection{余模与有理表示}
\label{b5:subsec:2103:01}

\begin{definition}\label{b5:def:rational-representation}
设 \(G=\operatorname{Spec}A\) 为域 \(k\) 上的有限型仿射群概形，\(V\) 为有限维 \(k\) 空间。群概形态射 \(G\to\operatorname{GL}(V)\) 称为一个有限维\term[youli-biaoshi]{有理表示}[rational representation]。
\end{definition}
\begin{theorem}\label{b5:thm:rational-comodule}
设 \(k\) 为域，\((A,\Delta,\varepsilon,S)\) 为有限生成的交换 \(k\) Hopf 代数，\(G=\operatorname{Spec}A\)。有限维有理 \(G\) 表示与有限维右 \(A\) 余模等价。具体地，给定有限维右余模 \((V,\delta)\) 及交换含幺 \(k\) 代数 \(R\)，若在 \(V\) 的一组 \(k\) 基 \(v_1,\ldots,v_n\) 下有
\[
\delta(v_j)=\sum_i v_i\otimes a_{ij},
\]
则对应的 \(R\) 点 \(g:A\to R\) 在 \(V\otimes_k R\) 上的作用矩阵为 \(\bigl(g(a_{ij})\bigr)\)，且
\[
\Delta(a_{ij})=\sum_\ell a_{i\ell}\otimes a_{\ell j},\qquad
\varepsilon(a_{ij})=\delta_{ij}.
\]
\end{theorem}
\begin{proof}
余模的两条公理分别给出所列矩阵系数等式。对 \(R\) 点求值，前式表示乘法相容，后式表示单位作用为恒等。对极给出逆矩阵 \(\bigl(S(a_{ij})\bigr)\)，因此这些矩阵在每个 \(R\) 上都可逆，定义到 \(\operatorname{GL}_n\) 的态射。反过来，群态射的矩阵系数必满足同样等式，定义余作用。换基只是同时改变这组系数的表达，同态相容条件也相同，所以得到基无关的范畴等价。
\end{proof}
\begin{lemma}\label{b5:lem:comodule-locally-finite}
设 \(k\) 为域，\(C\) 为 \(k\) 余代数，\(V\) 为右 \(C\) 余模，不要求有限维。每个 \(v\in V\) 都属于一个有限维 \(k\) 子余模。
\end{lemma}
\begin{proof}
写 \(\delta(v)=\sum_{i=1}^r v_i\otimes c_i\)，令 \(c_i\) 线性无关。选线性泛函 \(\phi_i\) 满足 \(\phi_i(c_j)=\delta_{ij}\)。令
\(W=\operatorname{span}\{v_i\}\)，亦即所有 \((1\otimes\phi)\delta(v)\) 的线性生成空间。
余结合律在最后一个因子上应用任意 \(\phi\)，给出
\[
\delta\bigl((1\otimes\phi)\delta(v)\bigr)
=\sum_i v_i\otimes(1\otimes\phi)\Delta(c_i)\in W\otimes C.
\]
故 \(W\) 稳定；余单位公式使 \(v=\sum_i\varepsilon(c_i)v_i\in W\)。
\end{proof}
\begin{corollary}\label{b5:thm:affine-group-linear}
域 \(k\) 上每个有限型仿射群概形都可以闭嵌入某个 \(\operatorname{GL}_n\)。
\end{corollary}
\begin{proof}
将 \(A=k[G]\) 本身看作余作用为 \(\Delta\) 的右余模。取 \(A\) 的有限个代数生成元，利用前一引理，选包含这些元及一的有限维子余模 \(W\)。在 \(W\) 的基 \(w_j\) 上写系数为 \(a_{ij}\)。这里使用 \(A\) 作为余代数的另一侧余单位等式 \((\varepsilon\otimes1)\Delta=1_A\)，而非右余模的 \((1\otimes\varepsilon)\Delta=1_A\)。由
\(w_j=(\varepsilon\otimes1)\Delta(w_j)=\sum_i\varepsilon(w_i)a_{ij}\)，
所有 \(w_j\) 都在这些矩阵系数的线性生成空间中。因此表示给出的
\(k\bigl[\operatorname{GL}(W)\bigr]\to A\) 的像包含全部代数生成元，是满射；这就是闭嵌入。
\end{proof}
这项证明同时说明有限生成坐标环的用途。余模局部有限性保证任何一个函数都受有限维矩阵系数控制，而有限生成性保证一次选择便能控制整个坐标环。

\subsection{权和代数环面}
\label{b5:subsec:2103:02}

\begin{theorem}\label{b5:thm:torus-weight-decomposition}
设 \(k\) 为任意特征的域，\(M\) 为有限生成交换群，\(G=D(M)=\operatorname{Spec}k[M]\)，\(e^m\) 为群代数 \(k[M]\) 中元素 \(m\in M\) 对应的基元。对有限维有理 \(G\) 表示，以及更一般的任意右 \(k[M]\) 余模 \((V,\delta)\)，令 \(V_m=\{v\in V\mid\delta(v)=v\otimes e^m\}\)。则有唯一的分解
\[
V=\bigoplus_{m\in M}V_m,\qquad
\delta(v)=v\otimes e^m\quad(v\in V_m).
\]
有限维时只有有限多个非零项。因此每个有理表示完全可约，简单表示是一维权；结论在任意特征下成立。
\end{theorem}
\begin{proof}
群代数的 \(e^m\) 是一组基，对每个 \(v\) 唯一写
\(\delta(v)=\sum_m v_m\otimes e^m\)，有限多个系数非零。余结合比较第三因子的 \(e^m\) 系数，得到
\(\delta(v_m)=v_m\otimes e^m\)；余单位使 \(v=\sum_m v_m\)。
若一组不同权向量的和为零，对其余作用比较不同 \(e^m\) 系数便知每项为零，所以是直和。表示同态保各权空间，一维子空间在同一权空间内都稳定，故分解为一维表示。唯一性也来自系数唯一性。
\end{proof}
对 \(\mathbb G_m^r\)，权是整数元组 \((m_1,\ldots,m_r)\)，作用字符为
\(t_1^{m_1}\cdots t_r^{m_r}\)。这不是任意复指数。比如一维交换 Lie 代数的任意复标量作用都给出 Lie 代数表示，却只有整数标量才能从 \(\mathbb G_m\) 的有理表示微分得到。
\ProseParagraphBreak
即使 \(G=\mu_p\) 在特征 \(p\) 中不约化，上述按 \(\mathbb Z/p\) 分次的结论仍成立。它与常数群 \(C_p\) 的模特征表示完全不同：后者的群代数是局部非半单代数，而这里讨论的是不同的群概形及其坐标余模。

\subsection{与 Lie 代数表示的关系及积分问题}
\label{b5:subsec:2103:03}

把一个 \(R\) 点在单位附近代入双数环 \(k[\epsilon]/(\epsilon^2)\)，便得到无穷小作用。具体地，对 \(\varepsilon:A\to k\)，线性映射 \(d:A\to k\) 满足
\[
d(ab)=\varepsilon(a)d(b)+d(a)\varepsilon(b)
\]
时，\(\varepsilon+\epsilon d\) 为双数值代数同态。这些 \(d\) 构成 \(\operatorname{Lie}(G)\) 的底向量空间。对两线性映射 \(d_1,d_2:A\to k\)，卷积为
\[
(d_1*d_2)(a)=\sum d_1(a_{(1)})d_2(a_{(2)}),
\qquad \Delta(a)=\sum a_{(1)}\otimes a_{(2)}.
\]
若 \(d_1,d_2\) 都满足上述导子等式，利用 \(\Delta\) 的乘法相容性与余单位等式，展开得到
\[
\begin{aligned}
(d_1*d_2)(ab)
&=\varepsilon(a)(d_1*d_2)(b)+(d_1*d_2)(a)\varepsilon(b)\\
&\quad+d_1(a)d_2(b)+d_2(a)d_1(b).
\end{aligned}
\]
交换 \(d_1,d_2\) 后，最后两项不变，因而卷积交换子 \([d_1,d_2]=d_1*d_2-d_2*d_1\) 仍满足导子等式。卷积的结合律使这个括号满足交替性与 Jacobi 恒等式，所以确实定义了 \(\operatorname{Lie}(G)\) 的 Lie 代数结构。

\ProseParagraphBreak
有理表示的余作用诱导 \(\rho_*(d)=(1\otimes d)\delta\)。余结合律给出
\[
\begin{aligned}
\rho_*(d_1)\rho_*(d_2)
&=(1\otimes d_1\otimes d_2)(\delta\otimes1)\delta\\
&=(1\otimes d_1\otimes d_2)(1\otimes\Delta)\delta\\
&=(1\otimes(d_1*d_2))\delta.
\end{aligned}
\]
两边交换 \(d_1,d_2\) 后相减，便有
\[
[\rho_*(d_1),\rho_*(d_2)]=\rho_*([d_1,d_2]).
\]
因此这是一个 Lie 代数表示。坐标方程在 \(I+\epsilon X\) 处线性化便恢复熟悉的矩阵条件，例如 \(\operatorname{Lie}(\operatorname{SL}_n)=\mathfrak{sl}_n\)。
\begin{proposition}\label{b5:prop:additive-rational-representations}
设 \(k\) 为特征零域，\(V\) 为 \(n\) 维 \(k\) 向量空间。加法群 \(\mathbb G_a\) 在 \(V\) 上的每个有理表示，在选定 \(k\) 基后唯一形如
\[
\rho(t)=\exp(tN)=\sum_{j\ge0}\frac{t^jN^j}{j!},
\]
其中 \(N\) 为幂零矩阵，故该和有限。每个有限维有理 \(\mathbb G_m\) 表示的微分则可对角化，特征值为整数。
\end{proposition}
\begin{proof}
加法群的矩阵条目为 \(t\) 的多项式，满足
\(\rho(t+u)=\rho(t)\rho(u)\)、\(\rho(0)=I\)。
对 \(u\) 求导并令 \(u=0\)，得到 \(\rho'(t)=\rho(t)N\)，其中 \(N=\rho'(0)\)。逐次比较多项式系数得到 \(A_j=N^j/j!\)。多项式次数有限，所以充分大的 \(N^j=0\)。反过来，幂零 \(N\) 的有限指数满足加法公式，给出有理表示。乘法群的结论由\cref{b5:thm:torus-weight-decomposition}中 \(t^m\) 的导数得到。
\end{proof}
因此从 Lie 代数表示到代数群有理表示会多出条件；即使两群的 Lie 代数都一维交换，\(\mathbb G_a\) 要求幂零作用，\(\mathbb G_m\) 要求整数半单作用。复 Lie 群的积分问题还涉及拓扑覆盖，留在附录C说明。在特征 \(p\) 中，群态射 \(t\mapsto t^p\) 的微分为零却不是零态射，微分还会进一步丢掉信息。

% !TeX root = ../../Book5.tex
% 纲要标题：### 21.4 代数群结构初步
% 纲要位置：工作记录/计划纲要.md，第 4409 行。

\section{代数群结构初步}
\label{b5:sec:2104}


% 纲要范围：
%
% * 幺幂群和约化群
% * 根数据的意义
% * 正特征行为和群概形方向
%
% ---
%

从矩阵例子进入一般代数群时，需要区分几个不同的结构条件。幺幂性控制元素的 Jordan 形，约化性排除连通正规幺幂部分；根数据又在根系以外保留了整数格。正特征时，坐标环的幂零信息和无穷小子群还会使点集直觉失效。

\subsection{幺幂群和约化群}
\label{b5:subsec:2104:01}

以下先在代数闭域 \(k\) 上讨论光滑仿射代数群，闭子群也取光滑闭子群。矩阵 \(u\) 若满足 \(u-I\) 幂零，就称为幺幂矩阵。一个光滑线性代数群在一个忠实表示中全部 \(k\) 点都幺幂时，称为幺幂群；这个概念的表示无关性是一般线性代数群理论的一项结构结论。本节的实例均在明确的矩阵模型中使用它。
\ProseParagraphBreak
上三角对角线全为一的群 \(U_n\) 是基本例子。令 \(J\) 为上三角矩阵代数中的严格上三角矩阵理想，\(U_n^{(r)}=1+J^r\)，即仅在距对角线至少 \(r\) 条的上对角线上允许非零条目的子群。因为 \(J\) 幂零，\(1+J^{r+s}\) 是 \(1+J\) 的正规子群。在商群中，若 \(X\in J^r\)、\(Y\in J^s\)，则 \(XY,YX\in J^{r+s}\)，所以
\[
(1+X)(1+Y)\equiv1+X+Y\pmod{J^{r+s}}
\quad(X\in J^r,\ Y\in J^s),
\]
与逆序相乘得到的 \((1+Y)(1+X)\) 相同。因此两个元素在商群中交换，其交换子属于 \(1+J^{r+s}\)，即 \(\bigl[U_n^{(r)},U_n^{(s)}\bigr]\subseteq U_n^{(r+s)}\)。因为 \(J^n=0\)，这个群幂零，特别可解。对 \(U_3\)，三个参数的乘法为
\[
(a,b,c)(a',b',c')
=(a+a',b+b',c+c'+ab'),
\]
可见中心方向记录了非交换性，不能把全部上三角参数都简单看成独立加法。
\begin{definition}\label{b5:def:reductive-group}
设 \(G\) 为代数闭域上的光滑连通仿射代数群。若 \(G\) 没有非平凡的光滑连通正规幺幂闭子群，就称为\term[yuehua-qun]{约化群}[reductive group]。若此外没有正维数的光滑连通正规可解闭子群，就称为半单代数群。这里的“约化群”不同于“坐标环约化”：后者只要求没有非零幂零元。
\end{definition}
例如环面没有非平凡幺幂元素，所以约化；它自身交换且有正维数时不半单。\(U_n\) 在 \(n\ge2\) 时光滑连通且非平凡，却是幺幂的，因而不约化。一般约化群的结构还涉及环面、根子群和正规子群；下一小节从整数格资料观察其中一部分。

\subsection{根数据的意义}
\label{b5:subsec:2104:02}

\begin{definition}\label{b5:def:root-datum}
设 \(X,Y\) 为两个有限秩自由交换群，给定完美整数配对
\(\langle\ ,\ \rangle:X\times Y\to\mathbb Z\)，以及有限子集
\(\Phi\subseteq X,\Phi^\vee\subseteq Y\) 及双射
\(\alpha\mapsto\alpha^\vee\)。要求对所有 \(\alpha\in\Phi\)，都有
\(\langle\alpha,\alpha^\vee\rangle=2\)，且反射
\[
s_\alpha(x)=x-\langle x,\alpha^\vee\rangle\alpha,\qquad
s_\alpha^\vee(y)=y-\langle\alpha,y\rangle\alpha^\vee
\]
分别保持 \(\Phi\) 和 \(\Phi^\vee\)，并对所有 \(\alpha,\beta\in\Phi\) 满足 \((s_\alpha\beta)^\vee=s_\alpha^\vee(\beta^\vee)\)。还要求约化条件 \(\Phi\cap\mathbb Q\alpha=\{\alpha,-\alpha\}\) 对每个 \(\alpha\in\Phi\) 成立，其中把 \(X\) 嵌入 \(X\otimes_{\mathbb Z}\mathbb Q\)。满足这些条件的资料称为\term[gen-shuju]{根数据}[root datum]；允许 \(\Phi\) 为空，也不要求它生成整个 \(X\otimes\mathbb Q\)，以便包含环面的中心方向。
\end{definition}
\symidx{\(X,Y,\Phi,\Phi^\vee\)：根数据中的字符格、余字符格、根与余根集合}
在连通约化群的结构理论中，若 \(T\) 是极大环面，这两个格具体为
\[
X=\operatorname{Hom}_{\mathrm{alg.grp}}(T,\mathbb G_m),\qquad
Y=\operatorname{Hom}_{\mathrm{alg.grp}}(\mathbb G_m,T).
\]
这里的同态是代数群态射，\(\chi\in X\) 称为字符，\(\nu\in Y\) 称为一参数子群。它们的复合是 \(\mathbb G_m\) 的群态射，故形如 \(t\mapsto t^n\)，其唯一整数指数就是配对：
\[
\chi(\nu(t))=t^{\langle\chi,\nu\rangle}.
\]
这一整数幂结论也可直接由 Laurent 多项式的 Hopf 等式看出：类群元只能是单项式 \(t^n\)。根系给出实向量空间里的角度与相对长度，根数据还记录哪些整数权给出极大环面上的字符。这些字符不必延拓为整个群的一维表示；例如 \(\operatorname{SL}_2\) 的环面有非平凡字符，而 \(\operatorname{SL}_2\) 自身没有非平凡的一维有理表示。因此同一个复半单 Lie 代数可以对应不同的代数群。
\ProseParagraphBreak
对 \(\operatorname{SL}_2\)，取 \(T=\bigl\{\operatorname{diag}(t,t^{-1})\bigr\}\)。若 \(\lambda(t)=t\)，则
\(X=\mathbb Z\lambda\)，根为 \(\pm2\lambda\)；一参数子群
\(\nu(t)=\operatorname{diag}(t,t^{-1})\)
满足 \(\langle\lambda,\nu\rangle=1\)，正余根就是 \(\nu\)。
对 \(\operatorname{PGL}_2\) 的对角环面，用对角条目比值 \(r\) 作坐标，字符格由 \(\alpha(r)=r\) 生成，根为 \(\pm\alpha\)。若 \(\eta(t)\) 是 \(\operatorname{diag}(t,1)\) 的射影类，则 \(\langle\alpha,\eta\rangle=1\)。商映射把 \(\operatorname{SL}_2\) 环面的参数送为 \(r=t^2\)，所以字符拉回为 \(\alpha\mapsto2\lambda\)，正余根 \(\nu\) 的像则为 \(2\eta\)。于是 \(\operatorname{PGL}_2\) 的余根是格 \(Y\) 的生成元的二倍，仍满足根余根配对为二。两者的抽象根系都为 \(A_1\)，但整数格中的嵌入不同。
\ProseParagraphBreak
这也解释了表示下降的一个必要条件。中心元 \(-I\in\operatorname{SL}_2(\mathbb C)\) 在
\(L_m=\operatorname{Sym}^m(\mathbb C^2)\)
上作用为 \((-1)^m\)。由\cref{b5:prop:special-projective-quotient}，代数群商 \(\operatorname{SL}_{2,\mathbb C}/\mu_2\) 同构于 \(\operatorname{PGL}_{2,\mathbb C}\)，它在复点上的商群也记为 \(\operatorname{PSL}_2(\mathbb C)\)。所以这个有理表示能下降到该商代数群当且仅当 \(m\) 为偶数。Lie 代数的最高权分类列出所有 \(m\ge0\)，却不会自动替每个中心商筛去不能下降的表示。

\subsection{正特征下的代数群与群概形}
\label{b5:subsec:2104:03}

在特征 \(p\) 中，对加法群的 Frobenius 态射
\[
F:\mathbb G_a\longrightarrow\mathbb G_a,\qquad t\longmapsto t^p
\]
对应坐标同态 \(k[t]\to k[t]\)、\(t\mapsto t^p\)。它的核为
\(\alpha_p=\operatorname{Spec}k[t]/(t^p)\)。
\symidx{\(\alpha_p\)：特征 \(p\) 中加法群 Frobenius 态射的核}
若 \(k\) 代数闭，点集上 \(F\) 为双射，但坐标同态不满，所以它不是代数群同构。其微分为零；核概形的切空间却一维，尽管它的 \(k\) 点只有一个。
\ProseParagraphBreak
同样，\(\mu_p=\operatorname{Spec}k[t,t^{-1}]/(t^p-1)\) 的切空间一维；在 \(t=1+\epsilon a\) 代入关系时，
\((1+\epsilon a)^p=1\)，对任意 \(a\) 都成立。于是 \(\alpha_p,\mu_p\) 与平凡点集之间的区别，必须由坐标环或所有代数上的点保留。
\ProseParagraphBreak
Frobenius 态射与非约化核说明，光滑性、约化性、可分性以及表示是否能由微分恢复，分别控制不同信息。后续的不变量理论将区分特征零的线性约化性与正特征下更弱的几何约化性；有限群的平均公式只在群阶可逆时适用，一般群概形需要相应的表示论方法。

% !TeX root = ../../Book5.tex
% 纲要标题：### 21.5 型、对合与古典群
% 纲要位置：工作记录/计划纲要.md，第 4415 行。
% 本轮补充的纲要细目：
% ---
% 写作范围：具体型、对合及四元数模型；不外推为一般扭曲形式分类。

\section{型、对合与古典群}
\label{b5:sec:2105}

% 纲要细目：
% * 从非退化型的分裂矩阵模型解释正交、辛、酉群及相似变换群，写清底域、特征和型的约定
% * 区别等距群、相似变换群、代数自同构群与中心商；代数群满态射不保证有理点映射满射
% * 选一个四元数的非分裂模型，并重述所需代数事实；不以一般扭曲形式分类或未写的非交换上同调作为前提
% #### 21.5.1 等距群、相似变换群与乘子
% #### 21.5.2 带对合代数的自同构群
% #### 21.5.3 分裂模型与四元数例子

正交群、辛群和酉群最初由保持型的矩阵定义。把型换成自同态代数上的伴随对合以后，同一批方程也能在非分裂代数中使用。本节先固定具体模型，再比较等距、相似变换和代数自同构；它们之间有联系，却不是同一个群。

\subsection{等距群、相似变换群与乘子}
\label{b5:subsec:2105:01}

设 \(F\) 为特征不为二的域，取 \(K=F\) 或一个可分二次扩张 \(K/F\)，并以 \(a\mapsto\bar a\) 分别表示恒等或非平凡共轭。给定 \(H\in\operatorname{GL}_n(K)\)、\(n\ge1\)，假设
\(H^*=\epsilon H\)、\(\epsilon\in\{1,-1\}\)，其中 \(X^*=\bar X^{\mathsf T}\)。型为
\(b(u,v)=\bar u^{\mathsf T}Hv\)，伴随对合为
\(\sigma_H(X)=H^{-1}X^*H\)。
\begin{definition}\label{b5:def:classical-similitude}
设 \(F\) 为特征不为二的域，\(K=F\) 或 \(K/F\) 为可分二次扩张，共轭分别为恒等或非平凡域自同构。取 \(H\in\operatorname{GL}_n(K)\)、\(n\ge1\)，满足 \(H^*=\epsilon H\)、\(\epsilon\in\{1,-1\}\)，其中 \(X^*=\bar X^{\mathsf T}\)。对非退化型 \(b(u,v)=\bar u^{\mathsf T}Hv\)，定义等距群与相似变换群
\[
\begin{aligned}
\operatorname{Isom}(b)&=\bigl\{g\in\operatorname{GL}_n(K)\mid g^*Hg=H\bigr\},\\
\operatorname{Sim}(b)&=\bigl\{g\in\operatorname{GL}_n(K)\mid
g^*Hg=cH\text{ 对某个 }c\in F^\times\bigr\}.
\end{aligned}
\]
相似变换的这个唯一标量 \(c\) 称为乘子，记为 \(\mu(g)\)。当 \(K=F\)、\(H\) 对称时，两群分别记为 \(\operatorname O(b)\)、\(\operatorname{GO}(b)\)；当 \(H\) 交替时，记为 \(\operatorname{Sp}(b)\)、\(\operatorname{GSp}(b)\)；二次扩张上的 Hermitian 情形记为 \(\operatorname U(b)\)、\(\operatorname{GU}(b)\)。
\end{definition}
这些式子不仅定义 \(F\) 上的抽象群。对任意交换 \(F\) 代数 \(R\)，把系数扩大到 \(K\otimes_FR\)，令共轭固定 \(R\)，便得到相同的矩阵方程；将非零乘子换为 \(R^\times\) 中的元，就定义相应仿射群概形。等距群是一般线性群的闭子群，相似群可在 \(\operatorname{GL}_n(K)\times\mathbb G_m\) 中用方程实现；\(K/F\) 有限维使这些都是有限个 \(F\) 坐标方程。
\begin{proposition}\label{b5:prop:classical-multiplier}
设 \(F\) 为特征不为二的域，\(K=F\) 或 \(K/F\) 为可分二次扩张，共轭分别为恒等或非平凡域自同构，\(X^*=\bar X^{\mathsf T}\)。取 \(n\ge1\)、\(H\in\operatorname{GL}_n(K)\)，满足 \(H^*=\epsilon H\)、\(\epsilon=\pm1\)。令 \(b\) 为 \(K^n\) 上的非退化型 \(b(u,v)=\bar u^{\mathsf T}Hv\)，\(\sigma_H(X)=H^{-1}X^*H\)，记 \(\operatorname{Sim}(b)\)、\(\operatorname{Isom}(b)\) 分别为相似变换群和等距群。则乘子给出群同态 \(\mu:\operatorname{Sim}(b)\to F^\times\)，其核为 \(\operatorname{Isom}(b)\)。并且
\[
g\in\operatorname{Sim}(b),\ \mu(g)=c
\quad\Longleftrightarrow\quad \sigma_H(g)g=cI.
\]
\end{proposition}
\begin{proof}
非退化型至少有一个非零配对值，故乘子唯一。两个相似变换连续作用使配对值乘两个标量的积，逆变换乘以逆标量，所以是同态，乘子一恰是等距条件。最后左乘 \(H^{-1}\) 即将矩阵方程改写为所示对合方程。
\end{proof}
乘子的实际像须另外判断。在正定实型上它只能为正数，且 \(\sqrt c I\) 实现每个正数；在分裂辛型
\(J=\begin{psmallmatrix}0&I_m\\-I_m&0\end{psmallmatrix}\)
上，\(\operatorname{diag}(cI_m,I_m)\) 实现每个 \(c\in F^\times\)。所以写核与像时不能一概补上“到 \(F^\times\) 满射”。

\subsection{带对合代数的自同构群}
\label{b5:subsec:2105:02}

\begin{proposition}\label{b5:prop:involution-automorphism-similitude}
设 \(F\) 为特征不为二的域，\(K=F\) 或 \(K/F\) 为可分二次扩张，共轭分别为恒等或非平凡域自同构，\(X^*=\bar X^{\mathsf T}\)。取 \(H\in\operatorname{GL}_n(K)\)、\(n\ge1\)，满足 \(H^*=\epsilon H\)、\(\epsilon=\pm1\)。令 \(A=M_n(K)\)、\(b(u,v)=\bar u^{\mathsf T}Hv\) 为 \(K^n\) 上的非退化型、\(\sigma(X)=H^{-1}X^*H\)，\(\operatorname{Sim}(b)\) 为型 \(b\) 的相似变换群。保持 \(\sigma\) 的 \(K\) 代数自同构组成的群记为
\(\operatorname{Aut}_K(A,\sigma)\)。内自同构
\(\operatorname{Int}(g)(X)=gXg^{-1}\)
给出群同构
\[
\operatorname{Sim}(b)/K^\times
\ \cong\ \operatorname{Aut}_K(A,\sigma),
\]
其中 \(K^\times\) 作为非零标量矩阵嵌入相似群，其乘子为 \(\bar a a\)。
\end{proposition}
\begin{proof}
矩阵代数的每个 \(K\) 代数自同构都是内自同构，这是第二卷定理17.3.1的 Skolem--Noether 结论在中心单代数 \(M_n(K)\) 上的应用。直接计算
\[
\sigma(gXg^{-1})=\sigma(g)^{-1}\sigma(X)\sigma(g).
\]
它等于 \(g\sigma(X)g^{-1}\) 对所有 \(X\) 成立，当且仅当
\(\sigma(g)g\) 与所有矩阵交换，即为某个 \(cI\)、\(c\in K^\times\)。对该等式再用 \(\sigma\)，得到 \(\bar c=c\)，故 \(c\in F^\times\)。由前一命题，这正是相似条件。内自同构的核是矩阵代数的中心单位 \(K^\times\)，而这些标量全属于相似群，遂得所示同构。
\end{proof}
这里比较的是指定域上的点群，并且自同构要求逐点固定 \(K\)。二次扩张上若允许自同构在中心上也作共轭，会得到另一个群。对一般带对合中心单代数，相同的计算仍把内自同构与条件 \(\sigma(g)g\in F^\times\) 联系起来；非分裂情形中的型及对合细节见第二卷附录B。
\ProseParagraphBreak
还须区分代数群满态射与域上点的满射。复数上，标量平方根能把任意可逆二阶矩阵缩放为行列式一矩阵，所以
\(\operatorname{SL}_2(\mathbb C)\to\operatorname{PGL}_2(\mathbb C)\)
满射。相应代数群态射在实数上也定义良好，但实点的像只含行列式为正的射影矩阵类：缩放把行列式乘一个正平方，不可能把负行列式改成一。\(\operatorname{diag}(-1,1)\) 的射影类便没有实的 \(\operatorname{SL}_2\) 原像。

\subsection{分裂模型与四元数例子}
\label{b5:subsec:2105:03}

在 \(F=\mathbb R\) 上，Hamilton 四元数代数
\[
\mathbb H=\mathbb R\oplus\mathbb Ri\oplus\mathbb Rj\oplus\mathbb Rk,
\quad i^2=j^2=k^2=-1,\quad ij=k=-ji
\]
带有共轭 \(\bar q\) 和范数
\(N(q)=q\bar q=q_0^2+q_1^2+q_2^2+q_3^2\)。
共轭反转乘法，故 \(N(uv)=N(u)N(v)\)。于是
\[
\operatorname{SL}_1(\mathbb H)(\mathbb R)=\bigl\{q\in\mathbb H\mid N(q)=1\bigr\}
\]
由一个四元二次方程定义为实代数群；逆由共轭给出，不需要除以变量。
\ProseParagraphBreak
复化后它变成分裂模型。把 \(i,j\) 分别送到
\[
\begin{pmatrix}\mathrm i&0\\0&-\mathrm i\end{pmatrix},
\qquad
\begin{pmatrix}0&1\\-1&0\end{pmatrix}.
\]
两矩阵平方为 \(-I\)、反交换，并且它们连同单位及乘积线性无关，所以给出
\(\mathbb H\otimes_{\mathbb R}\mathbb C\cong M_2(\mathbb C)\)。
计算一般元素的行列式恰为 \(N(q)\)，因此复化后的群为 \(\operatorname{SL}_2\)。在实点上，四元数单位球面是有界闭集，而 \(\operatorname{SL}_2(\mathbb R)\) 含任意大的 \(\operatorname{diag}(t,t^{-1})\)；这两个实模型的差异不由其共同的复化消除。
\begin{example}\label{b5:exm:quaternion-conjugation-group}
单位四元数通过 \(x\mapsto qxq^{-1}\) 作用于纯虚三维空间
\(\operatorname{Im}\mathbb H\)。它保持该空间，因为实部是 \((x+\bar x)/2\)，且共轭作用保持实部；它也保持正定范数。作用的核是 \(\{\pm1\}\)：若与 \(i,j,k\) 都交换，便在 \(\mathbb H\) 的中心 \(\mathbb R\) 中，再加范数一就只剩两个元。
例如 \(q=i\) 固定 \(i\) 并把 \(j,k\) 变号；而
\(q=(1+i)/\sqrt2\) 固定 \(i\)，将 \(j\) 送到 \(k\)、\(k\) 送到 \(-j\)。这是由非交换代数乘法产生的具体正交作用。
\end{example}
分裂矩阵型、非分裂四元数、等距群与代数自同构因而可以互相解释。它们也提醒我们：复数域上的根系或 Lie 代数类型只记录部分结构；改换基域、对合或中心商以后，点群与表示下降问题还需要重新核对。


\begin{chaptersummary}
余乘法规定张量积的作用，余单位规定平凡对象，对极使对偶和求逆成为可能。群代数中的元素按对角方式作用，包络代数中的 Lie 代数生成元按 Leibniz 方式作用；坐标代数则以反向的函数拉回描述群乘法。

\ProseParagraphBreak
有理表示等同于坐标余模，有限生成坐标环可由有限维矩阵系数恢复。环面表示按整数权分解，而微分只记录这些字符在单位处的一阶变化；正特征的 Frobenius 和非约化核说明一阶变化可能遗漏更多信息。

\ProseParagraphBreak
型的等距群是相似群的乘子核，保持伴随对合的内自同构则还要模去中心标量。四元数群在复化后成为熟悉的矩阵群，但其实际底域上的点仍可大不相同。这些区别是下一章讨论不变量与商时需要保留的对象条件。
\end{chaptersummary}

\BookExercises
\begin{optionalexercise}\label{b5:ex:21:01}
设 \((H,\Delta,\varepsilon,S)\) 为域 \(k\) 上的 Hopf 代数，证明 \(\varepsilon\bigl(S(h)\bigr)=\varepsilon(h)\)。提示：在 \(\operatorname{Hom}_k(H,k)\) 的卷积中比较余单位及其逆。
\end{optionalexercise}
\begin{solution}
对极等式应用 \(\varepsilon\)，得到 \((\varepsilon S)*\varepsilon=\varepsilon=\varepsilon*(\varepsilon S)\)，而余单位公理说明 \(\varepsilon\) 本来就是这个卷积代数的单位。单位的逆仍为单位，故 \(\varepsilon S=\varepsilon\)。
\end{solution}

\begin{optionalexercise}\label{b5:ex:21:02}
在 \(k[t]\) 上取 \(\Delta(t)=t\otimes1+1\otimes t\)。计算 \(\Delta(t^n)\)，并说明特征 \(p\) 时 \(t^p\) 仍为本原元。特征零时全部本原多项式有哪些？
\end{optionalexercise}
\begin{solution}
余乘法是代数同态，故 \(\Delta(t^n)=\sum_{j=0}^n\binom njt^j\otimes t^{n-j}\)。特征 \(p\) 下中间二项式系数为零，得到本原性。特征零时若非零本原多项式最高次为 \(d>1\)，其余乘法中 \(t^{d-1}\otimes t\) 系数为首项系数乘 \(d\)，非零，矛盾；常数本原要求 \(c=2c\)，所以为零。全部本原元恰为 \(kt\)。
\end{solution}

\begin{optionalexercise}\label{b5:ex:21:03}
设 \(G\) 为群。在 \(kG\) 中求满足 \(\Delta(u)=u\otimes u\)、\(\varepsilon(u)=1\) 的全部元素，称这类元素为类群元。
\end{optionalexercise}
\begin{solution}
写 \(u=\sum a_gg\)。左侧仅有 \(g\otimes g\)，而右侧不同 \(g,h\) 的系数为 \(a_ga_h\)。域中不能有两个非零系数，因此只有一个 \(a_g\ne0\)。比较对角系数得 \(a_g=a_g^2\)，余单位又要求它为一。故类群元恰好是群基 \(G\)。
\end{solution}

\begin{optionalexercise}\label{b5:ex:21:04}
设 \(G=C_2\)，在特征二中比较 \(kG\) 与 \(k^G\) 的底层代数、类群元及幂零性。
\end{optionalexercise}
\begin{solution}
前者同构于 \(k[u]/(u^2)\)，其中 \(u=s-1\)，有非零幂零元；后者为 \(k\times k\)，约化。前者类群元为 \(1,s\)。后者的类群元是函数 \(f(gh)=f(g)f(h)\)、\(f(1)=1\)，即字符 \(G\to k^\times\)；其 \(f(s)^2=1\) 在特征二只有 \(f(s)=1\)，所以只有常数一。两者不能因向量空间维数相同而认同。
\end{solution}

\begin{optionalexercise}\label{b5:ex:21:05}
在 \(U\bigl(\mathfrak{sl}_2(\mathbb C)\bigr)\) 的 Hopf 结构下，求 \(\Omega=h^2+2h+4fe=8C\) 的余乘法，其中 \(C\) 为本书按 Killing 型归一化的 Casimir 元。
\end{optionalexercise}
\begin{solution}
用生成元本原性展开得
\[
\Delta(\Omega)=\Omega\otimes1+1\otimes\Omega
+2h\otimes h+4f\otimes e+4e\otimes f.
\]
交叉项来自平方 \(h^2\) 与乘积 \(fe\)。因此 Casimir 中央并不意味着它是本原元；它在张量积上的特征值不能只把两因子特征值直接相加。
\end{solution}

\begin{optionalexercise}\label{b5:ex:21:06}
设 \(G\) 有限，给出 \(k^G\) 上左模与 \(G\)-分次空间的对应，并解释其张量乘法。
\end{optionalexercise}
\begin{solution}
点函数 \(\delta_g\) 是完整正交幂等元，因此 \(V=\bigoplus_gV_g\)、\(V_g=\delta_gV\)。反过来按各分次片投影定义作用。由 \(\Delta\delta_g=\sum_{ab=g}\delta_a\otimes\delta_b\)，张量积的 \(g\) 片是 \(\bigoplus_{ab=g}V_a\otimes W_b\)。若群非交换，交换张量因子一般不保持这个分次。
\end{solution}

\begin{optionalexercise}\label{b5:ex:21:07}
验证矩阵余代数的公式 \(\Delta(c_{ij})=\sum_\ell c_{i\ell}\otimes c_{\ell j}\)、\(\varepsilon(c_{ij})=\delta_{ij}\)。求其线性对偶代数。
\end{optionalexercise}
\begin{solution}
两次余乘法均为 \(\sum_{\ell,r}c_{ir}\otimes c_{r\ell}\otimes c_{\ell j}\)，重命名指标即一致；余单位直接只留下一个求和项。若对偶基为 \(e^{ij}\)，卷积满足 \(e^{ij}e^{rs}=\delta_{jr}e^{is}\)，所以对偶代数为矩阵代数。
\end{solution}

\begin{optionalexercise}\label{b5:ex:21:08}
在 \(\mathbb G_a\) 的正则余模 \(k[t]\) 中，求包含 \(t^d\) 的一个有限维子余模。特征零时证明最小这样的子余模为次数不超过 \(d\) 的多项式空间。
\end{optionalexercise}
\begin{solution}
二项式公式说明 \(W_d=\operatorname{span}\{1,t,\ldots,t^d\}\) 稳定。特征零时 \(\Delta(t^d)\) 的第二因子各次系数非零且独立，用相应系数泛函作用便得到每个 \(t^j\)，故任一含 \(t^d\) 的子余模都含 \(W_d\)。正特征下有些系数会消失，最小性可能改变。
\end{solution}

\begin{optionalexercise}\label{b5:ex:21:09}
给出 \(\mathbb G_m^2\) 在二维空间上的权 \((1,0),(0,1)\) 作用，计算对偶、二次对称幂和外平方的权。
\end{optionalexercise}
\begin{solution}
对偶的两个权为 \((-1,0)\)、\((0,-1)\)，因为字符取逆。对称平方中，基与权的对应为
\[
v_1^2:\ (2,0),\qquad
v_1v_2:\ (1,1),\qquad
v_2^2:\ (0,2).
\]
外平方 \(v_1\wedge v_2\) 的权为 \((1,1)\)。这些结论直接由字符乘法得到，在任意特征均成立。
\end{solution}

\begin{optionalexercise}\label{b5:ex:21:10}
证明 \(\mathbb G_m\to\mathbb G_m\) 的群概形态射恰为 \(t\mapsto t^n\)、\(n\in\mathbb Z\)。
\end{optionalexercise}
\begin{solution}
坐标同态由可逆 Laurent 多项式的像 \(u\) 决定，Hopf 条件要求其为类群元。写有限 Laurent 和，比较 \(\Delta(u)\) 与 \(u\otimes u\)，如群代数的系数论证可知 \(u=t^n\)。反向每个这样的公式显然满足全部结构等式。
\end{solution}

\begin{optionalexercise}\label{b5:ex:21:11}
在特征零域上，证明代数群同态 \(\mathbb G_a\to\mathbb G_m\) 必为平凡，而 \(\mathbb G_a\to\mathbb G_a\) 恰为 \(t\mapsto ct\)。
\end{optionalexercise}
\begin{solution}
前者对应把 \(k[u,u^{-1}]\) 映到 \(k[t]\)，但 \(k[t]\) 的单位只有非零常数，单位点条件使该常数为一。后者的坐标像必须本原，前题已证全部本原多项式为 \(ct\)，反向加法性显然。
\end{solution}

\begin{optionalexercise}\label{b5:ex:21:12}
设 \(k\) 特征为零，给 \(\mathbb G_a\) 一个三维有理表示，微分为单个三阶幂零 Jordan 块 \(N\)。写出表示矩阵并求不动空间。
\end{optionalexercise}
\begin{solution}
以 \(N=\begin{psmallmatrix}0&1&0\\0&0&1\\0&0&0\end{psmallmatrix}\)，有
\[
\rho(t)=\begin{pmatrix}1&t&t^2/2\\0&1&t\\0&0&1\end{pmatrix}.
\]
一个向量被所有参数固定，当且仅当各个 \(t\) 次系数作用为零，等价于 \(Nv=0\)。因此不动空间是第一坐标轴，维数一。
\end{solution}

\begin{optionalexercise}\label{b5:ex:21:13}
在特征 \(p\) 中，比较平凡的一维 \(\mathbb G_m\) 表示与字符 \(t^p\) 的表示：其微分与有理群表示是否相同？
\end{optionalexercise}
\begin{solution}
平凡字符的微分为零，\(t^p\) 的导数也为零，所以两个 Lie 代数表示相同。但坐标 Laurent 多项式 \(1\) 与 \(t^p\) 不同，故群概形表示不同；在代数闭域的点上也可选 \(t^p\ne1\)。微分在这里不忠实。
\end{solution}

\begin{optionalexercise}\label{b5:ex:21:14}
设 \(k\) 代数闭且特征为 \(p\)。计算 \(\alpha_p(k)\)、\(\alpha_p\bigl(k[\epsilon]/(\epsilon^2)\bigr)\)，并与平凡群概形比较。
\end{optionalexercise}
\begin{solution}
域中 \(a^p=0\) 只可能 \(a=0\)，所以第一点群平凡。双数环中 \(a+b\epsilon\) 的 \(p\) 次幂为 \(a^p\)，故全部 \(b\epsilon\) 都给点，组成与 \((k,+)\) 同构的群。平凡群概形在任何代数上只有一个点，所以两者不同。
\end{solution}

\begin{optionalexercise}\label{b5:ex:21:15}
设 \(A=k[\operatorname{GL}_n]\)。证明约束 \(\det X=1\) 生成 Hopf 理想，并求其单位处线性化。
\end{optionalexercise}
\begin{solution}
\(\Delta(\det X-1)=(\det X-1)\otimes\det X+1\otimes(\det X-1)\)，余单位为零，对极将它送为 \(-(\det X)^{-1}(\det X-1)\)。故为Hopf理想。模 \(\epsilon^2\) 有 \(\det(I+\epsilon Y)=1+\epsilon\operatorname{tr}Y\)，线性化条件为迹零，在任意特征成立。
\end{solution}

\begin{optionalexercise}\label{b5:ex:21:16}
求 \(U_3\) 的中心和交换子子群，并说明它与加法群三重直积不同。
\end{optionalexercise}
\begin{solution}
按正文坐标，交换子为 \((0,0,ab'-a'b)\)。若与全部 \((a',b',c')\) 交换，必须 \(a=b=0\)，故中心是全部 \((0,0,c)\)。取 \(a=1,b'=c,a'=b=0\) 可实现每个中心元，所以交换子子群等于中心。该群非交换，而三个加法群直积交换。
\end{solution}

\begin{optionalexercise}\label{b5:ex:21:17}
对实型 \(b\) 的矩阵 \(H=\operatorname{diag}(1,-1)\)，求相似乘子的像，并给出乘子为 \(-1\) 的矩阵。
\end{optionalexercise}
\begin{solution}
交换坐标的 \(P=\begin{psmallmatrix}0&1\\1&0\end{psmallmatrix}\) 满足 \(P^{\mathsf T}HP=-H\)。正乘子由实标量 \(\sqrt c I\) 实现，负乘子由 \(\sqrt{-c}P\) 实现，故像为全部 \(\mathbb R^\times\)。正定型只能有正乘子，所以型的选择影响实际像。
\end{solution}

\begin{optionalexercise}\label{b5:ex:21:18}
设 \(b\) 为非退化型，\(a\in K^\times\)。证明标量矩阵 \(aI\) 总给相似变换，但其在自同态代数上的内自同构恒等。由此区分三个群。
\end{optionalexercise}
\begin{solution}
有 \((aI)^*H(aI)=\bar aaH\)，所以乘子为 \(\bar aa\)。标量与每个矩阵交换，故内自同构恒等。标量矩阵 \(aI\) 属于 \(\operatorname{Isom}(b)\) 当且仅当 \(\bar aa=1\)；相似群包含更多标量，而对合代数自同构群已经把这些标量全都商掉。
\end{solution}

\begin{optionalexercise}\label{b5:ex:21:19}
证明实点映射 \(\operatorname{SL}_2(\mathbb R)\to\operatorname{PGL}_2(\mathbb R)\) 的像恰为正行列式射影类，核为 \(\{\pm I\}\)。
\end{optionalexercise}
\begin{solution}
标量缩放把行列式乘实平方，所以像必有正号。若 \(\det A>0\)，则 \((\det A)^{-1/2}A\) 行列式一，给原射影类的原像。核是标量 \(aI\) 且 \(a^2=1\)，正为两个元。负行列式类没有原像。
\end{solution}

\begin{optionalexercise}\label{b5:ex:21:20}
在 \(\mathbb H\) 中令 \(q=(1+i+j+k)/2\)。计算共轭在基 \(i,j,k\) 上的矩阵，以及该矩阵的行列式。
\end{optionalexercise}
\begin{solution}
四元数范数为一，所以 \(q^{-1}=(1-i-j-k)/2\)。使用 \(ij=k\)、\(jk=i\)、\(ki=j\) 及反向乘积变号，逐项相乘得到
\(qiq^{-1}=j,qjq^{-1}=k,qkq^{-1}=i\)。
矩阵为 \(\begin{psmallmatrix}0&0&1\\1&0&0\\0&1&0\end{psmallmatrix}\)，是三轮换矩阵，行列式一。它固定 \(i+j+k\)，在其正交平面上是三分之一整周的旋转；\(q\) 与 \(-q\) 仍给相同作用。
\end{solution}

\begin{optionalexercise}\label{b5:ex:21:21}
对 \(n\ge2,m\ge0\)，证明 \(\operatorname{Sym}^m(\mathbb C^n)\) 的自然有理 \(\operatorname{SL}_{n,\mathbb C}\) 表示能下降为中心商 \(\operatorname{PGL}_{n,\mathbb C}\) 的有理表示，当且仅当 \(n\mid m\)；并在 \(n=2\) 时把条件写成字符格条件。
\end{optionalexercise}
\begin{solution}
由\cref{b5:prop:special-projective-quotient}，\(\operatorname{SL}_{n,\mathbb C}\to\operatorname{PGL}_{n,\mathbb C}\) 是核为 \(\mu_n\) 的中心商，有理表示下降当且仅当在该核上平凡。
标量 \(\zeta I\)、\(\zeta^n=1\)，在每个纯 \(m\) 次张量及其对称商上按 \(\zeta^m\) 作用。
复数域上的 \(\mu_n\) 约化，且有本原 \(n\) 次单位根，所以这个子群概形上的作用平凡当且仅当 \(n\mid m\)。
这同时保证了下降表示的正则性，而不只是复点上的群作用。
对 \(n=2\)，环面最高权为 \(m\lambda\)，下降要求它属于根格 \(2\mathbb Z\lambda\)，即 \(m\) 偶数。
\end{solution}

\begin{optionalexercise}\label{b5:ex:21:22}
取一个交换Hopf代数 \(A\)。证明每个 \(R\) 点的逆由 \(gS\) 给出时，还满足 \((gh)^{-1}=h^{-1}g^{-1}\)，并解释为何不能据此认为 \(G(R)\) 交换。
\end{optionalexercise}
\begin{solution}
Hopf对极给卷积逆，任意结合群中双边逆唯一，故 \((gh)^{-1}=h^{-1}g^{-1}\)。\(A\) 交换保证卷积结果仍是代数同态，却不保证其余乘法余交换。以 \(A=k[\operatorname{GL}_2]\) 为例，点群仍为通常的非交换矩阵群。
\end{solution}

% !TeX root = ../../Book5.tex
% 纲要标题：## 第22章*　不变量理论与商
% 纲要位置：工作记录/计划纲要.md，第 4429 行。
% 本轮补充的纲要细目：
% > 前四节以有限群为主，可在掌握本卷第1、2章与多项式基础后阅读；所需分次与有限性事实就地重述。后几节再使用第21章的代数群语言，不要求先通读全部 Lie 代数理论。

\OptionalChapter{不变量理论与商}{b5:ch:22}
\BookChapterNumber{b5:ch:22}

\begin{chapterabstract}
本章研究群作用下的不变环：证明有限群情形的有限生成与次数界，计算生成元和关系，再以矩阵配对及商的几何意义讨论连续群作用。
\end{chapterabstract}

\begin{learninggoals}
\begin{itemize}
\item 区分形式多项式、点集函数、轨道分离及轨道闭包。
\item 证明有限群不变环有限生成，利用次数界找到生成集，并用 Molien 公式计算各次维数。
\item 为具体不变环求出生成元和全部关系，保持正确的分次。
\item 把 Schur--Weyl 对偶转化为配对不变量定理，理解协变式与相对不变量。
\item 解释仿射商合并非闭轨道的方式，并辨认线性约化与几何约化的不同作用。
\end{itemize}
\end{learninggoals}

交换两个变量 \(x,y\) 不改变 \(s=x+y\) 与 \(p=xy\)。例如 \(x^2+y^2=s^2-2p\)，更高次数的对称多项式也可以逐次消去最高项，改写成 \(s,p\) 的多项式。这里有两件不同的事：找到显然不变的候选元，以及证明它们已经生成全部不变量。对一般群作用还要再问，候选元之间是否存在关系，关系是否已经找全。本章从有限群的轨道多项式和平均算子入手，再用消去与 Hilbert 级数处理关系，最后研究连续群作用下商的几何意义。

\ProseParagraphBreak
第22.1—22.4节从有限群作用出发，利用第1、2章及多项式环、有限生成模的基本性质，求不变环的候选生成元并核对其全部关系。第22.5节的配对定理使用\cref{b5:thm:schur-weyl}。第22.6、22.7节则使用第21章的坐标环与有理余模；其中有限群商、矩阵闭轨道和线性约化有限生成在此证明，一般正特征约化群的有关结论引用文献中的结果。

\ProseParagraphBreak
不变环的计算涉及三个不同的问题：确定作用与底域，证明候选元生成整个不变环，以及求出它们之间的全部关系。本章在这些计算的基础上讨论商的几何意义。关于极化、几何约化以及商的进一步背景，可参阅\cite[Sections 1.2, 3.4 and 6.1; Chapters 7--9]{Dolgachev2003InvariantTheory}。

% !TeX root = ../../Book5.tex
% 纲要标题：### 22.1 群作用与不变多项式
% 纲要位置：工作记录/计划纲要.md，第 4433 行。
% * 由 \(k[V]=\operatorname{Sym}(V^*)\) 定义多项式环上的作用，说明逆元、分次及有限域上形式多项式与点集函数的区别
% * 以代数闭域上的几何轨道表述有限群的分离结论；用明确算例开始，不将一般无限群的轨道结论混入
% * 群阶可逆时完整证明 Reynolds 投影及其不变子环线性性，说明它一般不保持乘法
% #### 22.1.1 线性作用、函数上的作用与不变子环
% #### 22.1.2 有限群的轨道与不变量分离
% #### 22.1.3 平均法与 Reynolds 算子
\section{群作用与不变多项式}
\label{b5:sec:2201}

矩阵的元素会随基改变，特征多项式却保持不变。一个多项式的根可以重新编号，根的初等对称式却仍是原来的系数。不变量理论研究哪些多项式在给定作用下保持不变，以及这些多项式能辨认多少对象。

\subsection{线性作用、函数上的作用与不变子环}
\label{b5:subsec:2201:01}

\begin{definition}\label{b5:def:invariant-ring-revisited}
设 \(k\) 为域，\(V\) 为有限维 \(k\) 向量空间，群 \(G\) 通过线性自同构作用在 \(V\) 上。在 \(k[V]=\operatorname{Sym}(V^*)\) 上定义
\[
(g\cdot f)(v)=f(g^{-1}v).
\]
此处等式由线性代入解释为形式多项式恒等式。满足 \(g\cdot f=f\) 对所有 \(g\in G\) 成立的多项式组成子代数
\[
k[V]^G=\bigl\{\,f\in k[V]\mid g\cdot f=f\text{ 对所有 }g\in G\,\bigr\},
\]
称为作用的\term[bubian-huan]{不变环}[invariant ring]，其中的元素称为多项式不变量。
\end{definition}

逆元保证 \(g\cdot(h\cdot f)=(gh)\cdot f\)。先在 \(V^*\) 上取对偶作用，再乘法延拓，就得到同一定义。它保持全次数；若 \(S_d=\operatorname{Sym}^d(V^*)\)，则 \(k[V]^G=\bigoplus_{d\ge0}S_d^G\)。有限域上的形式多项式不能等同于点集函数，例如非零多项式 \(x^q-x\) 在 \(\mathbb F_q\) 每一点取零。因此本章的不变性始终指代数中的等式。

\begin{example}\label{b5:exm:invariant-sign-action}
设 \(\operatorname{char}k\ne2\)，阶二群由 \(s\) 生成，在 \(k^2\) 上同时变号。则 \(s\cdot x=-x,\ s\cdot y=-y\)，单项式 \(x^ay^b\) 不变当且仅当 \(a+b\) 为偶数。若两指数均偶，它是 \(x^2,y^2\) 的乘积；若均奇，先提出 \(xy\)。故
\[
k[x,y]^{\langle s\rangle}=k[x^2,xy,y^2].
\]
生成元满足 \((xy)^2=x^2y^2\)。找到生成元之后，还须求出它们之间的全部关系。
\end{example}

\subsection{有限群的轨道与不变量分离}
\label{b5:subsec:2201:02}

同一轨道上的不变量值必然相同。有限群情形的逆命题不需要除以群阶，但要在几何点上陈述。

\begin{proposition}\label{b5:prop:finite-invariants-separate}
设 \(k\) 为代数闭域，有限群 \(G\) 通过 \(k\) 代数自同构作用在有限型仿射 \(k\) 概形 \(X=\operatorname{Spec}A\) 上。若 \(x,y\in X(k)\) 不在同一轨道，则存在 \(F\in A^G\) 满足 \(F(x)=0,\ F(y)=1\)。
\end{proposition}
\begin{proof}
两个轨道是互不相交的有限点集。不同点在 \(A\) 中的极大理想两两互素，中国剩余定理给出 \(f\in A\)，使它在 \(Gx\) 上取零，在 \(Gy\) 上取一。令 \(F=\prod_{g\in G}g\cdot f\)。群作用只置换乘积因子，所以 \(F\) 不变；在 \(x\) 处每个因子为零，在 \(y\) 处每个因子为一。
\end{proof}

一般域上的几何轨道要扩到代数闭包后考察。一般无限群也没有上述结论：\(k^\times\) 在直线上按数乘作用时，零点与非零点属于不同轨道，却只有常数不变量。

\subsection{平均法与 Reynolds 算子}
\label{b5:subsec:2201:03}

\begin{proposition}\label{b5:prop:invariant-reynolds}
设 \(k\) 为域，有限群 \(G\) 通过 \(k\) 代数自同构作用在交换含幺 \(k\) 代数 \(A\) 上，\(A^G\) 为其不变子代数，且 \(|G|\) 在 \(k\) 中可逆。则
\[
\mathcal R:A\longrightarrow A^G,\qquad
\mathcal R(f)=\frac1{|G|}\sum_{g\in G}g\cdot f
\]
是 \(A^G\) 线性投影，即 \(\mathcal R|_{A^G}=1\)，且
\(\mathcal R(af)=a\mathcal R(f)\) 对 \(a\in A^G\) 成立。若作用保持分次，则 \(\mathcal R\) 保持次数。此映射称为\term[Reynolds-suanzi]{Reynolds 算子}[Reynolds operator]。
\end{proposition}
\begin{proof}
作用于求和式的群元素只重排指标，所以像落在 \(A^G\)。不变元的各求和项相同，平均后仍为原元。对不变的 \(a\)，每项 \(g\cdot(af)=a(g\cdot f)\)，提出 \(a\) 即得线性性。保次数则由各求和项直接得到。
\end{proof}

这个投影通常不保持乘法。变号作用下 \(\mathcal R(x)=0\)，但 \(\mathcal R(x^2)=x^2\)。其主要用处是处理系数：若 \(f=\sum_i a_if_i\)，而 \(f,f_i\) 都不变，则
\[
f=\sum_i\mathcal R(a_i)f_i.
\]
这使一个带任意系数的等式转化为带不变系数的等式。

% !TeX root = ../../Book5.tex
% 纲要标题：### 22.2 有限生成与次数界
% 纲要位置：工作记录/计划纲要.md，第 4445 行。
% * 逐步证明轨道多项式给出的整性、有限模性及不变子环的有限生成，并就地建立 Artin–Tate 引理
% * 有限群不变量的有限生成不要求群阶可逆，不能把平均法失效误作有限生成失效
% * 特征零下完整证明可取次数不超过群阶的齐次生成元，并说明有限搜索如何终止；更广特征的次数界另列证明范围
% * 模特征情形给具体比较；线性约化群的一般有限生成归22.7，不参与本节有限群证明
% #### 22.2.1 轨道多项式与整性
% #### 22.2.2 有限模性及 Artin–Tate 引理
% #### 22.2.3 有限群不变量的有限生成
% #### 22.2.4 Noether 次数界与有效计算
\section{有限生成与次数界}
\label{b5:sec:2202}

有限生成代数的子环未必有限生成。有限群作用提供了额外的整性条件；特征零时，平均法与极化还给出可以停止计算的明确次数。

\subsection{轨道多项式与整性}
\label{b5:subsec:2202:01}

\begin{proposition}\label{b5:prop:orbit-polynomial-integrality}
设 \(k\) 为域，有限群 \(G\) 通过 \(k\) 代数自同构作用在交换含幺 \(k\) 代数 \(A\) 上，\(A^G\) 为其不变子代数。则每个 \(a\in A\) 都在 \(A^G\) 上整。若 \(A\) 是有限生成 \(k\) 代数，则 \(A\) 是有限生成 \(A^G\) 模。
\end{proposition}
\begin{proof}
多项式
\[
P_a(T)=\prod_{g\in G}(T-g\cdot a)
\]
首一，其系数因群作用只置换因子而不变，且 \(P_a(a)=0\)。若 \(A=k[a_1,\ldots,a_n]\)，则也有 \(A=A^G[a_1,\ldots,a_n]\)。设 \(a_i\) 满足次数 \(d_i\) 的首一整性方程。反复用它降低指数，得到有限模生成集
\[
\{\,a_1^{e_1}\cdots a_n^{e_n}\mid0\le e_i<d_i\,\}.
\]
\end{proof}

轨道多项式允许重复因子，特征也可以整除群阶；这里没有使用平均。

\subsection{有限模性及 Artin--Tate 引理}
\label{b5:subsec:2202:02}

\begin{lemma}[Artin--Tate]\label{b5:lem:invariants-artin-tate}
设 \(k\) 为域，\(k\subseteq B\subseteq A\) 为交换含单位代数。若 \(A\) 是有限生成 \(k\) 代数，又是有限生成 \(B\) 模，则 \(B\) 是有限生成 \(k\) 代数。
\end{lemma}
\begin{proof}
取 \(A=k[a_1,\ldots,a_n]\)，并取 \(B\) 模生成元 \(e_1,\ldots,e_r\)，增补后可设 \(e_1=1\)。将各 \(a_i\) 与各 \(e_je_\ell\) 写成 \(e_j\) 的 \(B\) 线性组合，由这些有限多个系数生成的 \(k\) 子代数记为 \(B_0\)。模 \(M=\sum_j B_0e_j\) 包含一和全部 \(a_i\)，又因已收齐乘法结构系数而对乘法封闭。因此 \(A\subseteq M\subseteq A\)，即 \(A\) 是有限 \(B_0\) 模。

\(B_0\) 为有限生成 \(k\) 代数，由 Hilbert 基定理为 Noether 环。因此有限 \(B_0\) 模 \(A\) 的子模 \(B\) 有有限生成元 \(b_1,\ldots,b_s\)，从而 \(B=B_0[b_1,\ldots,b_s]\)，这是有限生成 \(k\) 代数。
\end{proof}

不能直接对尚不知为 Noether 环的 \(B\) 使用子模有限生成；引入 \(B_0\) 正是为了避免循环。

\subsection{有限群不变量的有限生成}
\label{b5:subsec:2202:03}

\begin{theorem}\label{b5:thm:finite-invariant-generation}
设 \(k\) 为任意域，有限群 \(G\) 通过 \(k\) 代数自同构作用在有限生成交换含幺 \(k\) 代数 \(A\) 上，\(A^G\) 为其不变子代数。则 \(A^G\) 是有限生成 \(k\) 代数，且 \(A\) 是有限生成 \(A^G\) 模。若 \(A=k[V]=\operatorname{Sym}_k(V^*)\) 来自有限维 \(k\) 线性表示 \(V\)，可以取齐次不变量作为代数生成元。
\end{theorem}
\begin{proof}
\cref{b5:prop:orbit-polynomial-integrality}给出有限模性，再对 \(B=A^G\) 使用\cref{b5:lem:invariants-artin-tate}。在线性作用下，各齐次分量分别不变，将已有有限生成元全部拆成齐次分量，所得有限集仍生成整个不变环。
\end{proof}

模特征中失效的是平均法，而非此定理。例如在特征 \(p>0\) 中，令阶 \(p\) 群在 \(k[x,y]\) 上由 \(\sigma(x)=x+y,\ \sigma(y)=y\) 生成。元素 \(z=x^p-xy^{p-1}\) 不变，而且
\begin{equation}\label{b5:eq:modular-transvection-invariants}
k[x,y]^{\langle\sigma\rangle}=k[y,z].
\end{equation}
事实上，首一关系 \(x^p=z+xy^{p-1}\) 使每个多项式唯一写为
\(\sum_{i=0}^{p-1}a_i(y,z)x^i\)。唯一性来自各项最高 \(x\) 次数模 \(p\) 不同。若某个正指标的 \(a_i\ne0\)，取最大的 \(i\)；在 \(\sigma(f)-f\) 中，基向量 \(x^{i-1}\) 的系数是 \(iya_i\ne0\)，矛盾。于是只剩 \(a_0\)。

\subsection{Noether 次数界与有效计算}
\label{b5:subsec:2202:04}

\begin{theorem}[特征零下的 Noether 次数界]\label{b5:thm:noether-degree-bound}
设 \(k\) 为特征零域，有限群 \(G\) 线性作用在有限维 \(k\) 向量空间 \(V\) 上，\(m=|G|\)，\(k[V]=\operatorname{Sym}_k(V^*)\) 为多项式函数代数，各线性函数的次数为一。则不变环 \(k[V]^G\) 由全次数不超过 \(m\) 的齐次不变量生成。
\end{theorem}
\begin{proof}
令 \(S=k[V]\)，并令 \(J\subseteq S\) 为次数在 \(1,\ldots,m\) 内的全部齐次不变量生成的理想。对任意线性形式 \(\ell\in V^*\)，有
\[
\prod_{g\in G}(T-g\cdot\ell)
=T^m+c_1T^{m-1}+\cdots+c_m,\qquad c_i\in S_i^G.
\]
代入 \(T=\ell\)，得到 \(\ell^m\in J\)。

对任意 \(\ell_1,\ldots,\ell_m\in V^*\)，容斥展开给出极化等式
\[
m!\,\ell_1\cdots\ell_m
=\sum_{I\subseteq\{1,\ldots,m\}}(-1)^{m-|I|}
\left(\sum_{i\in I}\ell_i\right)^m.
\]
遗漏某个指标的项相消，所有指标各出现一次的项共有 \(m!\) 个。特征零使 \(m!\) 可逆，因此 \(S_m\subseteq J\)，进而 \(S_d\subseteq J\) 对 \(d\ge m\) 成立。

在每个 \(S_i^G\)（\(1\le i\le m\)）取基，合并记为 \(f_1,\ldots,f_r\)，其生成的代数记为 \(B\)。对次数归纳。次数不超过 \(m\) 的不变量已在 \(B\) 中。若 \(f\in S_d^G,\ d>m\)，可写
\[
f=\sum_j a_jf_j,\qquad \deg a_j=d-\deg f_j<d,
\]
其中取齐次系数即可。应用\cref{b5:prop:invariant-reynolds}得
\(f=\sum_j\mathcal R(a_j)f_j\)。各系数是不变元且次数下降，由归纳假设属于 \(B\)，故 \(f\in B\)。最后拆分齐次分量便得到全部结论。
\end{proof}

因而只需在次数 \(1,\ldots,|G|\) 中列出单项式，计算其平均，并用线性代数选基，就能找到有限生成集。它通常冗余，且单项式数目可能很大；次数界保证终止，不保证快速。

\ProseParagraphBreak
上述证明也适用于特征大于 \(|G|\) 的域。群阶可逆而特征不大于群阶时，极化所需的阶乘仍可能为零，因而仅由本证明不能得到同一次数界。模特征则应另用具体构造，例如\eqref{b5:eq:modular-transvection-invariants}。

% !TeX root = ../../Book5.tex
% 纲要标题：### 22.3 Hilbert 级数与 Molien 公式
% 纲要位置：工作记录/计划纲要.md，第 4460 行。
% * 在复数域或明确的特征零范围内，由各次对称幂上的平均迹完整推导 Molien 公式
% * 计算循环群及小型反射群的 Hilbert 级数，解释各次系数而不只给出有理函数
% * 区分 Hilbert 级数、生成元与关系；正特征的域内迹不能无条件当作整数维数
% #### 22.3.1 齐次不变量的维数与特征标
% #### 22.3.2 Molien 公式
% #### 22.3.3 循环群与小型反射群的计算
\section{Hilbert 级数与 Molien 公式}
\label{b5:sec:2203}

有限生成告诉我们最终可以找到一张有限的生成元表，但没有立即给出每个次数有多少独立不变量。把这个维数问题转成迹的计算，就能直接使用有限群特征标。本节在复数域上讨论，以便将迹与整数维数明确对应。

\subsection{齐次不变量的维数与特征标}
\label{b5:subsec:2203:01}

\begin{definition}\label{b5:def:invariant-hilbert-series}
设 \(k\) 为域，\(A=\bigoplus_{d\ge0}A_d\) 为非负分次 \(k\) 代数，每个 \(A_d\) 都是有限维 \(k\) 向量空间。形式幂级数
\[
H_A(t)=\sum_{d\ge0}(\dim_k A_d)t^d\in\mathbb Z[\![t]\!]
\]
称为 \(A\) 的\term[Hilbert-jishu]{Hilbert 级数}[Hilbert series]。其分次属于数据的一部分；生成元的次数不是一时，不能擅自按次数一计算。
\end{definition}

例如次数为 \(a_1,\ldots,a_r>0\) 的多项式环有级数
\[
\prod_{i=1}^r(1-t^{a_i})^{-1},
\]
因为展开各几何级数就是按单项式的次数计数。这里不要求 \(t\) 是复数，也没有收敛问题。

\begin{lemma}\label{b5:lem:invariant-dimension-trace}
设有限群 \(G\) 作用在有限维复空间 \(W\) 上。则
\[
\dim W^G=\frac1{|G|}\sum_{g\in G}\operatorname{tr}(g\mid W).
\]
\end{lemma}
\begin{proof}
平均算子是像为 \(W^G\) 的幂等算子。按像与核分解 \(W\)，其矩阵为 \(\operatorname{diag}(I,0)\)，故迹等于像的维数；迹的线性性给出右端。
\end{proof}

应用于 \(W=\operatorname{Sym}^d(V^*)\)，便把第 \(d\) 次不变量的数目写成对称幂特征标的平均。模特征中，即使偶然可以构造别的投影，域内的迹也至多记录维数在该域中的像，不能直接恢复整数维数。

\subsection{Molien 公式}
\label{b5:subsec:2203:02}

\begin{theorem}[Molien 公式]\label{b5:thm:molien}
设有限群 \(G\) 线性作用在有限维复向量空间 \(V\) 上，\(g_V,g_{V^*}\) 分别表示 \(g\in G\) 在 \(V\) 及对偶空间上的算子，\(I\) 为相应恒等算子。记 \(H_{\mathbb C[V]^G}(t)\) 为不变环按多项式总次数分次的 Hilbert 级数，则
\[
H_{\mathbb C[V]^G}(t)
=\frac1{|G|}\sum_{g\in G}\frac1{\det(I-tg_{V^*})}
=\frac1{|G|}\sum_{g\in G}\frac1{\det(I-tg_V)}.
\]
两端均按 \(t=0\) 展开为形式幂级数。
\end{theorem}
\begin{proof}
每个 \(g\) 有有限阶，在复数域上可对角化。设其在 \(V^*\) 上的特征值为 \(\lambda_1,\ldots,\lambda_n\)。在特征向量的单项式基下，
\[
\sum_{d\ge0}\operatorname{tr}\bigl(g\mid\operatorname{Sym}^d(V^*)\bigr)t^d
=\sum_{a_1,\ldots,a_n\ge0}\lambda_1^{a_1}\cdots\lambda_n^{a_n}t^{a_1+\cdots+a_n}
=\prod_i(1-\lambda_it)^{-1}.
\]
现在逐次使用\cref{b5:lem:invariant-dimension-trace}，有限求和可以与形式级数相交换，得到第一式。对偶作用的矩阵为 \(g_V^{-T}\)，而 \(g\mapsto g^{-1}\) 置换群元素，所以求和后可把它换成 \(g_V\)，得到第二式。两式相等不表示两个对应单项对每个 \(g\) 都相等。
\end{proof}

计算时可按共轭类求和，因为同一共轭类的特征多项式相同。级数的系数虽写成单位根表达式的平均，最终必为非负整数，正是各次不变子空间的维数。

\subsection{循环群与小型反射群的计算}
\label{b5:subsec:2203:03}

令 \(C_m\) 在一维复空间上由本原 \(m\) 次单位根 \(\zeta\) 数乘。只有 \(x^{jm}\) 不变，故
\[
\frac1m\sum_{r=0}^{m-1}\frac1{1-\zeta^rt}
=\frac1{1-t^m}.
\]
直接展开左端也能看到这一点：\(t^d\) 的系数是 \(\dfrac1m\sum_r\zeta^{rd}\)，它在 \(m\mid d\) 时为一，其余为零。

\begin{example}\label{b5:exm:molien-sign-plane}
对同时变号的二维作用，单位元与非平凡元分别贡献 \((1-t)^{-2}\)、\((1+t)^{-2}\)，所以
\[
H(t)=\frac12\left(\frac1{(1-t)^2}+\frac1{(1+t)^2}\right)
=\frac{1+t^2}{(1-t^2)^2}
=\frac{1-t^4}{(1-t^2)^3}.
\]
因此次数 \(2r\) 的维数为 \(2r+1\)，奇数次维数为零。最后一个表达式与三个二次生成元、一个四次关系相容；仅由级数本身，还不能断定存在这样的展示。
\end{example}

再让 \(S_3\) 置换 \(\mathbb C^3\) 的坐标。单位元、三个换位、两个三轮换分别给出
\[
H(t)=\frac16\left(
\frac1{(1-t)^3}
+\frac3{(1-t)^2(1+t)}
+\frac2{1-t^3}
\right)
=\frac1{(1-t)(1-t^2)(1-t^3)}.
\]
这与初等对称多项式的次数 \(1,2,3\) 一致。在反射表示 \(x_1+x_2+x_3=0\) 上，平凡的一维方向被去掉；由于特征为零，分解为该平面与平凡直线是直和，故其级数为
\[
\frac1{(1-t^2)(1-t^3)}.
\]
例如二次与三次不变量空间各一维，六次则有两个独立候选 \(e_2^3,e_3^2\)。下一节说明如何把这种计数用于检验生成元及其关系。

% !TeX root = ../../Book5.tex
% 纲要标题：### 22.4 生成元、关系与 Gröbner 计算
% 纲要位置：工作记录/计划纲要.md，第 4472 行。
% * 在已证次数界内通过平均和线性代数搜索齐次生成元，证明生成性后再处理关系
% * 用消去计算核；未证明候选元生成整个不变环时，所得关系只描述其生成的子代数
% * Hilbert 级数核验必须先有保持次数的满射，不由两个级数相等直接断言环同构
% * 完整处理特征不为二时同时变号的不变环 \(k[x^2,xy,y^2]\cong k[u,v,w]/(uw-v^2)\)，保持生成元权二分次
% #### 22.4.1 齐次生成元的搜索与完备性
% #### 22.4.2 用消去求生成元之间的关系
% #### 22.4.3 用 Hilbert 级数核验展示
% #### 22.4.4 循环作用的完整算例
\section{生成元、关系与 Gröbner 计算}
\label{b5:sec:2204}

一个候选不变量很容易检验；困难在于证明没有遗漏其他生成元，以及没有遗漏生成元之间的关系。这两种完备性需要不同的理由。次数界解决第一项，消去和分次维数可以解决第二项。

\subsection{齐次生成元的搜索与完备性}
\label{b5:subsec:2204:01}

设 \(G\) 已由有限个矩阵给出，并已列出其全部元素，在可有效进行四则运算的特征零域 \(k\) 上作用于 \(S=k[x_1,\ldots,x_n]\)。以下是一种有限的计算方法。
\begin{algorithm}[htbp]
\caption{Reynolds 平均生成元搜索}
\label{b5:alg:reynolds-generator-search}
\begin{algorithmsteps}
{特征零的精确可计算域 \(k\)、有限群 \(G\) 的全部作用矩阵，以及 \(S=k[x_1,\ldots,x_n]\)。}
{\(S^G\) 的有限齐次代数生成集。}
\State 对每个 \(1\le d\le |G|\)，列出全部次数为 \(d\) 的单项式，计算其 Reynolds 平均。
\State 把所得多项式按单项式坐标写成列向量，经消元取出 \(S_d^G\) 的一组基。
\State 逐次处理各 \(d\)。把先前保留的生成元的所有次数为 \(d\) 的乘积展开，取其线性生成空间；从所取的基中保留补足该空间所需的元素。
\end{algorithmsteps}
\end{algorithm}
步骤一、二求得整个 \(S_d^G\)，因为投影把一个生成集送到像的生成集。步骤三仅删除已经属于先前生成子代数的元，不改变最终生成的代数。\cref{b5:thm:noether-degree-bound}保证在次数 \(|G|\) 后停止时已经生成 \(S^G\)。零平均必须删去，但不能只因两个候选次数相同就认为它们冗余。

\ProseParagraphBreak
此法以已知的有限群及精确系数运算为输入，输出一个有限齐次生成集；它不是关于任意给定矩阵群是否有限的判定算法。即使群很小，单项式空间也可能很大，故这里的终止性不应读成复杂度保证。

\subsection{用消去求生成元之间的关系}
\label{b5:subsec:2204:02}

\begin{proposition}\label{b5:prop:invariant-elimination-presentation}
设 \(k\) 为域，\(f_1,\ldots,f_r\in k[x_1,\ldots,x_n]\)，定义
\[
\phi:k[u_1,\ldots,u_r]\longrightarrow k[x_1,\ldots,x_n],
\qquad u_i\longmapsto f_i.
\]
在大多项式环 \(k[x_1,\ldots,x_n,u_1,\ldots,u_r]\) 中令
\[
J=(u_1-f_1,\ldots,u_r-f_r).
\]
则 \(\ker\phi=J\cap k[u_1,\ldots,u_r]\)。若 \(\mathcal G\) 是 \(J\) 在消去 \(x\) 变量的单项式序下的 Gröbner 基，则
\(\mathcal G\cap k[u_1,\ldots,u_r]\) 是这个核的 Gröbner 基。
\end{proposition}
\begin{proof}
代入 \(u_i=f_i(x)\) 定义从大环到 \(k[x]\) 的满射，其核为 \(J\)：按各 \(u_i-f_i\) 逐次除法，任意多项式模 \(J\) 化为其代入所得的多项式。把这张映射限制在 \(k[u]\) 就是 \(\phi\)，所以核为所述交。

消去序的条件是每个含 \(x\) 的单项式都大于所有只含 \(u\) 的单项式。若 \(h\in J\cap k[u]\) 非零，其首项被某个 \(g\in\mathcal G\) 的首项整除。于是 \(g\) 的首项只含 \(u\)，消去序又迫使 \(g\) 的所有项只含 \(u\)。故交集中的基元已经控制核的全部首项，这正是 Gröbner 基的定义。
\end{proof}

因而求关系的步骤是：引入独立变量 \(u_i\)，生成 \(J\)，在消去序下计算 Gröbner 基，然后保留不含 \(x\) 的基元。若 \(f_i\) 已证明生成 \(k[V]^G\)，得到的商环才是不变环；否则它只展示子代数 \(k[f_1,\ldots,f_r]\)。构造 \(J\) 本身不会证明候选生成元充分。

\subsection{用 Hilbert 级数核验展示}
\label{b5:subsec:2204:03}

\begin{lemma}\label{b5:lem:hilbert-surjection-test}
设 \(k\) 为域，\(A,B\) 为每个齐次分量都有限维的非负分次 \(k\) 代数，\(\psi:A\to B\) 为保持次数的满射，\(H_A,H_B\) 为相应 Hilbert 级数。若 \(H_A(t)=H_B(t)\)，则 \(\psi\) 是同构。
\end{lemma}
\begin{proof}
逐次取核得到
\[
0\longrightarrow(\ker\psi)_d\longrightarrow A_d\longrightarrow B_d\longrightarrow0.
\]
维数相等使每个核分量为零，而分次核是这些分量的直和，故整个核为零。
\end{proof}

例如 \(u_i\) 的次数为 \(a_i>0\)，齐次非零多项式 \(F\in k[u_1,\ldots,u_r]\) 的次数为 \(b\)。多项式环为整环，乘 \(F\) 的映射单射，所以
\[
H_{k[u]/(F)}(t)=\frac{1-t^b}{\prod_i(1-t^{a_i})}.
\]
这来自次数平移的短正合列，而不是从“一个关系”凭直觉减去一项。多个关系只有在另行建立正则序列等条件后，才可逐项作同样的乘法。

\subsection{循环作用的完整算例}
\label{b5:subsec:2204:04}

\begin{proposition}\label{b5:prop:sign-invariant-presentation}
设 \(k\) 为特征不为二的域，二阶循环群 \(C_2\) 的非单位元在 \(k[x,y]\) 上把 \(x,y\) 分别变为 \(-x,-y\)。则存在分次同构
\[
k[u,v,w]/(uw-v^2)\ \simeq\ k[x,y]^{C_2},
\qquad
(u,v,w)\longmapsto(x^2,xy,y^2),
\]
其中 \(x,y\) 的次数为一，\(u,v,w\) 的次数均为二。
\end{proposition}
\begin{proof}
\cref{b5:exm:invariant-sign-action}已按单项式的指数奇偶证明满射。关系 \(uw-v^2\) 显然落入核。模此关系，每个多项式都可化为
\[
A(u,w)+v\cdot B(u,w).
\]
其像的第一部分只含 \(x,y\) 指数均偶的单项式，第二部分只含指数均奇的单项式。这两组互不相交，各组内不同 \(u^aw^b\) 也给出不同单项式，故若像为零，所有系数为零。于是没有其他关系。
\end{proof}

在消去计算中，同一个核由首项为 \(v^2\) 的单元素 Gröbner 基 \(\{v^2-uw\}\) 控制，标准单项式正是上述两组。它们给出
\[
H(t)=\frac{1+t^2}{(1-t^2)^2}.
\]
在复数域上，这也与\cref{b5:exm:molien-sign-plane}相符。两种核验使用的资料不同：单项式法直接证明线性无关，Molien 法则在已有满射后比较维数。

\ProseParagraphBreak
这个例子还能推广。设 \(k\) 含本原 \(m\) 次单位根，且特征不整除 \(m\)，令生成元在坐标环上作用为
\[
x\longmapsto\zeta x,\qquad y\longmapsto\zeta^{-1}y.
\]
不变条件是 \(m\mid(a-b)\)。从 \(x^ay^b\) 先提出 \((xy)^{\min(a,b)}\)，余下是 \(x^m\) 或 \(y^m\) 的幂，因而
\begin{equation}\label{b5:eq:cyclic-surface-presentation}
k[x,y]^{C_m}\simeq k[u,v,w]/(uv-w^m),
\quad
(u,v,w)=(x^m,y^m,xy).
\end{equation}
为验证全部关系，将含 \(uv\) 的项用 \(w^m\) 消去，留下
\[
w^c,\qquad u^aw^c\ (a\ge1),\qquad v^bw^c\ (b\ge1).
\]
代入后的指数差分别为零、正的 \(m\) 倍、负的 \(m\) 倍，且每类内指数唯一决定原单项式，所以这些像线性无关。分次为 \(\deg u=\deg v=m,\deg w=2\)，于是级数为
\[
\frac{1-t^{2m}}{(1-t^m)^2(1-t^2)}.
\]
当 \(m=2\) 时，只是前三个生成元的名称与上例有所调换。

% !TeX root = ../../Book5.tex
% 纲要标题：### 22.5 经典不变量与协变式
% 纲要位置：工作记录/计划纲要.md，第 4487 行。
% * 就地陈述并解释对称多项式基本定理；一般有限反射群定理标明选读及外部证明范围
% * 从单矩阵特征多项式及特征零下的迹出发；同特征多项式不等于矩阵共轭，多个矩阵的迹字理论另列范围
% * 完整处理一般线性群的配对不变量基本定理及其实际依赖，选一个正交或辛群的 Gram 型例子
% * 以低次二元型区分不变量、相对不变量与协变式，明确一般线性群和特殊线性群下判别式的不同变换性质
% #### 22.5.1 对称多项式与反射群的例子
% #### 22.5.2 矩阵共轭下的特征多项式与迹
% #### 22.5.3 向量与对偶向量的配对不变量
% #### 22.5.4 二元型的判别式与协变式
\section{经典不变量与协变式}
\label{b5:sec:2205}

不同的群产生不同的不变量问题。置换坐标保留对称多项式，改变矩阵的基保留特征多项式，同时改变向量与对偶向量则保留配对。本节从这些具体作用中证明几类基本结果，并说明不变量之外还有相对不变量和协变式。

\subsection{对称多项式与反射群的例子}
\label{b5:subsec:2205:01}

这一基本定理的主要证明已在第一卷定理12.4.4给出。为使本节的对称群与矩阵例子能够独立阅读，下面回顾陈述及首项消去证明，并保留它在正特征中的适用范围。

\begin{theorem}[对称多项式基本定理]\label{b5:thm:invariant-symmetric-polynomials}
设 \(k\) 为域，\(S_n\) 置换 \(k[x_1,\ldots,x_n]\) 的变量。令 \(e_i\) 为 \(i\) 次初等对称多项式，则
\[
k[x_1,\ldots,x_n]^{S_n}=k[e_1,\ldots,e_n],
\]
且 \(e_1,\ldots,e_n\) 代数独立。此结论不要求特征零。
\end{theorem}
\begin{proof}
固定字典序 \(x_1>\cdots>x_n\)，先处理齐次对称多项式 \(f\)。其首单项式为 \(x_1^{a_1}\cdots x_n^{a_n}\)，必有 \(a_1\ge\cdots\ge a_n\)：若 \(a_i<a_j\) 且 \(i<j\)，互换这两个变量便出现更大的同系数单项式，矛盾。多项式
\[
e_1^{a_1-a_2}e_2^{a_2-a_3}\cdots e_{n-1}^{a_{n-1}-a_n}e_n^{a_n}
\]
恰有该首单项式，首系数为一。乘以 \(f\) 的首系数后相减，首项严格下降。全次数不变，而同次数只有有限多个单项式，所以有限步后余式为零。拆分齐次分量就证明生成性。

另一方面，不同的 \(e_1^{b_1}\cdots e_n^{b_n}\) 的首指数是
\[
(b_1+\cdots+b_n,\ b_2+\cdots+b_n,\ldots,b_n),
\]
它唯一确定 \((b_1,\ldots,b_n)\)。非零多项式关系中的最大首项因而不可能相消，所以不存在这样的关系。
\end{proof}

还可令符号置换群 \((C_2)^n\rtimes S_n\) 作用在特征不为二的 \(k^n\) 上。先对各坐标独立变号取不变量，得到 \(k[x_1^2,\ldots,x_n^2]\)；再对置换取不变量，得到
\[
k\bigl[e_1(x_1^2,\ldots,x_n^2),\ldots,e_n(x_1^2,\ldots,x_n^2)\bigr].
\]
其基本次数为 \(2,4,\ldots,2n\)。这里两步都给出了生成性的证明。这是反射群具有多项式不变环的典型例子；一般复有限反射群的多项式性及其逆命题属于 Chevalley--Shephard--Todd 定理\footnote{谢泼德（Geoffrey Colin Shephard，1927-2016）和托德（John Arthur Todd，1908-1994）都是英国数学家，研究几何及有限复反射群；舍瓦莱已在本卷前文介绍。}，不能仅由这两个例子推广得出。

\subsection{矩阵共轭下的特征多项式与迹}
\label{b5:subsec:2205:02}

\begin{theorem}\label{b5:thm:single-matrix-invariants}
设 \(k\) 为代数闭域，\(n\ge1\) 为整数，\(\operatorname{GL}_n(k)\) 在 \(M_n(k)\) 上按共轭作用，\(k[M_n]\) 表示矩阵空间的多项式函数代数。对 \(A\in M_n(k)\)，定义特征多项式的系数函数 \(c_i\) 为
\[
\det(TI-A)=T^n-c_1(A)T^{n-1}+\cdots+(-1)^n\cdot c_n(A).
\]
则 \(k[M_n]^{\operatorname{GL}_n}=k[c_1,\ldots,c_n]\)，且这些生成元代数独立。若 \(\operatorname{char}k=0\)，也可用
\(\operatorname{tr}(A),\ldots,\operatorname{tr}(A^n)\) 生成。
\end{theorem}
\begin{proof}
把不变量 \(f\) 限制到对角矩阵。置换矩阵共轭说明这个限制是对称多项式，故由\cref{b5:thm:invariant-symmetric-polynomials}等于某个 \(P(e_1,\ldots,e_n)\)。于是 \(f-P(c_1,\ldots,c_n)\) 在对角矩阵及其全部共轭矩阵上为零。

特征多项式有互异根的矩阵构成非空 Zariski 开集：其补集是判别式多项式的零集，而取互异对角元证明该多项式非零。仿射空间的非空开集稠密，这个开集中的矩阵都可对角化。因此差多项式在稠密集上为零，只能恒为零，得到生成性。任何 \(c_i\) 的关系限制到对角矩阵就是 \(e_i\) 的关系，故代数独立。

令 \(p_j(A)=\operatorname{tr}(A^j)\)，并置 \(c_0=1\)。在对角矩阵上使用 Newton 恒等式\footnote{牛顿（Isaac Newton，1642/1643-1727），英国数学家、物理学家，奠定经典力学并独立发展微积分。}，再由稠密性延拓，得到
\[
rc_r=\sum_{j=1}^r(-1)^{j-1}c_{r-j}p_j
\qquad(1\le r\le n).
\]
特征零允许逐次除以 \(r\)，故每个 \(c_r\) 都在前 \(r\) 个迹生成的代数中。反向也由这些恒等式递推得到。
\end{proof}

特征多项式相同并不意味共轭。例如零矩阵与非零的二阶幂零 Jordan 块有相同特征多项式，因而所有多项式不变量都相同，却有不同的秩。多个矩阵的同时共轭涉及 \(\operatorname{tr}(A_{i_1}\cdots A_{i_r})\) 等迹字；上面的对角化证明不能用于一组未必交换的矩阵，本节只处理单矩阵。

\subsection{向量与对偶向量的配对不变量}
\label{b5:subsec:2205:03}

\begin{theorem}[一般线性群的配对不变量]\label{b5:thm:gl-pairing-fft}
设 \(k\) 为特征零域，\(V\) 为有限维非零 \(k\) 向量空间，\(r,s\ge0\) 为整数。群 \(\operatorname{GL}(V)\) 在
\[
W=V^{\oplus r}\oplus(V^*)^{\oplus s}
\]
上同时按原作用与对偶作用。则 \(k[W]^{\operatorname{GL}(V)}\) 由配对
\[
p_{ji}(v_1,\ldots,v_r,\phi_1,\ldots,\phi_s)=\phi_j(v_i)
\]
生成。
\end{theorem}
\begin{proof}
先处理对每个向量、对偶向量变量都一次齐次的多线性不变量，共有 \(a\) 个向量变量和 \(b\) 个对偶向量变量。标量矩阵 \(tI\) 使函数值乘以 \(t^{a-b}\)。域 \(k\) 无限，若 \(a\ne b\)，选 \(t\) 使此因子不等于一，就知道该函数为零。因此只需 \(a=b=d\)。

这时一个多线性函数 \(F\) 唯一对应算子 \(A\in\operatorname{End}_k(V^{\otimes d})\)，使
\[
F(v_1,\ldots,v_d,\phi_1,\ldots,\phi_d)
=(\phi_1\otimes\cdots\otimes\phi_d)
\bigl(A(v_1\otimes\cdots\otimes v_d)\bigr).
\]
有限维张量配对的非退化性给出存在性与唯一性。\(F\) 不变当且仅当 \(A\) 与每个 \(g^{\otimes d}\) 交换。对 \(d\ge1\)，\cref{b5:thm:schur-weyl}在任意特征零域上给出的第二个中心化子等式适用于当前的 \(k,V\)，说明 \(A\) 是因子置换算子 \(P_\sigma\) 的线性组合。其对应函数为
\[
\prod_{j=1}^d\phi_j\bigl(v_{\sigma^{-1}(j)}\bigr),
\]
故每个多线性不变量都是配对乘积的线性组合。次数零的情形只给常数。

最后对一般不变量 \(f\) 取多重齐次分量；作用不混合各变量组，所以每个分量仍不变。设它在 \(v_i\) 上的次数为 \(a_i\)，在 \(\phi_j\) 上的次数为 \(b_j\)。将
\[
v_i=\sum_{\nu=1}^{a_i}t_{i\nu}v_{i\nu},\qquad
\phi_j=\sum_{\mu=1}^{b_j}u_{j\mu}\phi_{j\mu}
\]
代入，并取所有参数各一次的系数。所得函数 \(\widetilde f\) 对新增变量多线性，且不变，因为代入和取系数均与群作用交换。让同一组副本相等，得到
\[
\widetilde f=\left(\prod_i a_i!\right)
\left(\prod_j b_j!\right)f.
\]
这由齐次性及多项展开的多项式系数给出。特征零允许除以这些阶乘；先把 \(\widetilde f\) 写成配对乘积，再合并副本，便把 \(f\) 写成原配对的多项式。
\end{proof}

这个结论只断言生成性。若 \(\dim V=n\)，配对矩阵 \((p_{ji})\) 的秩至多 \(n\)，故其 \((n+1)\) 阶子式都是关系；证明这些子式生成全部关系，是另外的关系定理，不能由生成性直接得到。

\ProseParagraphBreak
在二维交错空间 \(V=k^2\) 中，标准交错型为
\(\omega(v,w)=\det(v,w)\)，保持它的群是 \(\operatorname{SL}_2\)。对两个向量组成的矩阵 \(X=(v,w)\)，有
\[
k[V\oplus V]^{\operatorname{SL}_2}=k\bigl[\omega(v,w)\bigr]
\]
（这里仍设特征零）。证明也很具体：若 \(\det X=t\ne0\)，矩阵
\(\operatorname{diag}(1,t)X^{-1}\) 的行列式为一，把 \(X\) 送到 \(\operatorname{diag}(1,t)\)。所以不变量 \(f\) 在可逆矩阵上等于多项式
\(F(\det X)\)，其中 \(F(t)=f\bigl(\operatorname{diag}(1,t)\bigr)\)；由稠密性，这个等式在全部矩阵上成立。这里的交错 Gram 矩阵只有一个独立条目，其余是零及其负值。

\subsection{二元型的判别式与协变式}
\label{b5:subsec:2205:04}

\begin{definition}\label{b5:def:relative-invariant-covariant}
设 \(k\) 为域，\(G\) 为 \(k\) 上的仿射代数群，有理作用在有限维 \(k\) 空间 \(U,W\) 上。多项式函数 \(f:U\to k\) 若作为代数态射满足
\(f(gu)=\chi(g)f(u)\)，其中 \(\chi:G\to\mathbb G_m\) 为代数群同态，则称 \(f\) 为具有特征 \(\chi\) 的\term[xiangdui-bubianliang]{相对不变量}[relative invariant]。若多项式映射 \(F:U\to W\) 作为代数态射满足 \(F(gu)=gF(u)\)，则称 \(F\) 为\term[xiebian-shi]{协变式}[covariant]。相对不变量也可看作取值于一维表示 \(k_\chi\) 的协变式，其中 \(G\) 在 \(k_\chi\) 上按 \(\chi\) 数乘作用。
\end{definition}
\symidx{\(k_\chi\)：代数群字符 \(\chi\) 对应的一维有理表示}

设 \(k\) 的特征零，令 \(\operatorname{GL}_2\) 在二元型上作用为
\((g\cdot F)(X,Y)=F\bigl(g^{-1}(X,Y)^T\bigr)\)。对二次型
\[
F=aX^2+bXY+cY^2,\qquad
Q=\begin{pmatrix}a&b/2\\b/2&c\end{pmatrix},
\]
变换后矩阵为 \(g^{-T}Qg^{-1}\)，故判别式
\(\Delta=b^2-4ac=-4\det Q\) 满足
\[
\Delta(g\cdot F)=(\det g)^{-2}\Delta(F).
\]
因此它在一般线性群下是相对不变量，在特殊线性群下是不变量。

\ProseParagraphBreak
一个实际的协变构造是 Hessian。对次数 \(d\ge2\) 的二元型，置
\[
\mathcal H(F)=
\det\begin{pmatrix}F_{XX}&F_{XY}\\F_{XY}&F_{YY}\end{pmatrix}.
\]
链式法则给出
\[
\mathcal H(g\cdot F)(X,Y)
=(\det g)^{-2}\,
\mathcal H(F)\bigl(g^{-1}(X,Y)^T\bigr).
\]
故在 \(\operatorname{SL}_2\) 下，\(\mathcal H\) 是从 \(d\) 次型空间到 \(2d-4\) 次型空间的协变式；在 \(\operatorname{GL}_2\) 下须把值域作用另乘 \((\det g)^{-2}\)。例如
\[
F=aX^3+3bX^2Y+3cXY^2+dY^3
\]
给出
\[
\mathcal H(F)=36\bigl((ac-b^2)X^2+(ad-bc)XY+(bd-c^2)Y^2\bigr).
\]
它随型变成另一个型，通常不是标量不变量。区分“系数次数”和“关于 \(X,Y\) 的次数”也很必要：这里前者为二，后者也为二；对一般 \(d\)，二者分别为二和 \(2d-4\)。

% !TeX root = ../../Book5.tex
% 纲要标题：### 22.6 仿射商与轨道闭包
% 纲要位置：工作记录/计划纲要.md，第 4502 行。
% * 从有限群不变环及有限模性建立仿射商，写清泛性质、底域及几何点条件
% * 用环面作用的具体例子区分轨道与商点；一般闭轨道定理逐项标明假设和证明状态
% * 通过矩阵的半单部分与轨道闭包解释多项式不变量不能区分全部共轭轨道
% #### 22.6.1 由不变环构造仿射商
% #### 22.6.2 有限群商与有限态射
% #### 22.6.3 代数环面作用与轨道闭包
% #### 22.6.4 矩阵共轭与闭轨道
\section{仿射商与轨道闭包}
\label{b5:sec:2206}

不变环可以看作商空间上的函数环，但“商”不一定就是原有轨道组成的集合。有限群作用时二者十分接近；连续变化的代数群作用则会把不同轨道合并。下面先从坐标环说明商的确切含义，再计算这些合并怎样发生。

\subsection{由不变环构造仿射商}
\label{b5:subsec:2206:01}

\begin{definition}\label{b5:def:affine-invariant-quotient}
设 \(G\) 为域 \(k\) 上的仿射代数群，作用在仿射 \(k\) 概形 \(X=\operatorname{Spec}A\) 上，\(\delta:A\to A\otimes_k k[G]\) 为诱导的右余作用。若不变子代数 \(B=A^G\) 是有限生成 \(k\) 代数，则把
\[
X/\!/G=\operatorname{Spec}B,\qquad
\pi:X\longrightarrow X/\!/G
\]
称为由不变环定义的\term[fangshe-shang]{仿射商}[affine quotient]；其中 \(\pi\) 由包含 \(B\hookrightarrow A\) 给出。不变元是余作用满足 \(\delta(b)=b\otimes1\) 的元，这一定义也适用于非代数闭域和群概形。
\end{definition}

\begin{proposition}\label{b5:prop:affine-quotient-universal}
设 \(k\) 为域，仿射代数群 \(G\) 作用在 \(X=\operatorname{Spec}A\) 上，且 \(B=A^G\) 是有限生成 \(k\) 代数。令 \(\pi:X\to X/\!/G=\operatorname{Spec}B\) 为包含 \(B\hookrightarrow A\) 定义的态射。对任意仿射 \(k\) 概形 \(Y=\operatorname{Spec}C\)，每个 \(G\) 不变态射 \(f:X\to Y\) 都唯一分解为
\[
X\xrightarrow{\pi}X/\!/G\xrightarrow{\overline f}Y.
\]
这里 \(G\) 在 \(Y\) 上平凡作用，不变性是作用态射与投影复合后相等。
\end{proposition}
\begin{proof}
态射 \(f\) 对应代数同态 \(f^*:C\to A\)。不变性恰表示其每个像元素满足余作用等式，因而像包含在 \(B\) 中。限制值域得到唯一同态 \(C\to B\)，反变地就是所需 \(\overline f\)。唯一性来自 \(B\hookrightarrow A\) 为单射。
\end{proof}

这里证明的是相对于仿射目标的泛性质；它没有自行断言对任意概形目标也成立。即使 \(\pi\) 存在，同一个商点也可能代表多个轨道。

\subsection{有限群商与有限态射}
\label{b5:subsec:2206:02}

\begin{theorem}\label{b5:thm:finite-group-quotient}
设 \(k\) 为代数闭域，有限群 \(G\) 作用在仿射代数集 \(X=\operatorname{Spec}A\) 上。则 \(B=A^G\) 有限生成，商态射
\(\pi:X\to\operatorname{Spec}B\) 有限且满射；在 \(k\) 点上，其每个纤维恰为一个 \(G\) 轨道。
\end{theorem}
\begin{proof}
\cref{b5:thm:finite-invariant-generation}给出 \(B\) 有限生成及 \(A\) 为有限 \(B\) 模；后者正是仿射态射有限的定义。整扩张的上覆性质给出素谱满射。回顾该性质在此的理由：对 \(\mathfrak p\in\operatorname{Spec}B\)，局部化 \(A_{B\setminus\mathfrak p}\) 非零且在 \(B_{\mathfrak p}\) 上整；取其极大理想，收缩是 \(B_{\mathfrak p}\) 的极大理想，所以为 \(\mathfrak pB_{\mathfrak p}\)。收缩为极大的依据是整扩张中，上环为域时下环也为域：任意非零下环元素的逆在上环中满足首一整性方程，乘以适当幂就把该逆写回下环。再收缩到 \(A\) 即得覆盖 \(\mathfrak p\) 的素理想。

对 \(B\) 的极大理想，覆盖它的素理想亦为极大，商域是 \(k\) 的有限扩张；代数闭性使它仍是 \(k\)，所以 \(k\) 点上的映射也满射。同一轨道显然有同样的不变量值。若两个点不在同一轨道，\cref{b5:prop:finite-invariants-separate}给出在它们上取零、一的不变量，故不会映到同一商点。于是各纤维恰为轨道。
\end{proof}

对同时变号作用，商是方程 \(uw=v^2\) 定义的二次锥，商映射为
\[
(x,y)\longmapsto(x^2,xy,y^2).
\]
在特征不为二的代数闭域上，每个非原点纤维有两个点，原点纤维只有一个点。商可以产生奇点，即使原来的平面光滑；这是不能把“有限商”理解成处处局部同构的一个简单理由。

\subsection{代数环面作用与轨道闭包}
\label{b5:subsec:2206:03}

设 \(k\) 为代数闭域，\(\mathbb G_m\) 在 \(\mathbb A^2\) 上作用为
\[
t\cdot(x,y)=(tx,t^{-1}y).
\]
由单项式的权，\(k[x,y]^{\mathbb G_m}=k[xy]\)，所以仿射商映射是 \(\pi(x,y)=xy\)。当 \(c\ne0\) 时，曲线 \(xy=c\) 是一个轨道：给定两点，取 \(t=x'/x\) 即可互相变换，而且该轨道为闭集。

\ProseParagraphBreak
零纤维却包含三个轨道：原点、去掉原点的 \(x\) 轴、去掉原点的 \(y\) 轴。后两个轨道都不是闭集，其闭包都包含原点；只有原点是闭轨道。因此商点零合并了三个轨道，而保留了它们共同趋向的闭轨道。

\ProseParagraphBreak
若改成 \(t\cdot(x,y)=(tx,ty)\)，所有非恒定单项式都有非零权，不变环只有 \(k\)，仿射商只有一点。每条过原点直线去掉原点都给出一个轨道，其闭包仍包含原点。在射影空间中这些直线却是不同的点；要保留这种信息，需要后面带有线性化的射影构造，不能直接沿用这个仿射商。

\subsection{矩阵共轭与闭轨道}
\label{b5:subsec:2206:04}

\begin{proposition}\label{b5:prop:matrix-closed-orbits}
设 \(k\) 为代数闭域。矩阵 \(A\in M_n(k)\) 的共轭轨道为 Zariski 闭集，当且仅当 \(A\) 半单，即可对角化。每个轨道闭包包含一个且仅一个半单共轭轨道，它由 \(A\) 的特征多项式确定。
\end{proposition}
\begin{proof}
在 Jordan 标准形中写 \(A=D+N\)，其中 \(D\) 在各块上为标量，\(N\) 具有通常的超对角一。在一个大小 \(r\) 的块上用
\(\operatorname{diag}(t^{r-1},t^{r-2},\ldots,1)\) 共轭，便把 \(N\) 变为 \(tN\)。合并各块得到：\(t\ne0\) 时 \(D+tN\) 与 \(A\) 共轭，\(t=0\) 时等于 \(D\)。由多项式映射 \(t\mapsto D+tN\)，\(D\) 在轨道闭包中；若 \(A\) 不是半单的，二者不共轭，原轨道因此不闭。

反过来，设 \(D\) 的不同特征值为 \(\lambda_1,\ldots,\lambda_r\)，其固定特征多项式为 \(P(T)\)。矩阵集合
\[
\left\{\,B\in M_n(k)\ \middle|\
\det(TI-B)=P(T),\
\prod_{i=1}^r(B-\lambda_iI)=0\,\right\}
\]
是闭集，因为它由矩阵条目的多项式方程定义。第二条件使最小多项式无重根，故 \(B\) 可对角化；第一条件又固定各特征值的重数，所以这个闭集恰为 \(D\) 的共轭轨道。半单轨道于是闭。

特征多项式系数在一个轨道闭包上恒定。任何落在闭包中的半单矩阵均具有 \(A\) 的特征多项式，因而与 \(D\) 共轭。这给出唯一性。
\end{proof}

\cref{b5:thm:single-matrix-invariants}与此命题相互解释：所有不变量只记住半单共轭类，而不会区分具有相同特征值和重数的不同 Jordan 块。一般约化群仿射作用也有以闭轨道解释商纤维的理论；这里只对上面的环面模型和矩阵共轭完整证明该现象，不由个别模型推断任意代数群都有同样性质。

% !TeX root = ../../Book5.tex
% 纲要标题：### 22.7* 约化群与几何不变量理论初步
% 纲要位置：工作记录/计划纲要.md，第 4776 行。
% * 在第21章有理表示语言下完整证明所采用的线性约化群有限生成版本，不默认所有群的不变环均有限生成
% * 区分特征零与正特征；Haboush 定理给出旗簇与 Steinberg 输入明确的证明概要，幂提升与 Nagata 有限生成完整证明，不用有限群平均代替
% * 给出一个含明确线性化的低维射影例子；完整 Hilbert–Mumford 判据、模空间与变分 GIT 留给后续专著
% #### 22.7.1 线性约化性与 Hilbert 有限生成定理
% #### 22.7.2 正特征、几何约化性及适用范围
% #### 22.7.3 稳定性与射影商的一个例子
% ---
\OptionalSection{约化群与几何不变量理论初步}{b5:sec:2207}

有限群平均可以把系数投影到不变部分。对于无限代数群，有理表示的完全可约性也能给出这样的投影。本节先用这一投影证明有限生成，再介绍正特征的几何约化性与射影商的构造。

\subsection{线性约化性与 Hilbert 有限生成定理}
\label{b5:subsec:2207:01}

\begin{definition}\label{b5:def:linearly-reductive}
设 \(G\) 为域 \(k\) 上的仿射群概形。若每个有限维有理 \(G\) 表示都可分解为简单有理表示的直和，则称 \(G\) 为\term[xianxing-yuehua]{线性约化群}[linearly reductive group]。对可能无限维的有理表示，本节使用右 \(k[G]\) 余模的含义。
\end{definition}

\begin{lemma}\label{b5:lem:reductive-reynolds}
设 \(G\) 为域 \(k\) 上的线性约化仿射群概形，\(M\) 为有理 \(G\) 表示，\(M^G\) 为其不动子空间。则有自然的 \(G\) 等变投影
\(\mathcal R_M:M\to M^G\)。取不变量保持有理表示的短正合列。若 \(A\) 是有理 \(G\) 代数，则 \(\mathcal R_A\) 是 \(A^G\) 线性的。
\end{lemma}
\begin{proof}
有限维时，将 \(M\) 分解为简单模。平凡简单模之和恰为 \(M^G\)，其余简单模之和记为 \(M'\)。没有非零同态能在平凡简单模与非平凡简单模之间来回映射：非零简单模同态必为同构。因此这两个等型部分及投影均不依赖分解的选择，并且任意等变同态保持它们。

由\cref{b5:lem:comodule-locally-finite}，一般 \(M\) 中每个向量都属于有限维子表示。两个这样的子表示都包含在它们的有限维和中；投影的自然性使定义相容。因此有限维投影黏合为 \(\mathcal R_M\)，且仍对所有同态自然。

取不变量本来左正合。对等变满射 \(M\to N\) 和 \(n\in N^G\)，选原像 \(m\)；自然性说明 \(\mathcal R_M(m)\) 映到 \(\mathcal R_N(n)=n\)，所以右端也满射。最后，若 \(a\in A^G\)，乘法映射 \(x\mapsto ax\) 是等变线性映射；对它用自然性即得
\(\mathcal R_A(ax)=a\mathcal R_A(x)\)。
\end{proof}

\begin{theorem}[线性约化群的有限生成]\label{b5:thm:reductive-hilbert-finiteness}
设 \(k\) 为域，\(G\) 为 \(k\) 上的线性约化群，\(A\) 为有限生成交换 \(k\) 代数，具有有理 \(G\) 作用。则 \(A^G\) 是有限生成 \(k\) 代数。
\end{theorem}
\begin{proof}
先设 \(A=S=\operatorname{Sym}(W)\)，其中 \(W\) 有限维，作用保持标准分次。令 \(I\subseteq S\) 为全部正次齐次不变量生成的理想。Hilbert 基定理说明 \(S\) 为 Noether 环，所以可从这些不变量中选出有限个齐次元 \(f_1,\ldots,f_r\) 生成 \(I\)。确实，先取有限理想生成元，再展开它们所用的有限多个原始不变量即可。

对不变量的次数归纳。常数不需要生成；若 \(f\in S_d^G,\ d>0\)，由 \(f\in I\) 可写
\[
f=\sum_i a_if_i,\qquad \deg a_i=d-\deg f_i<d,
\]
只保留相应齐次部分即可。\cref{b5:lem:reductive-reynolds}以及分次投影的自然性给出
\[
f=\sum_i\mathcal R_S(a_i)f_i.
\]
这些系数不变且次数下降，由归纳假设属于 \(k[f_1,\ldots,f_r]\)，所以 \(S^G=k[f_1,\ldots,f_r]\)。

一般情形先取 \(A\) 的有限个代数生成元。余模局部有限性使它们包含在有限维 \(G\) 稳定子空间 \(W\subseteq A\) 中。乘法延拓给出等变满射
\(\operatorname{Sym}(W)\to A\)。引理的不变量正合性又给出
\(\operatorname{Sym}(W)^G\to A^G\) 满射。前者已经有限生成，故其商 \(A^G\) 也有限生成。
\end{proof}

分裂环面在任意特征都线性约化：\cref{b5:thm:torus-weight-decomposition}把表示分成一维权表示，Reynolds 算子就是取权零部分。有限群在群阶可逆时则由 Maschke 定理得到这一性质。这里的一般证明没有给出类似 \(|G|\) 的次数上界；有限生成与有效的次数估计仍应区分。

\subsection{正特征、几何约化性及适用范围}
\label{b5:subsec:2207:02}

\cref{b5:def:reductive-group}中的约化性排除非平凡的光滑连通正规幺幂闭子群，是群本身的结构条件；线性约化性则是全部有限维有理表示的分解条件。在代数闭域上，特征零时这两种条件对光滑连通仿射代数群等价，正特征时不能这样替换。

\ProseParagraphBreak
例如在特征 \(p>0\) 的代数闭域上，考虑 \(\operatorname{SL}_2\) 的 \(p\) 次对称幂 \(k[X,Y]_p\)。子空间
\[
U=kX^p+kY^p
\]
因 \((aX+bY)^p=a^pX^p+b^pY^p\) 而稳定。假设存在稳定补空间 \(L\)。对角环面的权分别为 \(p,p-2,\ldots,-p\)，作为代数群的特征它们仍互不相同；中间权 \(p-2\) 迫使 \(X^{p-1}Y\in L\)。然而使 \(X\mapsto X,\ Y\mapsto Y+tX\) 的幺幂元将它送为
\[
X^{p-1}Y+tX^p.
\]
取 \(t\ne0\)，稳定性便迫使 \(X^p\in L\cap U\)，矛盾。因此这个表示不完全可约，即使 \(\operatorname{SL}_2\) 是约化群。

\begin{definition}\label{b5:def:geometrically-reductive}
设 \(k\) 为代数闭域，\(G\) 为 \(k\) 上的仿射代数群。若对每个有限维有理 \(G\) 表示 \(V\) 及每个非零固定向量 \(v\in V^G\)，都有某个正次齐次不变量 \(f\in k[V]^G\) 满足 \(f(v)\ne0\)，则称 \(G\) 为\term[jihe-yuehua]{几何约化群}[geometrically reductive group]。
\end{definition}

线性约化时可将不变直线 \(kv\) 等变地投影出来，得到在 \(v\) 上非零的不变线性函数，因此必几何约化。几何约化允许使用更高次数，这是它能够适应正特征的关键区别。

\begin{theorem}[Haboush 定理]\label{b5:thm:haboush-geometric-reductivity}
设 \(k\) 为代数闭域，\(G\) 为 \(k\) 上的连通约化仿射代数群。则 \(G\) 几何约化。
\end{theorem}
\begin{proofsketch}
给定非零固定向量 \(v\in V^G\)，要构造的具体对象是一个有限维非零表示 \(E\) 及等变线性映射
\[
\Theta:V\longrightarrow\operatorname{End}(E),\qquad
\Theta(v)=I_E,
\]
其中 \(G\) 在自同态空间上按共轭作用。一旦做到，\(\det\circ\Theta\) 就是次数 \(\dim E>0\) 的不变多项式，并在 \(v\) 处取值一。以下构造始终保留 \(v\mapsto1\mapsto I_E\)。不能一般用迹代替行列式：特征 \(p\) 中若 \(p\mid\dim E\)，则 \(\operatorname{tr}(I_E)=0\)，而 \(\det(I_E)=1\) 仍成立。

约化群的结构理论给出忠实平坦的有限满同态
\[
 C\times G_{\mathrm{sc}}\longrightarrow G,
\]
其中 \(C\) 为 \(G\) 的最大中心环面，\(G_{\mathrm{sc}}\) 为导出群的半单单连通覆盖，核为有限中心子群概形。坐标环的拉回单射使拉回表示中的不变向量与不变多项式保持不变。环面的有理表示按权分解：先投影到 \(C\) 的零权子空间，投影与 \(G_{\mathrm{sc}}\) 作用交换，且保持给定固定向量。若在该子空间上找到 \(G_{\mathrm{sc}}\) 不变的齐次多项式，将它沿投影拉回，便得到乘积群的不变量。因此只须处理半单单连通群。固定最大环面 \(T\)。

给定 \(0\ne v\in V^G\)，选 \(\ell\in V^*\) 使 \(\ell(v)=1\)。矩阵系数映射
\[
 \phi:V\longrightarrow k[G],\qquad
 \phi(w)(g)=\ell(g^{-1}w)
\]
对于左平移 \((h\cdot f)(g)=f(h^{-1}g)\) 等变，且 \(\phi(v)=1\)。右 \(T\) 作用与左 \(G\) 作用交换；取右权零部分的 Reynolds 投影，便得到 \(G\) 等变映射 \(\pi\phi:V\to k[G]^T\)，仍把 \(v\) 送到 \(1\)。

关键的旗簇构造给出有限维 \(G\) 子模的递增并
\[
 k[G]^T=\bigcup_{n\ge0}M_n,\qquad
 M_n\simeq E_{n\rho}\otimes E_{n\rho},\qquad
 E_{n\rho}=H^0(G/B,\mathcal L_{n\rho}),
\]
其中 \(B\) 为 Borel 子群，\(\rho\) 为正根半和。取相反的 Borel 子群 \(B^-\)，利用 \(G/T\) 在 \(G/B\times G/B^-\) 中的开轨道，并清去大胞腔截面的分母，便得到这些子模。常数 \(1\) 对应 \(G\) 同态 \(E_{n\rho}^*\to E_{n\rho}\)，其像是由最高权向量生成的简单子模 \(E'_{n\rho}\)。当 \(E'_{n\rho}=E_{n\rho}\) 时，该同态为同构，于是 \(M_n\) 可识别为按共轭作用的 \(\operatorname{End}(E_{n\rho})\)，且 \(1\) 对应恒等自同态。

在 \(G=\operatorname{SL}_2\) 中，这些分母可以写成实际函数。\(G/T\) 由两条不同直线组成的有序对表示，因而是 \(\mathbb P^1\times\mathbb P^1\) 中的开集。取两组齐次坐标 \(u=(u_0,u_1)\)、\(z=(z_0,z_1)\)，令
\[
s(u,z)=u_0z_1-u_1z_0.
\]
两直线不同恰为 \(s\ne0\)。截面 \(a\in H^0(\mathbb P^1\times\mathbb P^1,\mathcal O(n,n))\) 是两组变量各 \(n\) 次齐次的多项式，因此 \(a/s^n\) 与齐次代表的缩放无关，在这个开集上给出正则函数。任一这样的函数都可以清分母到某个 \(a/s^n\)，而将 \(a\) 换成 \(as\) 给出子模的递增包含；常数一正来自截面 \(s^n\)。若记 \(E_n=H^0(\mathbb P^1,\mathcal O(n))\)，张量积 \(E_n\otimes E_n\) 识别为 \(\operatorname{Hom}(E_n^*,E_n)\)，于是 \(s^n\) 先给出等变映射 \(Q_n:E_n^*\to E_n\)。它此时还不必可逆。只有在所选的 Steinberg 情形证明 \(Q_n\) 是同构后，才可用
\(A\mapsto A Q_n^{-1}\) 把这个空间与 \(\operatorname{End}(E_n)\) 等变识别，并将 \(Q_n\) 送到 \(I_{E_n}\)。

特征零时 \(E_{n\rho}\) 不可约。特征 \(p>0\) 时，取
\[
 q=p^a,\qquad n=q-1,\qquad a\ge1.
\]
这里使用旗簇线丛、整模型及 Steinberg 表示\footnote{斯坦伯格（Robert Steinberg，1922-2014），加拿大裔美国数学家，研究代数群及其有限群的表示论。}的外部结果。Serre 消失将充分大 \(n\) 的截面模与特征零整模型的约化联系起来；该特征零模型的 Weyl 维数计算为 \((n+1)^N\)，其中 \(N\) 为正根数。在所用分裂整模型的有限群 \(G(\mathbb F_q)\) 上，最高权子模 \(E'_{(q-1)\rho}\) 的限制包含 Steinberg 模 \(\operatorname{St}_q\)。它不可约且维数为 \(q^N\)，故充分大的 \(q\) 满足
\[
 q^N=\dim\operatorname{St}_q
 \le\dim E'_{(q-1)\rho}
 \le\dim E_{(q-1)\rho}
 =q^N.
\]
因此 \(E'_{(q-1)\rho}=E_{(q-1)\rho}\)。指数 \(q-1\) 共尾，故这些不可约模给出所需的 \(\operatorname{End}(E)\) 子模递增并。上述线丛、整模型与有限群表示的比较采用\cite[Section 3, Proposition 2, pp. 142--143]{Demazure1976MumfordConjecture}中的代数群结果。

\(\pi\phi(V)\) 有限维，故包含在某个这样的子模中。选相应同构 \(\sigma\) 使 \(\sigma(1)=I_E\)，则
\[
 f(w)=\det\bigl(\sigma\pi\phi(w)\bigr)
\]
是正次齐次多项式，因行列式保持共轭而 \(G\) 不变，且
\(f(v)=\det(I_E)=1\)。这就得到几何约化定义要求的多项式；不需除以 \(\dim E\)。

这里省略了约化群的中心覆盖结构、大胞腔截面的构造，以及 Steinberg 模、Serre 消失与整模型的比较。矩阵系数、\(k[G]^T\) 的递增并和行列式论证见\cite[Section 3, Propositions 1--2, pp. 140--143]{Demazure1976MumfordConjecture}；原论文为\cite{Haboush1975GeometricallyReductive}。这些输入使正特征的行列式构造成立，却不能一般地给出不变线性投影。
\end{proofsketch}

\begin{lemma}[不变量的幂提升]\label{b5:lem:geometric-reductive-power-lifting}
设 \(k\) 为代数闭域，\(G\) 为几何约化仿射代数群，\(A\) 为有理 \(G\) 作用的交换 \(k\) 代数，\(I\subseteq A\) 为 \(G\) 稳定理想。对每个 \(u\in(A/I)^G\)，存在整数 \(m>0\) 及 \(b\in A^G\)，使 \(b\) 在 \(A/I\) 中的像为 \(u^m\)。
\end{lemma}
\begin{proof}
\(u=0\) 时取 \(b=0\)。否则选原像 \(a\in A\)，令 \(W\) 为其 \(G\) 平移生成的有限维子表示。所有平移在商中都映到 \(u\)，故商映射限制为等变满射 \(q:W\to ku\)；以 \(u\) 识别 \(ku=k\)，把 \(q\) 看作 \(W^*\) 中的非零固定向量。几何约化给出齐次不变量 \(F\in k[W^*]^G=\operatorname{Sym}(W)^G\)，满足 \(F(q)\ne0\)。除以这个非零标量，可设 \(F(q)=1\)。

包含 \(W\subseteq A\) 延拓为等变代数映射 \(\operatorname{Sym}(W)\to A\)。令 \(b\) 为 \(F\) 的像，便有 \(b\in A^G\)。若 \(F\) 的次数为 \(m>0\)，每个 \(w\in W\) 在商中等于 \(q(w)u\)，故 \(b\) 的商像恰为 \(F(q)u^m=u^m\)。
\end{proof}

\begin{theorem}[Nagata 有限生成定理]\label{b5:thm:geometric-reductive-finite-generation}
设 \(k\) 为代数闭域，\(G\) 为几何约化仿射代数群，\(A\) 为有限生成交换 \(k\) 代数，具有有理 \(G\) 作用。则 \(A^G\) 有限生成。
\end{theorem}
\begin{proof}
\(A=0\) 时结论显然，以下设 \(A\ne0\)，记 \(R=A^G\)。\cref{b5:lem:geometric-reductive-power-lifting}说明，对任意真 \(G\) 稳定理想 \(I\)，都有整扩张
\[
 B_I=R/(I\cap R)\subseteq D_I=(A/I)^G.
\]
若 \(D_I\) 是有限生成 \(k\) 代数，其有限个代数生成元均整于 \(B_I\)，便使 \(D_I\) 成为有限生成 \(B_I\) 模。\cref{b5:lem:invariants-artin-tate}于是给出 \(B_I\) 有限生成。

先设 \(A\) 非负分次、\(A_0=k\)，且作用保持分次。\(A\) 为 Noether 环，对齐次稳定理想作 Noether 归纳，可假设每个非零真齐次稳定理想 \(I\) 所给商的 \(D_I\) 已有限生成。这里允许 \(A\) 及其商带幂零元。若 \(R\) 没有正次元，则 \(R=k\)；否则取正次齐次 \(0\ne f\in R\)。

若 \(f\) 在 \(A\) 中不为零因子，则
\(fA\cap R=fR\)：若 \(fa\) 不变，其余作用给出
\((f\otimes1)(\delta(a)-a\otimes1)=0\)，而 \(A\otimes_k k[G]\) 作为 \(A\) 模自由，可以消去 \(f\)。取 \(I=fA\)，前述整性与 Artin--Tate 论证给出 \(R/fR\) 有限生成。取其有限个齐次代数生成元在 \(R\) 中的齐次提升 \(r_1,\ldots,r_s\)。对齐次 \(r\in R\)，先用这些提升写出模 \(fR\) 的表达，再将余项写为 \(fr'\)；\(r'\) 的次数低于 \(r\)。按次数归纳得到 \(R=k[f,r_1,\ldots,r_s]\)。

若 \(f\) 在 \(A\) 中为零因子，\(fb\) 不变不再允许直接消去 \(f\) 并断言 \(b\) 不变。乘以 \(f\) 的线性映射的核恰为 \(\operatorname{Ann}_A(f)\)，所以必须先取这个商；具体得到 \(R\)-模同构
\[
(A/\operatorname{Ann}_A(f))^G
\xrightarrow{\ \bar b\mapsto fb\ } fA\cap R.
\]
这个映射良定义且单射，因为两个代表的乘积之差为零恰好表示它们相差一个零化 \(f\) 的元。它的像恰为右侧：\(fb\) 不变等价于 \(\delta(b)-b\otimes1\) 属于 \(\operatorname{Ann}_A(f)\otimes_k k[G]\)，也就是 \(\bar b\) 不变。所用的核等式来自 \(A\otimes_k k[G]\) 作为 \(A\) 模自由。现在分别考察模去以下两个理想的商：
\[
 I_1=fA,\qquad I_2=\operatorname{Ann}_A(f)
\]
它们均为非零真齐次稳定理想。令 \(B_j=R/(I_j\cap R)\)、\(D_j=(A/I_j)^G\)。归纳假设及第一段的论证使 \(B_1,B_2\) 有限生成，且 \(D_j\) 为有限生成 \(B_j\) 模。将两个 \(B_j\) 的有限个代数生成元提升到 \(R\)，由这些提升生成有限型子代数 \(C\subseteq R\)。于是 \(C\to B_j\) 均满射，尤其 \(D_2\) 是有限生成 \(C\) 模。

第一个商 \(A/fA\) 控制任意不变量模去 \(fA\) 后的部分，使它可以由 \(C\) 中的元表示；第二个商 \(A/\operatorname{Ann}_A(f)\) 则通过上述同构控制剩余的 \(fA\cap R\)。它们处理的是同一个不变量的余类与余项。

取 \(D_2\) 的有限个 \(C\) 模生成元 \(\bar b_1,\ldots,\bar b_m\)，并取它们在 \(A\) 中的提升 \(b_i\)。\(\bar b_i\) 不变意味着 \(\delta(b_i)-b_i\otimes1\) 落在 \(I_2\otimes_k k[G]\)，所以 \(fb_i\in R\)。对任意 \(r\in R\)，由 \(C\to B_1\) 满射可选 \(c\in C\)，使 \(r-c=fb\)。余作用等式给出 \((f\otimes1)(\delta(b)-b\otimes1)=0\)；乘以 \(f\otimes1\) 的映射的核是 \(I_2\otimes_k k[G]\)，故 \(\bar b\in D_2\)。写 \(\bar b=\sum_i c_i\bar b_i\)，其中 \(c_i\in C\)，乘以 \(f\) 即得
\[
 r=c+\sum_i c_i fb_i,\qquad
 R=C+\sum_i Cfb_i.
\]
因此 \(R\) 由 \(C\) 的代数生成元与有限多个 \(fb_i\) 生成。两种情形都完成归纳。

一般 \(A\) 中，取包含 \(1\) 与有限个代数生成元的有限维稳定子空间 \(W\)。令 \(F_0A=k\)，\(F_dA\) 为至多 \(d\) 个 \(W\) 中元素的乘积所张成的空间，得到穷尽的稳定滤过。令 \(G\) 固定变量 \(z\)，构造 Rees 代数
\[
 \mathcal R_F(A)=\bigoplus_{d\ge0}F_dA\,z^d\subseteq A[z].
\]
它由 \(z\) 及 \(W\) 的有限基所给的 \(w_i z\) 生成，零次为 \(k\)，并具有保持分次的有理 \(G\) 作用，故已证情形给出 \(\mathcal R_F(A)^G\) 有限生成。代入 \(z=1\) 在不变环上满射到 \(A^G\)：对 \(a\in A^G\)，取 \(d\) 使 \(a\in F_dA\)，则 \(az^d\) 是其不变量原像。因此 \(A^G\) 也有限生成。
\end{proof}

双商归纳与幂提升的标准论证见\cite[Section 3.4, Theorem 3.3 and Lemma 3.5]{Dolgachev2003InvariantTheory}；上面的最后一步用稳定滤过的 Rees 代数归约到分次情形。结合\cref{b5:thm:haboush-geometric-reductivity}，便得到约化群在有限型交换代数上的不变环有限生成。

\subsection{稳定性与射影商的一个例子}
\label{b5:subsec:2207:03}

设代数闭域 \(k\) 上的线性约化群 \(G\) 线性作用于 \(V\)。这个线性作用诱导 \(\mathbb P(V)\) 及线丛 \(\mathcal O(-1)\) 上的作用；该线丛在点 \([v]\) 的纤维是直线 \(kv\)。取对偶便给出 \(\mathcal O(1)\) 的线性化。选择提升到线丛的作用称为选择线性化；同一个射影作用可以有不同选择。

\ProseParagraphBreak
线性化与齐次多项式的联系由截面空间直接给出：对 \(d\ge0\)，
\[
H^0(\mathbb P(V),\mathcal O(d))=\operatorname{Sym}^d(V^*),
\qquad f(av)=a^df(v).
\]
缩放代表向量使函数值按 \(d\) 次幂变化，正与 \(\mathcal O(d)\) 纤维的变换相容，所以齐次多项式给出射影空间上的截面。给定线性化也就规定了群在这些截面上的作用；不变齐次多项式对应不变截面。截面在 \([v]\) 处是否消失不依赖代表向量。因此半稳定性就是要求某个正次数的不变截面在该点不消失；若所有这些截面都消失，就没有可供商映射使用的 Proj 坐标。

\begin{definition}\label{b5:def:projective-semistability}
设 \(k\) 为代数闭域，线性约化群 \(G\) 线性作用在有限维 \(k\) 空间 \(V\) 上，并在 \(\mathbb P(V)\) 上取此作用诱导的 \(\mathcal O(1)\) 线性化。点 \([v]\in\mathbb P(V)\) 若存在正次齐次不变量 \(f\in k[V]^G\) 使 \(f(v)\ne0\)，称为\term[banwending-dian]{半稳定点}[semistable point]。若还可选择这样的 \(f\)，使轨道在仿射开集
\(D_+(f)=\bigl\{\,[w]\mid f(w)\ne0\,\bigr\}\) 中闭，且稳定子群有限，则称为\term[wending-dian]{稳定点}[stable point]。
\end{definition}

非半稳定点称为不稳定点；半稳定而非稳定的点称为严格半稳定点。这些条件都与线性化有关。先说明不变齐次环如何给出半稳定集上的商映射，再用一个环面作用计算其纤维。

\begin{proposition}\label{b5:prop:projective-invariant-quotient-construction}
设 \(k\) 为代数闭域，线性约化群 \(G\) 线性作用在有限维 \(k\) 向量空间 \(V\) 上，在 \(\mathbb P(V)\) 上取诱导的 \(\mathcal O(1)\) 线性化。记 \(S=k[V]\)、\(R=S^G\)，\(\mathbb P(V)^{\mathrm{ss}}\) 为半稳定点构成的开子集。则 \(R\) 是有限生成分次代数，并有自然的 \(G\) 不变态射
\[
 \pi:\mathbb P(V)^{\mathrm{ss}}\longrightarrow\operatorname{Proj}R.
\]
对每个非零正次齐次 \(f\in R\)，其限制
\(D_+(f)\to\operatorname{Spec}(R_f)_0\) 是由该仿射开集的不变环定义的仿射商。下标零表示齐次局部化的次数零部分。
\end{proposition}
\begin{proof}
\cref{b5:thm:reductive-hilbert-finiteness}给出有限生成，作用保持次数使 \(R\) 保持分次。半稳定定义说明 \(\mathbb P(V)^{\mathrm{ss}}\) 正是这些 \(D_+(f)\) 的并。按 Proj 的仿射图构造，\(D_+(f)\) 的坐标环为 \(A_f=(S_f)_0\)，目标相应图的坐标环为 \(B_f=(R_f)_0\)。

因为 \(f\) 不变，\(A_f^G=B_f\)。确实，\(A_f\) 的元可写成 \(a/f^m\)，其中 \(a\) 齐次且 \(\deg a=m\deg f\)；它不变时，其余作用满足 \(\delta(a)/(f^m\otimes1)=(a\otimes1)/(f^m\otimes1)\)。\(S\) 是整环，且 \(S\otimes_k k[G]\) 作为 \(S\) 模自由，所以 \(f\otimes1\) 不为零因子。清分母得到 \(\delta(a)=a\otimes1\)，即 \(a\in R\)，反向显然。这也适用于非光滑的线性约化群概形。\cref{b5:prop:affine-quotient-universal}于是给出每张图上的仿射商。

两张图的交是 \(D_+(ff')\)。在次数零环中，若 \(\deg f=r\)、\(\deg f'=s\)，这一交由不变函数 \((f')^r/f^s\) 的非零条件截出；两张图之间的坐标变换就是进一步局部化。因此上述商态射在交上相容，粘合为 \(\pi\)。这证明所陈述的射影构造；关于稳定点上的轨道商及完整数值判据，还须另作论证。
\end{proof}

\begin{example}\label{b5:exm:projective-torus-quotient}
设 \(k\) 为代数闭域，\(G=\mathbb G_m\) 在 \(V=k^3\) 上作用为
\[
t\cdot(x,y,z)=(tx,y,t^{-1}z).
\]
取这个线性作用给出的 \(\mathcal O(1)\) 线性化。单项式的权为 \(a-c\)，所以
\[
k[V]^G=k[y,xz],\qquad \deg y=1,\quad\deg(xz)=2.
\]
正次不变量同时为零的射影点恰为
\([1:0:0]\)、\([0:0:1]\)。因此半稳定集是去掉这两点的开集。

记 \(R=k[y,xz]\)，两张仿射图的坐标环可直接计算为
\[
(R_y)_0=k\left[\frac{xz}{y^2}\right],\qquad
(R_{xz})_0=k\left[\frac{y^2}{xz}\right].
\]
第一式中，次数零单项式把 \(xz\) 的每一次幂配上 \(y^{-2}\)；第二式中，\(y\) 的指数必须为偶数。两图交上的这两个坐标互为倒数，因而拼成 \(\mathbb P^1\)。这也就是取第二个 Veronese 子环 \(k[y^2,xz]\) 的效果：把原来的偶数次数除以二，两个生成元的重分次均为一，次数零局部化及其交保持相同。商映射在本例具体为
\[
[x:y:z]\longmapsto[y^2:xz].
\]
两坐标在半稳定集上不同时为零，故这是正则映射。

\ProseParagraphBreak
例如经过 \([1:1:1]\) 的轨道可写成
\(t\mapsto[t:1:t^{-1}]=[t^2:t:1]\)、\(t\ne0\)。它在整个 \(\mathbb P^2\) 中的闭包是曲线 \(y^2=xz\)，参数 \(t=0\) 与 \(t=\infty\) 分别补上 \([0:0:1]\) 与 \([1:0:0]\)。这两个端点都已从半稳定集删除，因此该轨道在半稳定集内闭。相比之下，经过 \([1:1:0]\) 的轨道写成 \([t:1:0]\)，其闭包中的 \([0:1:0]\) 仍半稳定，故该轨道在半稳定集内不闭。\cref{b5:fig:projective-torus-boundaries}标出这两种不同的边界情形。
% !TeX root = ../Book5.tex
\begin{figure}[htbp]
\centering
\begin{tikzpicture}[x=1cm,y=1cm,every node/.style={font=\small}]
  \draw[AlgebraBlue,line width=.9pt]
    (-2,0) .. controls (-1.3,1.7) and (1.3,1.7) .. (2,0);
  \draw[->,AlgebraBlue] (-.8,1.22)--(-1.1,1.07);
  \draw[->,AlgebraBlue] (.8,1.22)--(1.1,1.07);
  \draw[fill=white] (-2,0) circle[radius=2pt];
  \draw[fill=white] (2,0) circle[radius=2pt];
  \draw[red!70!black] (-2.08,-.08)--(-1.92,.08);
  \draw[red!70!black] (-2.08,.08)--(-1.92,-.08);
  \draw[red!70!black] (1.92,-.08)--(2.08,.08);
  \draw[red!70!black] (1.92,.08)--(2.08,-.08);
  \fill (0,1.275) circle[radius=2pt];
  \node[above] at (0,1.275) {\([1:1:1]\)};
  \node[above] at (-1.8,.35) {\(t\to0\)};
  \node[above] at (1.8,.35) {\(t\to\infty\)};
  \node[below,align=center] at (-2,-.08) {\([0:0:1]\)\\不稳定，删除};
  \node[below,align=center] at (2,-.08) {\([1:0:0]\)\\不稳定，删除};
  \node at (0,.45) {\(y^2=xz\)};
  \begin{scope}[shift={(7,0)}]
    \draw[AlgebraBlue,line width=.9pt] (-2,0)--(2,0);
    \draw[->,AlgebraBlue] (-.8,0)--(-1.2,0);
    \fill (-2,0) circle[radius=2pt];
    \draw[fill=white] (2,0) circle[radius=2pt];
    \draw[red!70!black] (1.92,-.08)--(2.08,.08);
    \draw[red!70!black] (1.92,.08)--(2.08,-.08);
    \fill (0,0) circle[radius=2pt];
    \node[above] at (0,0) {\([1:1:0]\)};
    \node[above] at (-1.2,.25) {\(t\to0\)};
    \node[below,align=center] at (-2,-.08) {\([0:1:0]\)\\半稳定，保留};
    \node[below,align=center] at (2,-.08) {\([1:0:0]\)\\不稳定，删除};
    \node at (0,1.05) {\(z=0\)};
  \end{scope}
\end{tikzpicture}
\caption[射影轨道的边界]{两个射影轨道闭包及其边界点的示意。
左侧轨道的两个边界点都不在半稳定集中，右侧轨道却留下一个半稳定边界点。
空心叉号表示删除的点，实心点表示保留的点；曲线形状仅用于示意闭包中的包含关系。}
\label{b5:fig:projective-torus-boundaries}
\end{figure}


若 \(xz\ne0\)，稳定子群有限：当 \(y\ne0\) 时为单位群，当 \(y=0\) 时满足 \(t^2=1\)，作为群概形也是有限的。轨道在 \(D_+(xz)\) 中闭。事实上，给定比例 \(y^2/(xz)\)，在代数闭域上可用射影倍数及群参数把任意两个 \(xz\ne0\) 的点互相变换；所以该轨道是 \(D_+(xz)\) 中一个闭纤维的几何点集。这里不能把轨道与概形纤维直接等同：令 \(a=x/z\)、\(b=y/z\)，则这张图的坐标环为 \(k[a,a^{-1},b]\)，商坐标为 \(b^2/a\)。商点 \([0:1]\) 的概形纤维因此有坐标环
\[
 k[a,a^{-1},b]/(b^2),
\]
其中 \(b\) 是非零幂零元；约化后才得到轨道的坐标环 \(k[a,a^{-1}]\)。

若 \(xz=0,\ y\ne0\)，则 \([0:1:0]\) 是半稳定固定点，稳定子群为整个 \(\mathbb G_m\)，因而严格半稳定。对 \(xy\ne0\) 的点 \([x:y:0]\)，正则族
\[
 \mathbb A^1\longrightarrow\mathbb P(V)^{\mathrm{ss}},\qquad
 t\longmapsto[tx:y:0]
\]
在 \(t\ne0\) 时属于该点的轨道，在 \(t=0\) 时取值 \([0:1:0]\)。由 \(\mathbb G_m\) 在 \(\mathbb A^1\) 中稠密，这个固定点属于轨道闭包。对 \(yz\ne0\) 的点 \([0:y:z]\)，用 \(t\mapsto[0:y:tz]\) 得到同样结论。任一不变量在 \(xz=0\) 上都是 \(y\) 的多项式，所以凡包含原来的点的 \(D_+(f)\) 也包含这个固定点；这些非闭轨道上的点均严格半稳定。因此稳定集恰为 \(xz\ne0\)，而商点 \([1:0]\) 对应一个固定轨道与两个退化到它的非闭轨道。
\end{example}

这个例子已经显示射影商保留了仿射商未必保留的方向信息，也显示半稳定集必须先选出来。完整的 Hilbert--Mumford 判据\footnote{芒福德（David Bryant Mumford，1937-），美国数学家，在代数几何和几何不变量理论方面作出重要贡献；希尔伯特已在本卷前文介绍。}、一般商的几何性质与改变线性化时的变化，可以继续阅读\cite[Chapters 7--9]{Dolgachev2003InvariantTheory}；这些内容不由一个环面例子的计算代替。


\begin{chaptersummary}
有限群的轨道多项式给出整性，Artin--Tate 引理再给出不变环的有限生成；这个论证在模特征也成立。群阶可逆时，Reynolds 投影能够改造系数；特征零的极化进一步将生成次数限制在群阶以内。Molien 公式计算各齐次分量的维数，但不能独自代替生成元与关系的证明。

\ProseParagraphBreak
消去把代入同态的核写成一个理想的交。候选关系给出保持分次的满射之后，Hilbert 级数相等才保证没有遗漏关系。循环作用的曲面、矩阵的特征多项式及向量与对偶向量的配对这几类算例，分别采用单项式、稠密性和张量中心化子的方法。

\ProseParagraphBreak
有限群商的几何纤维是轨道，一般代数群的商则可能合并多个轨道。矩阵共轭下，每个轨道闭包中有唯一半单轨道，不变量正好记住它。线性约化性使 Reynolds 投影推广到有理表示，并支持 Hilbert 的有限生成证明；正特征的一般约化群需要几何约化理论。射影商还依赖线性化和半稳定点的选择。
\end{chaptersummary}

\BookExercises

\begin{optionalexercise}\label{b5:ex:22-coordinate-versus-points}
在 \(\mathbb F_q[x]\) 中，令乘法群 \(\mathbb F_q^\times\) 通过 \(x\mapsto ax\) 作用。求形式多项式的不变环，并给出其中在所有 \(\mathbb F_q\) 点上取零的非零元。
\end{optionalexercise}
\begin{solution}
单项式 \(x^d\) 不变当且仅当 \(q-1\mid d\)，因为 \(\mathbb F_q^\times\) 为循环群。因此不变环为 \(\mathbb F_q[x^{q-1}]\)。非零多项式 \(x^{q-1}(x^{q-1}-1)\) 在零点及每个非零点上都为零；在 \(q=2\) 时也得到 \(x(x-1)\)。所以到点集函数的求值映射并非单射。
\end{solution}

\begin{optionalexercise}\label{b5:ex:22-sign-one-coordinate}
设 \(\operatorname{char}k\ne2\)，\(s(x)=-x,s(y)=y\)。求 \(k[x,y]^{\langle s\rangle}\)，并写出 \(x^3+x^2y+xy^2+y^3\) 的 Reynolds 平均。
\end{optionalexercise}
\begin{solution}
不变单项式恰有偶数的 \(x\) 指数，所以不变环为 \(k[x^2,y]\)；二者代数独立，因为其不同幂积给出不同单项式。平均删去奇数 \(x\) 指数项，结果为 \(x^2y+y^3\)。
\end{solution}

\begin{optionalexercise}\label{b5:ex:22-normal-invariants}
设有限群 \(G\) 作用于交换 \(k\) 代数 \(A\)，且 \(N\trianglelefteq G\)。证明 \(A^N\) 上有良定的 \(G/N\) 作用，并证明 \((A^N)^{G/N}=A^G\)。
\end{optionalexercise}
\begin{solution}
若 \(a\in A^N\)，则对 \(n\in N\) 有 \(n(ga)=g(g^{-1}ng)a=ga\)，故 \(ga\in A^N\)。代表元 \(gn\) 与 \(g\) 在 \(A^N\) 上相同，因此商群作用良定。一个 \(a\in A^N\) 被每个陪集固定，当且仅当被每个 \(g\in G\) 固定，这正给出所述等式。
\end{solution}

\begin{optionalexercise}\label{b5:ex:22-orbit-polynomial}
设 \(C_3\) 循环置换 \(x,y,z\)。写出 \(x\) 的轨道多项式，并据此给出 \(k[x,y,z]\) 在不变环上的有限个模生成元；说明该组未必最小。
\end{optionalexercise}
\begin{solution}
轨道多项式为
\[
T^3-(x+y+z)T^2+(xy+yz+zx)T-xyz.
\]
三个变量都满足它，因此 \(x^ay^bz^c\)、\(0\le a,b,c<3\) 这 \(27\) 个元生成该模。首一关系保证可以逐次降低每个指数。这仅给上界；三元和等不变量还能产生额外线性关系，例如 \(z=(x+y+z)-x-y\)，所以不能据此声称模秩为 \(27\)。
\end{solution}

\begin{optionalexercise}\label{b5:ex:22-modular-affine}
设 \(\operatorname{char}k=p>0\)，阶 \(p\) 群在 \(k[x]\) 上按 \(\sigma(x)=x+1\) 作用。求不变环。
\end{optionalexercise}
\begin{solution}
令 \(z=x^p-x\)，则 \(\sigma(z)=z\)。每个多项式唯一写成 \(\sum_{i=0}^{p-1}a_i(z)x^i\)：首一关系给出存在，不同 \(i\) 的最高 \(x\) 次数模 \(p\) 不同给出唯一。若最大非零正指标为 \(i\)，差 \(\sigma(f)-f\) 的 \(x^{i-1}\) 基系数是 \(i\cdot a_i(z)\ne0\)，矛盾。因此 \(k[x]^{C_p}=k[z]\)。
\end{solution}

\begin{optionalexercise}\label{b5:ex:22-noether-sharp}
在复数域上构造一个阶为 \(m\) 的有限群线性作用，使不变环不能由次数严格小于 \(m\) 的元生成。
\end{optionalexercise}
\begin{solution}
取 \(C_m\) 在一维上由本原 \(m\) 次单位根作用。不变环为 \(\mathbb C[x^m]\)。所有小于 \(m\) 次的不变量都是常数，不能生成 \(x^m\)。因此统一的 Noether 上界在这族例子中达到。
\end{solution}

\begin{optionalexercise}\label{b5:ex:22-molien-scalar-three}
让 \(C_2\) 在 \(\mathbb C^3\) 上同时变号。求 Hilbert 级数及次数 \(2r\) 的系数，并指出二次生成元有多少个。
\end{optionalexercise}
\begin{solution}
Molien 公式给出
\[
H(t)=\frac12\bigl((1-t)^{-3}+(1+t)^{-3}\bigr)
=\frac{1+3t^2}{(1-t^2)^3}.
\]
次数 \(2r\) 的系数为 \(\binom{2r+2}{2}\)，奇数次为零。二次单项式有六个。每个偶数次单项式可把其变量因子两两配对，故六个二次元确能生成；级数却不等于 \((1-t^2)^{-6}\)，说明它们有关系。
\end{solution}

\begin{optionalexercise}\label{b5:ex:22-molien-reflection}
让 \(C_2\) 在 \(\mathbb C^2\) 上由 \(\operatorname{diag}(-1,1)\) 作用。用 Molien 公式求级数，并解释为什么它与同时变号作用不同。
\end{optionalexercise}
\begin{solution}
结果为
\[
\frac12\left(\frac1{(1-t)^2}+\frac1{1-t^2}\right)
=\frac1{(1-t)(1-t^2)}.
\]
它对应自由生成元 \(y,x^2\)，各次维数为 \(\lfloor d/2\rfloor+1\)。同时变号没有一次不变量且有二次元之间的关系；这两个作用虽来自同一个群，却不是同一个表示。
\end{solution}

\begin{optionalexercise}\label{b5:ex:22-cubic-cyclic}
令 \(C_3\) 在 \(\mathbb C[x,y]\) 上作用为 \(x\mapsto\zeta x,y\mapsto\zeta^{-1}y\)。在次数六中列出一组不变量基，并与展示 \(uv=w^3\) 比较。
\end{optionalexercise}
\begin{solution}
单项式 \(x^ay^{6-a}\) 的权为 \(2a-6\)，故 \(a=0,3,6\)。基为 \(y^6,x^3y^3,x^6\)。展示中它们分别为 \(v^2,w^3,u^2\)，而看似第四个的 \(uv\) 等于 \(w^3\)。带权次数分别为 \(\deg u=\deg v=3,\deg w=2\)，不能把 \(uv=w^3\) 当成通常三变量次数下的齐次关系。
\end{solution}

\begin{optionalexercise}\label{b5:ex:22-cusp-elimination}
对同态 \(k[u,v]\to k[t]\)、\((u,v)\mapsto(t^2,t^3)\)，求核及像的 Hilbert 级数，取 \(\deg t=1\)。
\end{optionalexercise}
\begin{solution}
关系为 \(v^2-u^3\)。模它后每个多项式形如 \(A(u)+v\cdot B(u)\)，其像分别含偶次幂 \(t^{2a}\) 和奇次幂 \(t^{2b+3}\)，互不相交且各自独立，故核就是该主理想。级数为
\[
\frac{1-t^6}{(1-t^2)(1-t^3)}
=1+t^2+t^3+t^4+\cdots.
\]
消去计算求出了子代数 \(k[t^2,t^3]\) 的生成关系。要把它识别为某个群作用的不变环，还须给出该群作用并证明相应的不变性与生成性。
\end{solution}

\begin{optionalexercise}\label{b5:ex:22-equal-series-not-isomorphism}
给出两个 Hilbert 级数相同但不同构的非负分次 \(k\) 代数，说明为什么不能省略\cref{b5:lem:hilbert-surjection-test}中的满射。
\end{optionalexercise}
\begin{solution}
取 \(A=k[x,y]/(xy)\) 与 \(B=k[x,y]/(x^2)\)，变量次数一。两者级数均为 \((1-t^2)/(1-t)^2\)。但 \(A\) 约化：若多项式幂被 \(xy\) 整除，两个素元 \(x,y\) 都整除该多项式，故它本身被 \(xy\) 整除；而 \(B\) 中非零的 \(x\) 平方为零。因此两环不同构。若存在保持次数的满射，级数判别反而会迫使同构，说明本例中不存在这种满射。
\end{solution}

\begin{optionalexercise}\label{b5:ex:22-symmetric-reduction}
在三个变量中把 \(x^2y^2+x^2z^2+y^2z^2\) 写成初等对称多项式，并由根为 \(1,2,3\) 的首一多项式的系数计算它在这组三根上的值。
\end{optionalexercise}
\begin{solution}
展开 \(e_2^2\) 得
\(x^2y^2+x^2z^2+y^2z^2+2xyz(x+y+z)\)，故所求为 \(e_2^2-2e_1e_3\)。根 \(1,2,3\) 给出 \(e_1=6,e_2=11,e_3=6\)，所以值为 \(121-72=49\)，也等于 \(4+9+36\)。
\end{solution}

\begin{optionalexercise}\label{b5:ex:22-signed-permutation}
在特征不为二的域上，求独立变号及坐标互换生成的二维群的不变环和 Hilbert 级数。
\end{optionalexercise}
\begin{solution}
先对变号取不变量得到 \(k[x^2,y^2]\)，再对互换取不变量得到 \(k[x^2+y^2,x^2y^2]\)。对称多项式基本定理说明两个生成元代数独立，次数为二、四，故级数为 \((1-t^2)^{-1}(1-t^4)^{-1}\)。
\end{solution}

\begin{optionalexercise}\label{b5:ex:22-newton-three}
在特征零域上，用 \(p_i=\operatorname{tr}(A^i)\)、\(1\le i\le3\)，表示三阶矩阵的特征多项式系数。
\end{optionalexercise}
\begin{solution}
Newton 递推依次给出
\[
c_1=p_1,\qquad c_2=\frac{p_1^2-p_2}{2},\qquad
c_3=\frac{p_1^3-3p_1p_2+2p_3}{6}.
\]
因此特征多项式为 \(T^3-c_1T^2+c_2T-c_3\)。分母说明这个直接计算在特征二、三中不能照搬；系数 \(c_i\) 本身仍可由行列式定义。
\end{solution}

\begin{optionalexercise}\label{b5:ex:22-pairing-rank-one}
设 \(k\) 为特征零域，\(V=k\)，取两个向量变量 \(a,b\) 及两个对偶变量 \(c,d\)。求 \(\operatorname{GL}(V)\) 作用下配对不变量的全部关系。
\end{optionalexercise}
\begin{solution}
四个配对为 \(u=ac,v=bc,w=ad,z=bd\)，满足 \(uz-vw=0\)。由配对定理，它们生成全部不变环。消去 \(uz\) 后留下不同时含 \(u,z\) 的单项式。若两个这样的单项式有相同的 \(a,b,c,d\) 指数，其 \(u,v,w,z\) 指数差必为 \((r,-r,-r,r)\)；若 \(r>0\)，第一个同时含 \(u,z\)，若 \(r<0\)，第二个同时含 \(u,z\)，均矛盾。因此标准单项式的像互不相同，核就是 \((uz-vw)\)。
\end{solution}

\begin{optionalexercise}\label{b5:ex:22-unbalanced-polarization}
设 \(k\) 特征零，\(V\ne0\)。证明 \(\operatorname{GL}(V)\) 在 \(V^{\oplus r}\) 上的多项式不变量只有常数，并说明为什么 \(\operatorname{SL}(V)\) 未必如此。
\end{optionalexercise}
\begin{solution}
次数 \(d>0\) 的齐次函数在标量矩阵 \(tI\) 下使函数值乘以 \(t^d\)。无限域中可取 \(t^d\ne1\)，不变性迫使函数为零。若 \(\dim V=n\)、\(r\ge n\)，在 \(\operatorname{SL}(V)\) 下，前 \(n\) 个向量的行列式是非恒定不变量，因为左乘群矩阵只使它乘以行列式一。
\end{solution}

\begin{optionalexercise}\label{b5:ex:22-binary-quadratic}
在特征零域上，计算 \(F=X^2+XY+Y^2\) 在 \(g=\operatorname{diag}(2,1)\) 下的新型及判别式，并核验相对不变因子。
\end{optionalexercise}
\begin{solution}
在特征零域中，
\[
g\cdot F=\frac14X^2+\frac12XY+Y^2.
\]
原判别式为 \(-3\)，新判别式为 \(1/4-1=-3/4\)，恰等于 \((\det g)^{-2}(-3)\)。若把它误当作一般线性群不变量，两数就会发生矛盾。
\end{solution}

\begin{optionalexercise}\label{b5:ex:22-hessian-cubic}
在特征零域上，计算 \(F=X^3+Y^3\) 与 \(F=X^3\) 的 Hessian，说明该构造如何辨认这两个型。
\end{optionalexercise}
\begin{solution}
第一种的二阶导数矩阵为 \(\operatorname{diag}(6X,6Y)\)，Hessian 为 \(36XY\)；第二种矩阵为 \(\operatorname{diag}(6X,0)\)，Hessian 为零。因此在特征零时，协变性质说明二者不可能在 \(\operatorname{GL}_2\) 下互相变换，因为变换不能把零 Hessian 变为非零 Hessian。
\end{solution}

\begin{optionalexercise}\label{b5:ex:22-finite-quotient-real}
令 \(C_2\) 在实直线上变号。写出实代数簇的商态射，并说明为什么它在实点上不满射，这是否与有限群商定理矛盾。
\end{optionalexercise}
\begin{solution}
不变环为 \(\mathbb R[x^2]\)，商态射为 \(x\mapsto u=x^2\)。其在实点上的像为非负实数，遗漏负数。素谱态射仍满射，例如 \(u+1\) 上方有极大理想 \((x^2+1)\)，其剩余域为 \(\mathbb C\)；它不是实点。正文的 \(k\) 点满射结论假定 \(k\) 代数闭，因此没有矛盾。
\end{solution}

\begin{optionalexercise}\label{b5:ex:22-torus-weight-two}
令 \(\mathbb G_m\) 在 \(k^2\) 上作用为 \(t\cdot(x,y)=(t^2x,t^{-1}y)\)，\(k\) 代数闭。求仿射不变环，描述非零商纤维及零纤维的轨道。
\end{optionalexercise}
\begin{solution}
单项式不变当且仅当 \(2a-b=0\)，故不变环为 \(k[xy^2]\)。若 \(xy^2=c\ne0\)，取 \(t=y/y'\) 就把任意一点送到另一点，因为 \(x'=x(y/y')^2\)，所以非零纤维是一个轨道。零纤维是两坐标轴，分为原点及两条去原点的轴；在代数闭域中平方映射满射，故去原点的 \(x\) 轴也只有一个轨道。
\end{solution}

\begin{optionalexercise}\label{b5:ex:22-nilpotent-degeneration}
设 \(k\) 为特征零的代数闭域。在 \(M_3(k)\) 中写出从大小三的幂零 Jordan 块到零矩阵的显式共轭退化，并说明秩为什么不能由这个作用下的多项式不变量给出。
\end{optionalexercise}
\begin{solution}
取 \(D(t)=\operatorname{diag}(t^2,t,1)\)，则 \(D(t)J_3D(t)^{-1}=tJ_3\)。当 \(t\to0\) 的代数族取值为零，而 \(t\ne0\) 时都属于同一轨道。若秩由一个不变多项式给出，该多项式在所有 \(t\ne0\) 上为二，因而作为 \(t\) 的多项式在零处也应为二；实际零矩阵秩为零，矛盾。
\end{solution}

\begin{optionalexercise}\label{b5:ex:22-reynolds-ideal-contraction}
设 \(G\) 线性约化，有理作用在有限生成交换代数 \(A\) 上，\(B=A^G\)。对理想 \(J\subseteq B\) 证明 \(JA\cap B=J\)。
\end{optionalexercise}
\begin{solution}
包含 \(J\subseteq JA\cap B\) 显然。若 \(b\in JA\cap B\)，写 \(b=\sum_i j_i a_i\)，其中 \(j_i\in J\)。应用 \(B\) 线性的 Reynolds 投影得 \(b=\sum_i j_i\mathcal R_A(a_i)\in J\)，所以等号成立。
\end{solution}

\begin{optionalexercise}\label{b5:ex:22-positive-characteristic-splitting}
在特征二的代数闭域上，将正文 \(\operatorname{SL}_2\) 表示 \(k[X,Y]_2\) 的不可分裂性写成一个三维矩阵计算。
\end{optionalexercise}
\begin{solution}
取基 \(X^2,XY,Y^2\)，子模 \(U=\langle X^2,Y^2\rangle\)。对角元 \(\operatorname{diag}(t,t^{-1})\) 的权为 \(2,0,-2\)，因此任何稳定一维补空间都必须为 \(kXY\)。但上幺幂变换 \(Y\mapsto Y+X\) 将 \(XY\) 送到 \(XY+X^2\)，其矩阵的第二列含非零的第一分量，所以 \(kXY\) 不稳定。这里权 \(2\) 与权 \(-2\) 是不同的环面特征，不能因域的特征二而将它们当作相同。
\end{solution}

\begin{optionalexercise}\label{b5:ex:22-projective-twist}
在 \(\mathbb P^1\) 上考虑 \(t\cdot[x:y]=[tx:y]\)。分别使用线性提升 \((x,y)\mapsto(tx,y)\) 和 \((x,y)\mapsto(x,t^{-1}y)\) 给出的 \(\mathcal O(1)\) 线性化，求半稳定集。
\end{optionalexercise}
\begin{solution}
两个线性提升相差整体标量 \(t^{-1}\)，所以给出同一射影作用。第一个提升的不变多项式环为 \(k[y]\)，半稳定集为 \(y\ne0\)；第二个为 \(k[x]\)，半稳定集为 \(x\ne0\)。因此同一射影作用在不同线性化下，半稳定集确可不同；不能只由点上的群作用确定它。
\end{solution}

